\documentclass[10pt,a4paper]{article}
\usepackage[T1]{fontenc}\usepackage[utf8]{inputenc}
\usepackage[a4paper,margin=2cm]{geometry}
\usepackage{booktabs,longtable,array,amsmath,amssymb,graphicx,xcolor,multirow,threeparttable}
\usepackage[hidelinks]{hyperref}
\newcommand{\SV}[1]{\mathrm{SV}_{#1}}
\renewcommand{\arraystretch}{1.1}
\setlength{\LTcapwidth}{\textwidth}
\setlength{\parindent}{0pt}\setlength{\parskip}{4pt}
\renewcommand{\thetable}{S\arabic{table}}
\title{\bfseries Supplementary Material\\[4pt]\large Multi-dimensional spectral geometry of biological knowledge in single-cell transformer representations}
\author{Ihor Kendiukhov}
\date{}
\begin{document}
\maketitle
Each table below summarises one analysis; the complete values are in the file named in its caption (folder \texttt{supplementary\_tables/}).
Extractions: lung (Tabula Sapiens lung, primary), lung (ext.) (lung atlas of Travaglini et al., replication), kidney~B, immune and kidney~A
(see the main text, Table~1). $L_\ell$ is layer $\ell$ (0--11); $\SV{k}$ is the $k$-th direction of the singular value decomposition of the centred
observed-gene matrix at that layer. Unless stated otherwise intervals are 95\%; network-statistic intervals use vertex subsampling with delete-$d$
scaling; edge-statistic intervals use the two-way (factor $\times$ target) pigeonhole bootstrap; node-statistic intervals resample genes.

\setcounter{table}{0}
\begingroup
\footnotesize
\setlength{\tabcolsep}{2.6pt}
\renewcommand{\arraystretch}{0.97}
\setlength{\LTleft}{0pt}\setlength{\LTright}{0pt}
\begin{longtable}{@{\extracolsep{\fill}}>{\scriptsize}l>{\scriptsize}l*{6}{>{\scriptsize}r}@{}}
\caption[Spectral concentration with depth]{\textbf{Spectral concentration with depth, in every extraction.} Singular values $S_1 \geq S_2 \geq \dots$ of the centred observed-gene matrix ($|V| \times 512$) at each layer. ER, effective rank $\exp(-\sum_k p_k \log p_k)$ with $p_k = S_k/\sum_j S_j$ (Roy and Vetterli, 2007); $\mathrm{ER}_{\mathrm{var}}$, exponential of the Shannon entropy of the variance shares $q_k = S_k^2/\sum_j S_j^2$; PR, participation ratio $(\sum_k S_k^2)^2/\sum_k S_k^4$; $\SV{1}$, share of variance $q_1$ on the first direction (\%); $\SV{1\text{--}10}$, cumulative share of the first ten (\%). \textbf{A}, the five main extractions at five layers ($|V|$ = 5{,}648, 5{,}616, 4{,}780, 2{,}039 and 1{,}972 genes for lung, lung (ext.), kidney~B, immune and kidney~A). Fold, $L_0$ value over $L_{11}$ value for ER, $\mathrm{ER}_{\mathrm{var}}$ and PR, $L_{11}$ over $L_0$ for the variance shares. Intervals for $\SV{1}$: 95\% gene-bootstrap percentile intervals ($B = 200$; genes resampled with replacement). No bootstrap interval is given for ER: resampling genes with replacement duplicates rows and biases it downwards, so its sampling variability is taken from the without-replacement subsamples of panel~D. \textbf{B}, all sixteen extractions. Cell draw, an independent draw of 256 immune cells; donor, 256 cells of one donor; B/CD4~T/CD8~T cells, 256 cells of one type; ckpt, one draw of 192 kidney cells (the 8{,}143-gene axis of kidney~B) passed through the whole-human checkpoint used for every other extraction, the kidney-fine-tuned and the brain scGPT checkpoint. Rises, number of the eleven consecutive layer steps at which ER / $\mathrm{ER}_{\mathrm{var}}$ / PR increased (0, strictly decreasing). \textbf{C}, identifiability and controls. Near-degenerate pairs: number of layers (of 12) at which the 95\% basic bootstrap interval of the gap ratio $S_k/S_{k+1}$ ($B = 200$ gene resamples) reaches 1; last column, count and percentage over all 192 extraction-layers of the sixteen extractions (all 2{,}112 pairs in the last of these rows). The nine such pairs among $\SV{1}$--$\SV{4}$ ($k \leq 3$) of the five main extractions are flagged again at $B = 1{,}000$. Feature shuffle: ER after each of the 512 features is permuted independently across genes (mean of 20 draws; SD $\leq 0.12$), divided by the observed ER, and the shuffled ER itself at $L_{11}$; last column, range over the sixteen extractions. Null direction: singular values below $10^{-3} S_1$ per layer, and the absolute cosine between the last right singular vector and the uniform vector $\mathbf{1}/\sqrt{512}$ at $L_0$ / $L_{11}$, with the relative size of the last singular value. \textbf{D}, detection controls. Matched: 1{,}300 genes subsampled without replacement within common deciles of $\log n_{\mathrm{obs}}$ pooled over the five extractions, mean $\pm$ SD over 200 draws; $n_{\mathrm{obs}} \geq 3$: the spectrum at $L_{11}$ after excluding genes contributed by fewer than three sampled cells (thresholds 5 and 10 in the CSV). Last column, range over the five extractions. Full values: \texttt{Table\_S1\_compression.csv}.}\label{tab:S1}\\
\toprule
\endfirsthead
\multicolumn{8}{@{}l}{\scriptsize\itshape Table~\thetable\ (continued)}\\
\toprule
\endhead
\midrule
\multicolumn{8}{r@{}}{\scriptsize\itshape continued on next page}\\
\endfoot
\bottomrule
\endlastfoot
\multicolumn{8}{@{}l}{\scriptsize\textbf{A\quad Spectral summaries at five layers, five main extractions}}\\
\midrule
Extraction & Measure & $L_0$ & $L_3$ & $L_6$ & $L_9$ & $L_{11}$ & Fold \\
\midrule
lung & ER & 457.4 & 435.1 & 402.4 & 354.2 & 324.1 & 1.41$\times$ \\*
 & $\mathrm{ER}_{\mathrm{var}}$ & 257.2 & 199.9 & 128.4 & 69.6 & 47.3 & 5.43$\times$ \\*
 & PR & 86.4 & 65.3 & 36.0 & 18.4 & 12.3 & 7.04$\times$ \\*
 & $\SV{1}$ (\%) & 7.3 & 7.3 & 11.2 & 16.9 & 24.0 & 3.27$\times$ \\*
 & \quad 95\% CI & 7.2--7.5 & 7.2--7.6 & 10.7--11.8 & 16.0--18.2 & 22.5--25.4 & \\*
 & $\SV{1\text{--}10}$ (\%) & 24.1 & 30.1 & 40.8 & 54.0 & 61.3 & 2.54$\times$ \\
\addlinespace[2pt]
lung (ext.) & ER & 459.4 & 439.0 & 409.8 & 363.3 & 340.8 & 1.35$\times$ \\*
 & $\mathrm{ER}_{\mathrm{var}}$ & 263.1 & 206.3 & 140.5 & 77.1 & 65.2 & 4.03$\times$ \\*
 & PR & 87.8 & 63.7 & 37.2 & 18.2 & 17.8 & 4.94$\times$ \\*
 & $\SV{1}$ (\%) & 7.3 & 7.7 & 11.4 & 18.7 & 17.9 & 2.45$\times$ \\*
 & \quad 95\% CI & 7.1--7.5 & 7.5--8.0 & 11.1--11.8 & 18.3--19.2 & 17.5--18.4 & \\*
 & $\SV{1\text{--}10}$ (\%) & 23.7 & 29.4 & 38.4 & 50.6 & 54.6 & 2.31$\times$ \\
\addlinespace[2pt]
kidney~B & ER & 456.2 & 434.0 & 401.0 & 352.8 & 324.1 & 1.41$\times$ \\*
 & $\mathrm{ER}_{\mathrm{var}}$ & 257.3 & 196.7 & 127.5 & 70.5 & 50.6 & 5.09$\times$ \\*
 & PR & 88.5 & 63.3 & 36.1 & 19.1 & 14.8 & 5.96$\times$ \\*
 & $\SV{1}$ (\%) & 7.0 & 7.1 & 10.9 & 16.5 & 17.6 & 2.50$\times$ \\*
 & \quad 95\% CI & 6.9--7.2 & 7.0--7.3 & 10.6--11.2 & 16.1--17.0 & 17.1--18.1 & \\*
 & $\SV{1\text{--}10}$ (\%) & 24.0 & 30.8 & 41.2 & 53.8 & 60.8 & 2.53$\times$ \\
\addlinespace[2pt]
immune & ER & 442.4 & 420.5 & 388.5 & 342.8 & 309.6 & 1.43$\times$ \\*
 & $\mathrm{ER}_{\mathrm{var}}$ & 236.9 & 186.5 & 125.5 & 74.7 & 48.7 & 4.86$\times$ \\*
 & PR & 90.1 & 68.1 & 40.2 & 22.8 & 15.2 & 5.93$\times$ \\*
 & $\SV{1}$ (\%) & 5.6 & 6.5 & 10.3 & 14.3 & 18.6 & 3.32$\times$ \\*
 & \quad 95\% CI & 5.4--6.0 & 6.2--7.0 & 9.8--11.0 & 13.7--15.0 & 16.7--20.9 & \\*
 & $\SV{1\text{--}10}$ (\%) & 25.4 & 31.2 & 40.7 & 52.2 & 61.3 & 2.42$\times$ \\
\addlinespace[2pt]
kidney~A & ER & 437.4 & 415.2 & 381.0 & 329.2 & 297.1 & 1.47$\times$ \\*
 & $\mathrm{ER}_{\mathrm{var}}$ & 209.5 & 158.0 & 98.1 & 52.0 & 34.8 & 6.01$\times$ \\*
 & PR & 64.8 & 46.3 & 25.5 & 14.3 & 10.7 & 6.08$\times$ \\*
 & $\SV{1}$ (\%) & 8.1 & 9.8 & 14.7 & 20.6 & 23.1 & 2.85$\times$ \\*
 & \quad 95\% CI & 7.6--8.6 & 9.4--10.3 & 14.2--15.5 & 19.7--21.6 & 22.3--24.1 & \\*
 & $\SV{1\text{--}10}$ (\%) & 28.7 & 35.8 & 46.9 & 60.6 & 68.4 & 2.38$\times$ \\
\midrule
\multicolumn{8}{@{}l}{\scriptsize\textbf{B\quad From the first to the last layer, all sixteen extractions}}\\
\midrule
 & & & \multicolumn{3}{c}{\scriptsize Fold $L_0/L_{11}$} & $\SV{1}$ (\%) & Rises \\
\cmidrule(lr){4-6}
Extraction & Design & $|V|$ & ER & $\mathrm{ER}_{\mathrm{var}}$ & PR & $L_0 \to L_{11}$ & ER/$\mathrm{ER}_{\mathrm{var}}$/PR \\
\midrule
lung & primary & 5{,}648 & 1.41 & 5.43 & 7.04 & 7.3 $\to$ 24.0 & 0/0/0 \\
lung (ext.) & replication & 5{,}616 & 1.35 & 4.03 & 4.94 & 7.3 $\to$ 17.9 & 0/0/1 \\
kidney~B & second tissue & 4{,}780 & 1.41 & 5.09 & 5.96 & 7.0 $\to$ 17.6 & 0/0/0 \\
immune & continuity & 2{,}039 & 1.43 & 4.86 & 5.93 & 5.6 $\to$ 18.6 & 0/0/1 \\
kidney~A & continuity & 1{,}972 & 1.47 & 6.01 & 6.08 & 8.1 $\to$ 23.1 & 0/0/1 \\
\addlinespace[2pt]
immune & cell draw, seed 43 & 2{,}080 & 1.43 & 5.17 & 6.68 & 6.0 $\to$ 21.9 & 0/0/0 \\
 & cell draw, seed 44 & 1{,}972 & 1.41 & 4.60 & 5.59 & 5.7 $\to$ 17.5 & 0/0/0 \\
 & donor TSP14 & 1{,}964 & 1.47 & 6.35 & 8.14 & 7.4 $\to$ 30.0 & 0/0/0 \\
 & donor TSP21 & 1{,}765 & 1.37 & 3.96 & 5.08 & 5.6 $\to$ 16.3 & 0/0/1 \\
 & donor TSP25 & 1{,}791 & 1.38 & 4.04 & 5.03 & 6.0 $\to$ 16.6 & 0/0/1 \\
 & B cells & 1{,}398 & 1.47 & 8.37 & 13.35 & 6.4 $\to$ 38.2 & 0/0/0 \\
 & CD4 T cells & 1{,}465 & 1.40 & 4.64 & 6.27 & 5.9 $\to$ 18.7 & 0/0/0 \\
 & CD8 T cells & 1{,}595 & 1.44 & 6.37 & 9.77 & 6.1 $\to$ 30.1 & 0/0/0 \\
\addlinespace[2pt]
kidney (192 cells) & whole-human ckpt & 4{,}923 & 1.40 & 5.08 & 5.94 & 7.1 $\to$ 16.8 & 0/0/0 \\
 & kidney ckpt & 4{,}923 & 2.49 & 23.88 & 13.04 & 8.8 $\to$ 34.0 & 0/0/0 \\
 & brain ckpt & 4{,}923 & 1.64 & 9.73 & 10.56 & 7.0 $\to$ 28.9 & 0/0/2 \\
\midrule
\multicolumn{8}{@{}l}{\scriptsize\textbf{C\quad Near-degenerate adjacent singular values, feature shuffle and the null direction}}\\
\midrule
\multicolumn{2}{@{}l}{\scriptsize } & lung & lung (ext.) & kidney~B & immune & kidney~A & All 16 \\
\midrule
\multicolumn{2}{@{}l}{\scriptsize \itshape Near-degenerate pairs, $S_k/S_{k+1}$} & & & & & & \\
\multicolumn{2}{@{}l}{\scriptsize \quad $k = 1$ ($\SV{1}/\SV{2}$)} & 0 & 0 & 0 & 4 & 0 & 8/192 (4.2\%) \\
\multicolumn{2}{@{}l}{\scriptsize \quad $k = 2$ ($\SV{2}/\SV{3}$)} & 0 & 0 & 1 & 0 & 1 & 6/192 (3.1\%) \\
\multicolumn{2}{@{}l}{\scriptsize \quad $k = 3$ ($\SV{3}/\SV{4}$)} & 1 & 1 & 0 & 1 & 0 & 24/192 (12.5\%) \\
\multicolumn{2}{@{}l}{\scriptsize \quad $k = 4$ ($\SV{4}/\SV{5}$)} & 0 & 0 & 1 & 1 & 3 & 25/192 (13.0\%) \\
\multicolumn{2}{@{}l}{\scriptsize \quad $k = 5$ ($\SV{5}/\SV{6}$)} & 0 & 0 & 2 & 5 & 1 & 42/192 (21.9\%) \\
\multicolumn{2}{@{}l}{\scriptsize \quad $k = 6$ ($\SV{6}/\SV{7}$)} & 3 & 0 & 0 & 2 & 2 & 34/192 (17.7\%) \\
\multicolumn{2}{@{}l}{\scriptsize \quad $k = 7$ ($\SV{7}/\SV{8}$)} & 0 & 2 & 2 & 2 & 4 & 31/192 (16.1\%) \\
\multicolumn{2}{@{}l}{\scriptsize \quad $k = 8$ ($\SV{8}/\SV{9}$)} & 2 & 2 & 2 & 9 & 4 & 67/192 (34.9\%) \\
\multicolumn{2}{@{}l}{\scriptsize \quad $k = 9$ ($\SV{9}/\SV{10}$)} & 5 & 3 & 6 & 3 & 3 & 60/192 (31.2\%) \\
\multicolumn{2}{@{}l}{\scriptsize \quad $k = 10$ ($\SV{10}/\SV{11}$)} & 4 & 3 & 1 & 3 & 6 & 74/192 (38.5\%) \\
\multicolumn{2}{@{}l}{\scriptsize \quad $k = 11$ ($\SV{11}/\SV{12}$)} & 5 & 1 & 5 & 2 & 3 & 58/192 (30.2\%) \\
\multicolumn{2}{@{}l}{\scriptsize \quad all pairs, $k = 1$--$11$} & 20/132 & 12/132 & 20/132 & 32/132 & 27/132 & 429/2{,}112 (20.3\%) \\
\addlinespace[2pt]
\multicolumn{2}{@{}l}{\scriptsize \itshape Feature shuffle} & & & & & & \\
\multicolumn{2}{@{}l}{\scriptsize \quad shuffled/observed ER, $L_{0}$} & 1.09$\times$ & 1.09$\times$ & 1.10$\times$ & 1.11$\times$ & 1.12$\times$ & 1.06--1.12 \\
\multicolumn{2}{@{}l}{\scriptsize \quad shuffled/observed ER, $L_{11}$} & 1.48$\times$ & 1.40$\times$ & 1.48$\times$ & 1.52$\times$ & 1.57$\times$ & 1.40--2.51 \\
\multicolumn{2}{@{}l}{\scriptsize \quad shuffled ER, $L_{11}$} & 479.6 & 478.7 & 479.8 & 470.8 & 467.5 & 454.1--479.9 \\
\addlinespace[2pt]
\multicolumn{2}{@{}l}{\scriptsize \itshape Null direction} & & & & & & \\
\multicolumn{2}{@{}l}{\scriptsize \quad singular values $< 10^{-3} S_1$} & 1 & 1 & 1 & 1 & 1 & 1 at all 192 \\
\multicolumn{2}{@{}l}{\scriptsize \quad $|\cos|$ with uniform, $L_0$/$L_{11}$} & \multicolumn{5}{c}{\scriptsize 0.9993 / 0.9991 in each of the five} & \\
\multicolumn{2}{@{}l}{\scriptsize \quad $S_{512}/S_1$ ($L_0$, $L_{11}$)} & \multicolumn{5}{c}{\scriptsize $3.3$--$7.0 \times 10^{-8}$} & \\
\midrule
\multicolumn{8}{@{}l}{\scriptsize\textbf{D\quad Detection controls, five main extractions}}\\
\midrule
\multicolumn{2}{@{}l}{\scriptsize } & lung & lung (ext.) & kidney~B & immune & kidney~A & Range \\
\midrule
\multicolumn{2}{@{}l}{\scriptsize ER, all $V$, $L_{11}$ (panel A)} & 324.1 & 340.8 & 324.1 & 309.6 & 297.1 & 297.1--340.8 \\
\multicolumn{2}{@{}l}{\scriptsize ER, matched 1{,}300 genes, $L_{0}$} & 432.0 $\pm$ 0.4 & 435.7 $\pm$ 0.3 & 433.1 $\pm$ 0.4 & 431.4 $\pm$ 0.3 & 427.4 $\pm$ 0.3 & 427.4--435.7 \\
\multicolumn{2}{@{}l}{\scriptsize ER, matched 1{,}300 genes, $L_{11}$} & 296.3 $\pm$ 1.1 & 318.3 $\pm$ 0.5 & 302.8 $\pm$ 0.8 & 299.5 $\pm$ 0.8 & 290.0 $\pm$ 0.6 & 290.0--318.3 \\
\multicolumn{2}{@{}l}{\scriptsize ER, matched, fold $L_0/L_{11}$} & 1.46$\times$ & 1.37$\times$ & 1.43$\times$ & 1.44$\times$ & 1.47$\times$ & 1.37--1.47 \\
\addlinespace[2pt]
\multicolumn{2}{@{}l}{\scriptsize genes with $n_{\mathrm{obs}} \geq 3$} & 4{,}221 & 4{,}529 & 3{,}528 & 1{,}341 & 1{,}346 & 1{,}341--4{,}529 \\
\multicolumn{2}{@{}l}{\scriptsize $n_{\mathrm{obs}} \geq 3$: ER, $L_{11}$} & 335.1 & 340.0 & 325.2 & 318.1 & 301.2 & 301.2--340.0 \\
\multicolumn{2}{@{}l}{\scriptsize $n_{\mathrm{obs}} \geq 3$: PR, $L_{11}$} & 20.5 & 17.6 & 13.4 & 20.6 & 9.3 & 9.3--20.6 \\
\multicolumn{2}{@{}l}{\scriptsize $n_{\mathrm{obs}} \geq 3$: $\SV{1}$ (\%), $L_{11}$} & 13.7 & 17.7 & 18.8 & 15.5 & 27.0 & 13.7--27.0 \\
\end{longtable}
\endgroup

\setcounter{table}{1}
\begingroup
\footnotesize
\setlength{\LTleft}{0pt}\setlength{\LTright}{0pt}
\begin{longtable}{@{}p{\textwidth}@{}}
\caption[The early compartment axis]{\textbf{The early compartment axis.} Each of the first twelve directions ($\SV{1}$--$\SV{12}$) of the observed-gene decomposition at each layer, tested for separating genes directly annotated to extracellular space (GO:0005615) or cytosol (GO:0005829) from all other observed genes. AUROC, sign-free AUROC $\max(A, 1 - A)$; $z$, AUROC standardised by the mean and SD of its permutation null; $p_{\mathrm{adj}}$, Westfall--Young max-$T$ adjusted $p$ over the 288 tests of one extraction (12 directions $\times$ 12 layers $\times$ 2 terms; 2{,}000 joint permutations of the gene labels), 0.0005 $= 1/2{,}001$ being its floor. \textbf{A}, the cell (layer, direction) with the largest AUROC for each term; $n$, observed genes annotated to the term; under each extraction, $|V|$ / genes with any GO annotation. CI, 95\% gene-bootstrap percentile interval ($B = 2{,}000$, direction held fixed), extracellular cells only. 95th percentile of the null maximum of $z$: 4.68, 4.66, 4.59, 4.67, 4.55 (row order). Strat., AUROC within deciles of $\log n_{\mathrm{obs}}$, oriented like the AUROC (own max-$T$ family, labels permuted within deciles; $p_{\mathrm{adj}}$ at the floor except kidney~A cytosol, 0.0355); $|\rho|$, absolute Spearman correlation of the coordinate with $\log n_{\mathrm{obs}}$; Annot., AUROC of the same cell among annotated genes only ($p_{\mathrm{adj}}$ at the floor); OR, Fisher odds ratio (Haldane) for the term in the 27\% pole where it is enriched (two-sided $p$) and in the opposite pole. \textbf{B}, extracellular AUROC of the compartment direction ($\SV{1}$; $\SV{2}$ in immune) at every layer; for kidney~A the row shows $\SV{1}$, which is not the lung compartment direction ($r$ and lung-matched cell in D); $p_{\mathrm{adj}}$ as in A; each layer is decomposed separately; permutation-null mean 0.510--0.520. \textbf{C}, extracellular space on one common gene set per extraction (observed genes in the scGPT vocabulary and the ESM2 file; $n$, with extracellular genes in parentheses). Probe, out-of-fold AUROC of L2 logistic regression on the first 256 principal components (5 outer folds, regularisation by inner 5-fold cross-validation), with 95\% gene-bootstrap percentile interval ($B = 2{,}000$ resamples of the out-of-fold scores); $\Delta$, paired difference from scGPT $L_3$ on the same resamples; detection-stratified probe AUROCs differ from these by $\leq 0.013$ (cytosol probes in the CSV). Own-SVD, sign-free AUROC of $\SV{1}$ of the representation's own decomposition of these genes, and the best of its first twelve directions (index; optimistic by construction). Input-token embedding, scGPT's gene-embedding table, before any cell is processed; ESM2-15B, the 5{,}120-dimensional per-protein embedding of the 15-billion-parameter protein language model ESM2. \textbf{D}, on the extracellular cell of A, oriented so that extracellular genes score high: AUROC of cytosol, mitochondrion (GO:0005739) and endoplasmic-reticulum lumen (GO:0005788) against all other genes (below 0.5, the opposite end; $p_{\mathrm{adj}}$ for cytosol as in A, for the other two over the 288 tests of these two terms, same permutations); $r$, Pearson correlation with the lung coordinate on the $n$ shared genes; lung-matched cell, the extraction's cell most correlated with the lung coordinate ($|r|$), and its extracellular AUROC oriented like the lung coordinate. Full values: \texttt{Table\_S2\_compartment.csv}.}\label{tab:S2}\\
\toprule
\endfirsthead
\multicolumn{1}{@{}l}{\scriptsize\itshape Table~\thetable\ (continued)}\\
\toprule
\endhead
\multicolumn{1}{r@{}}{\scriptsize\itshape continued on next page}\\
\endfoot
\bottomrule
\endlastfoot
\scriptsize\setlength{\tabcolsep}{2.4pt}\renewcommand{\arraystretch}{1.05}%
\begin{tabular*}{\textwidth}{@{\extracolsep{\fill}}lllrcrrrrrrr@{}}
\multicolumn{12}{@{}l}{\textbf{A\quad Best cell for each term (288 tests per extraction)}}\\
\midrule
Extraction & Term ($n$) & Cell & AUROC & 95\% CI & $z$ & $p_{\mathrm{adj}}$ & Strat. & $|\rho|$ & Annot. & OR, enriched pole ($p$) & OR, opposite \\
\midrule
lung & extracellular (591) & $L_{3}$ $\SV{1}$ & 0.773 & 0.754--0.792 & 34.5 & 0.0005 & 0.776 & 0.12 & 0.779 & 5.52 ($7.9\times10^{-81}$) & 0.13 \\
5{,}648 / 5{,}343 & cytosol (1{,}989) & $L_{4}$ $\SV{2}$ & 0.633 & -- & 26.5 & 0.0005 & 0.589 & 0.53 & 0.615 & 1.96 ($7.8\times10^{-28}$) & 0.42 \\
\addlinespace[1.5pt]
lung (ext.) & extracellular (554) & $L_{3}$ $\SV{1}$ & 0.756 & 0.736--0.777 & 30.8 & 0.0005 & 0.767 & 0.20 & 0.763 & 5.02 ($1.0\times10^{-68}$) & 0.16 \\
5{,}616 / 5{,}239 & cytosol (1{,}978) & $L_{4}$ $\SV{2}$ & 0.640 & -- & 26.2 & 0.0005 & 0.594 & 0.40 & 0.619 & 2.14 ($8.3\times10^{-35}$) & 0.37 \\
\addlinespace[1.5pt]
kidney~B & extracellular (354) & $L_{2}$ $\SV{1}$ & 0.707 & 0.681--0.733 & 20.4 & 0.0005 & 0.700 & 0.25 & 0.719 & 3.77 ($1.1\times10^{-31}$) & 0.24 \\
4{,}780 / 4{,}544 & cytosol (1{,}876) & $L_{11}$ $\SV{2}$ & 0.603 & -- & 19.2 & 0.0005 & 0.554 & 0.58 & 0.580 & 1.80 ($5.4\times10^{-19}$) & 0.51 \\
\addlinespace[1.5pt]
immune & extracellular (258) & $L_{0}$ $\SV{2}$ & 0.760 & 0.732--0.788 & 21.6 & 0.0005 & 0.761 & 0.00 & 0.752 & 4.61 ($1.5\times10^{-28}$) & 0.09 \\
2{,}039 / 1{,}389 & cytosol (427) & $L_{0}$ $\SV{1}$ & 0.714 & -- & 21.2 & 0.0005 & 0.615 & 0.64 & 0.610 & 3.52 ($9.4\times10^{-28}$) & 0.17 \\
\addlinespace[1.5pt]
kidney~A & extracellular (146) & $L_{0}$ $\SV{1}$ & 0.759 & 0.724--0.793 & 15.9 & 0.0005 & 0.804 & 0.56 & 0.658 & 3.79 ($4.2\times10^{-14}$) & 0.04 \\
1{,}972 / 1{,}352 & cytosol (466) & $L_{3}$ $\SV{1}$ & 0.740 & -- & 24.8 & 0.0005 & 0.563 & 0.80 & 0.586 & 3.25 ($2.4\times10^{-25}$) & 0.04 \\
\end{tabular*}\\
\midrule
\scriptsize\setlength{\tabcolsep}{2.4pt}\renewcommand{\arraystretch}{1.05}%
\begin{tabular*}{\textwidth}{@{\extracolsep{\fill}}ll*{12}{r}@{}}
\multicolumn{14}{@{}l}{\textbf{B\quad Extracellular AUROC of the compartment direction at every layer}}\\
\midrule
Extraction & & $L_{0}$ & $L_{1}$ & $L_{2}$ & $L_{3}$ & $L_{4}$ & $L_{5}$ & $L_{6}$ & $L_{7}$ & $L_{8}$ & $L_{9}$ & $L_{10}$ & $L_{11}$ \\
\midrule
lung & AUROC & 0.730 & 0.751 & 0.762 & 0.773 & 0.751 & 0.646 & 0.633 & 0.612 & 0.594 & 0.591 & 0.629 & 0.611 \\
\quad $\SV{1}$ & $p_{\mathrm{adj}}$ & 0.0005 & 0.0005 & 0.0005 & 0.0005 & 0.0005 & 0.0005 & 0.0005 & 0.0005 & 0.0005 & 0.0005 & 0.0005 & 0.0005 \\
\addlinespace[1.5pt]
lung (ext.) & AUROC & 0.715 & 0.725 & 0.747 & 0.756 & 0.729 & 0.667 & 0.672 & 0.625 & 0.599 & 0.556 & 0.529 & 0.548 \\
\quad $\SV{1}$ & $p_{\mathrm{adj}}$ & 0.0005 & 0.0005 & 0.0005 & 0.0005 & 0.0005 & 0.0005 & 0.0005 & 0.0005 & 0.0005 & 0.0015 & 0.9170 & 0.0370 \\
\addlinespace[1.5pt]
kidney~B & AUROC & 0.682 & 0.693 & 0.707 & 0.703 & 0.666 & 0.626 & 0.623 & 0.587 & 0.565 & 0.532 & 0.508 & 0.523 \\
\quad $\SV{1}$ & $p_{\mathrm{adj}}$ & 0.0005 & 0.0005 & 0.0005 & 0.0005 & 0.0005 & 0.0005 & 0.0005 & 0.0005 & 0.0035 & 0.9705 & 1.0000 & 1.0000 \\
\addlinespace[1.5pt]
immune & AUROC & 0.760 & 0.740 & 0.643 & 0.563 & 0.620 & 0.647 & 0.632 & 0.588 & 0.583 & 0.546 & 0.542 & 0.596 \\
\quad $\SV{2}$ & $p_{\mathrm{adj}}$ & 0.0005 & 0.0005 & 0.0005 & 0.1824 & 0.0005 & 0.0005 & 0.0005 & 0.0015 & 0.0040 & 0.8211 & 0.9465 & 0.0005 \\
\addlinespace[1.5pt]
kidney~A & AUROC & 0.759 & 0.744 & 0.714 & 0.707 & 0.722 & 0.722 & 0.721 & 0.701 & 0.690 & 0.671 & 0.659 & 0.646 \\
\quad $\SV{1}$ & $p_{\mathrm{adj}}$ & 0.0005 & 0.0005 & 0.0005 & 0.0005 & 0.0005 & 0.0005 & 0.0005 & 0.0005 & 0.0005 & 0.0005 & 0.0005 & 0.0005 \\
\end{tabular*}\\
\midrule
\scriptsize\setlength{\tabcolsep}{2.4pt}\renewcommand{\arraystretch}{1.05}%
\begin{tabular*}{\textwidth}{@{\extracolsep{\fill}}lrrrr@{}}
\multicolumn{5}{@{}l}{\textbf{C\quad Linear probe and own-decomposition AUROC for extracellular space, common gene sets}}\\
\midrule
Representation & lung & lung (ext.) & kidney~B & immune \\
 & $n$ = 5{,}413 (591) & $n$ = 5{,}306 (554) & $n$ = 4{,}602 (354) & $n$ = 1{,}554 (258) \\
\midrule
\multicolumn{5}{@{}l}{\itshape Probe, out-of-fold AUROC (95\% CI)}\\
\quad scGPT $L_0$ & 0.837 (0.819--0.855) & 0.834 (0.813--0.853) & 0.799 (0.771--0.824) & 0.832 (0.803--0.859) \\
\quad scGPT $L_3$ & 0.836 (0.818--0.854) & 0.831 (0.810--0.851) & 0.806 (0.779--0.832) & 0.824 (0.792--0.852) \\
\quad scGPT $L_{11}$ & 0.833 (0.815--0.850) & 0.829 (0.808--0.848) & 0.799 (0.772--0.825) & 0.827 (0.797--0.854) \\
\quad scGPT input-token embedding & 0.838 (0.820--0.856) & 0.834 (0.814--0.854) & 0.797 (0.770--0.824) & 0.826 (0.796--0.853) \\
\quad ESM2-15B & 0.900 (0.884--0.914) & 0.910 (0.895--0.924) & 0.868 (0.845--0.890) & 0.900 (0.877--0.921) \\
\addlinespace[1.5pt]
\multicolumn{5}{@{}l}{\itshape $\Delta$ probe AUROC from scGPT $L_3$ (95\% CI)}\\
\quad scGPT $L_0$ & $+$0.001 ($-$0.005, $+$0.006) & $+$0.002 ($-$0.003, $+$0.008) & $-$0.007 ($-$0.015, $+$0.001) & $+$0.008 ($-$0.004, $+$0.021) \\
\quad scGPT $L_{11}$ & $-$0.003 ($-$0.011, $+$0.005) & $-$0.003 ($-$0.011, $+$0.005) & $-$0.007 ($-$0.018, $+$0.004) & $+$0.003 ($-$0.010, $+$0.016) \\
\quad scGPT input-token embedding & $+$0.002 ($-$0.004, $+$0.008) & $+$0.003 ($-$0.004, $+$0.010) & $-$0.008 ($-$0.018, $+$0.002) & $+$0.002 ($-$0.009, $+$0.012) \\
\quad ESM2-15B & $+$0.064 ($+$0.048, $+$0.081) & $+$0.079 ($+$0.062, $+$0.095) & $+$0.062 ($+$0.040, $+$0.085) & $+$0.076 ($+$0.050, $+$0.105) \\
\addlinespace[1.5pt]
\multicolumn{5}{@{}l}{\itshape Own-SVD: $\SV{1}$ AUROC; best of $\SV{1}$--$\SV{12}$ (index)}\\
\quad scGPT $L_0$ & 0.746; 0.746 (1) & 0.730; 0.730 (1) & 0.694; 0.694 (1) & 0.635; 0.715 (3) \\
\quad scGPT $L_3$ & 0.779; 0.779 (1) & 0.766; 0.766 (1) & 0.724; 0.724 (1) & 0.642; 0.706 (4) \\
\quad scGPT $L_{11}$ & 0.604; 0.705 (3) & 0.528; 0.722 (2) & 0.525; 0.660 (3) & 0.542; 0.729 (6) \\
\quad scGPT input-token embedding & 0.545; 0.758 (2) & 0.530; 0.743 (2) & 0.505; 0.705 (2) & 0.516; 0.723 (3) \\
\quad ESM2-15B & 0.542; 0.758 (8) & 0.539; 0.762 (8) & 0.542; 0.685 (8) & 0.557; 0.751 (8) \\
\end{tabular*}\\
\midrule
\scriptsize\setlength{\tabcolsep}{2.4pt}\renewcommand{\arraystretch}{1.05}%
\begin{tabular*}{\textwidth}{@{\extracolsep{\fill}}llrrrrlr@{}}
\multicolumn{8}{@{}l}{\textbf{D\quad Other compartments on the extracellular cell, and replication of the direction}}\\
\midrule
 & & \multicolumn{3}{c}{Oriented AUROC ($p_{\mathrm{adj}}$)} & & \multicolumn{2}{c}{Lung-matched cell} \\
\cmidrule(lr){3-5}\cmidrule(l){7-8}
Extraction & Cell & Cytosol & Mitochondrion & ER lumen & $r$ with lung ($n$) & Cell ($|r|$) & Extracellular ($p_{\mathrm{adj}}$) \\
\midrule
lung & $L_{3}$ $\SV{1}$ & 0.392 (0.0005) & 0.400 (0.0005) & 0.667 (0.0005) & -- & -- & -- \\
lung (ext.) & $L_{3}$ $\SV{1}$ & 0.395 (0.0005) & 0.401 (0.0005) & 0.660 (0.0005) & 0.985 (4{,}535) & $L_{3}$ $\SV{1}$ (0.985) & 0.756 (0.0005) \\
kidney~B & $L_{2}$ $\SV{1}$ & 0.426 (0.0005) & 0.438 (0.0010) & 0.613 (0.0420) & 0.970 (4{,}023) & $L_{2}$ $\SV{1}$ (0.970) & 0.707 (0.0005) \\
immune & $L_{0}$ $\SV{2}$ & 0.465 (0.9330) & 0.420 (0.7406) & 0.665 (0.0315) & 0.920 (856) & $L_{1}$ $\SV{2}$ (0.929) & 0.740 (0.0005) \\
kidney~A & $L_{0}$ $\SV{1}$ & 0.690 (0.0005) & 0.578 (0.0525) & 0.798 (0.0005) & 0.118 (840) & $L_{0}$ $\SV{2}$ (0.886) & 0.568 (0.5027) \\
\end{tabular*}\\
\end{longtable}
\endgroup

\setcounter{table}{2}
\begingroup\footnotesize
\setlength{\tabcolsep}{3pt}
\begin{longtable}{@{}p{\textwidth}@{}}
\caption{\textbf{H1: physical-interaction partners share spectral coordinates (significance, nulls, hubs, co-expression, co-complex pairs, pole size and STRING confidence).} $r$, scalar assortativity of a spectral coordinate over the STRING v12 physical edges (combined score $\geq 700$) whose two ends are observed genes ($V$; $n_V$ genes, $m$ edges), i.e.\ the Pearson correlation of the coordinate across the two ends of every edge. Cell, one direction at one layer; the family is $\SV{1}$--$\SV{10}$ at $L_0$--$L_{11}$ (120 cells). MS, Maslov--Sneppen degree-preserving rewiring of the same edges (double-edge swaps; burn-in $10m$ accepted swaps, then $m$ between draws); $N$, joint null draws applied to all 120 cells; Sig., number of cells (of 120) with two-sided Westfall--Young max-$T$ adjusted $p < 0.05$ (the family minimum is at the permutation floor $1/(N+1)$ in every family shown, and every significant cell has $r > 0$); $z = (r - \text{null mean})/\text{null SD}$; strongest cell, the cell with the largest $|z|$ under MS. 95\% CI: vertex subsampling ($n_V/2$ genes without replacement, $B = 500$) with delete-$d$ scaling. \textbf{(A)} Robustness extractions: two further 256-cell immune draws, one draw per donor for three donors, three immune cell types, and one further 192-cell draw of kidney~B's cells embedded with each of three scGPT checkpoints (whole-human, kidney, brain); no interval was computed for them. \textbf{(B)} Label shuffle, permutation of coordinates across $V$; degree strata, permutation within deciles of degree in the $V$ network or in the genome-wide STRING network; detection strata, permutation within deciles of $\log n_{\mathrm{obs}}$, the number of sampled cells that contributed to a gene. \textbf{(C)} The 5\% or 10\% most connected genes of $V$ are removed with their edges; genes, number removed; edges kept, fraction of the $m$ edges remaining; max.\ change, largest decrease of $r$ when one gene and its edges are removed; jackknife SE, delete-1 jackknife standard error of $r$. \textbf{(D)} Upper part: every STRING edge is compared with five gene pairs of $V$ that have no physical STRING link at any score, drawn from the same decile of $|$co-expression$|$ (Pearson correlation of log-normalised expression over the extraction's cells); AUROC stratified by decile, for three scores: cosine of the two genes' centred 512-dimensional vectors at the strongest cell's layer, $-|\Delta x|$ (negative absolute difference of the two genes' coordinates on the strongest cell), and $|$co-expression$|$ itself (the co-expression left after matching); 95\% CI as above; mean $|$coexpr$|$, over edges / matched non-edges; 50 bins, the same with 50 quantile bins (immune: 5,687 of 6,255 non-edges available); 3,000 cells, matching on co-expression computed on 3,000 cells of the same population. Lower part: co-complex pairs within $V$ (ComplexPortal; hu.MAP~2.0; each also without the pairs that are STRING physical $\geq 700$ edges) as the edge set of the same 120-cell family; $r$ and $z$ at the STRING strongest cell. \textbf{(E)} Co-pole rate, fraction of edges whose two ends fall in the same pole (top or bottom $K$ genes of the coordinate, $K = \lfloor (K/n)\,n_V \rfloor$); ratio, observed rate / null mean ($N$ as in A). Upper part: MS null; the interval is for the ratio of the observed rate to its configuration-model expectation, which differs from the MS null mean by 0.1--0.9\% in these cells. Lower part: $K/n = 0.27$ under both nulls. \textbf{(F)} $\rho$, Spearman correlation of a link's STRING score with the cosine of its two genes' centred, L2-normalised 512-dimensional vectors; $\rho_{\mathrm{adj}}$, partial Spearman correlation controlling for the ranks of the summed log degree (score $\geq 150$ network within $V$), of the summed $\log(1 + n_{\mathrm{obs}})$ and of $|$co-expression$|$ (3,000 cells); 95\% CI as above. Band AUROC, links of each combined-score band against one non-linked pair per link (no physical link at any score) matched on the degree deciles of both ends, drawn once (no interval). Rows `all 12': ranges over layers; for $\rho_{\mathrm{adj}}$ the smallest lower confidence bound over layers in brackets; monotone, layers at which the four band AUROCs increase strictly. MS chains: mean acceptance 0.90--0.96, lag-1 autocorrelation of the strongest cell's null statistic 0.09--0.14. Vertex-subsampling intervals with the pre-registered scaling factor $\sqrt{m/n}$, which omits the finite-population term and so gives half-widths smaller by about a factor $\sqrt{2}$, are in the CSV (columns ending \texttt{\_prereg}). Full values: \texttt{Table\_S3\_interaction.csv}.}\label{tab:S3}\\
\toprule
\endfirsthead
\multicolumn{1}{@{}l}{\footnotesize\itshape Table~\thetable\ (continued)}\\
\toprule
\endhead
\midrule
\multicolumn{1}{@{}r}{\footnotesize\itshape continued on next page}\\
\endfoot
\bottomrule
\endlastfoot
\textbf{A}\enspace H1 family against Maslov--Sneppen rewiring (MS), per extraction\\*
{\scriptsize\begin{tabular*}{\textwidth}{@{\extracolsep{\fill}}lrrrrlrcrrr@{}}
Extraction & $n_V$ & $m$ & $N$ & Sig. & Strongest cell & $r$ & 95\% CI & Null mean & $z$ & Median $r$ \\
\midrule
lung & 5,648 & 14,474 & 10,000 & 120 & $L_{3}\,\SV{1}$ & 0.435 & 0.388--0.487 & $-$0.001 & 53.1 & 0.181 \\
lung (ext.) & 5,616 & 14,132 & 10,000 & 120 & $L_{8}\,\SV{2}$ & 0.406 & 0.343--0.462 & $-$0.001 & 49.2 & 0.170 \\
kidney~B & 4,780 & 12,434 & 2,000 & 120 & $L_{9}\,\SV{3}$ & 0.404 & 0.351--0.454 & $-$0.001 & 45.9 & 0.209 \\
immune & 2,039 & 1,251 & 2,000 & 118 & $L_{4}\,\SV{6}$ & 0.566 & 0.465--0.688 & $-$0.003 & 20.9 & 0.318 \\
kidney~A & 1,972 & 1,390 & 2,000 & 120 & $L_{5}\,\SV{2}$ & 0.702 & 0.602--0.809 & $-$0.011 & 28.2 & 0.353 \\
\addlinespace[2pt]\multicolumn{11}{@{}l}{\itshape Robustness extractions (MS, $N = 1{,}000$)}\\
immune, draw (seed 43) & 2,080 & 1,293 & 1,000 & 117 & $L_{6}\,\SV{5}$ & 0.562 & -- & $-$0.004 & 20.7 & 0.266 \\
immune, draw (seed 44) & 1,972 & 1,263 & 1,000 & 117 & $L_{4}\,\SV{6}$ & 0.577 & -- & $-$0.005 & 20.3 & 0.306 \\
immune, donor TSP14 & 1,964 & 1,030 & 1,000 & 111 & $L_{6}\,\SV{7}$ & 0.559 & -- & $-$0.004 & 19.0 & 0.294 \\
immune, donor TSP21 & 1,765 & 1,050 & 1,000 & 109 & $L_{2}\,\SV{2}$ & 0.525 & -- & $-$0.005 & 16.9 & 0.247 \\
immune, donor TSP25 & 1,791 & 1,064 & 1,000 & 107 & $L_{2}\,\SV{2}$ & 0.521 & -- & $-$0.005 & 16.8 & 0.228 \\
immune, B cells & 1,398 & 615 & 1,000 & 102 & $L_{2}\,\SV{2}$ & 0.547 & -- & $-$0.007 & 14.0 & 0.265 \\
immune, CD4 T cells & 1,465 & 749 & 1,000 & 109 & $L_{1}\,\SV{2}$ & 0.502 & -- & $-$0.003 & 14.2 & 0.251 \\
immune, CD8 T cells & 1,595 & 896 & 1,000 & 109 & $L_{3}\,\SV{3}$ & 0.588 & -- & $-$0.004 & 18.4 & 0.257 \\
kidney cells, whole-human ckpt. & 4,923 & 12,689 & 1,000 & 120 & $L_{9}\,\SV{3}$ & 0.383 & -- & 0.000 & 44.2 & 0.216 \\
kidney cells, kidney ckpt. & 4,923 & 12,689 & 1,000 & 120 & $L_{2}\,\SV{3}$ & 0.418 & -- & $-$0.001 & 47.4 & 0.187 \\
kidney cells, brain ckpt. & 4,923 & 12,689 & 1,000 & 120 & $L_{1}\,\SV{4}$ & 0.302 & -- & $-$0.002 & 35.3 & 0.147 \\
\end{tabular*}}\\
\midrule
\textbf{B}\enspace Null mean of $r$ at the strongest cell under each null ($N = 10{,}000$)\\*
{\scriptsize\begin{tabular*}{\textwidth}{@{\extracolsep{\fill}}lrrrrrrrr@{}}
& \multicolumn{4}{c}{lung, $L_3\,\SV{1}$ ($r = 0.435$)} & \multicolumn{4}{c}{lung (ext.), $L_8\,\SV{2}$ ($r = 0.406$)} \\
\cmidrule(lr){2-5}\cmidrule(l){6-9}
Null & Null mean & Null SD & $z$ & Sig. & Null mean & Null SD & $z$ & Sig. \\
\midrule
Label shuffle over $V$ & $-$0.001 & 0.008 & 53.1 & 120 & $-$0.001 & 0.008 & 49.0 & 120 \\
Degree strata ($V$ network) & 0.005 & 0.009 & 50.4 & 120 & 0.007 & 0.009 & 45.3 & 120 \\
Degree strata (genome-wide STRING) & 0.021 & 0.010 & 41.0 & 120 & 0.025 & 0.010 & 36.5 & 120 \\
Detection strata ($\log n_{\mathrm{obs}}$) & 0.000 & 0.008 & 52.0 & 117 & 0.005 & 0.009 & 45.5 & 120 \\
MS rewiring (primary) & $-$0.001 & 0.008 & 53.1 & 120 & $-$0.001 & 0.008 & 49.2 & 120 \\
\end{tabular*}}\\
\midrule
\textbf{C}\enspace Hub removal (MS null on the remaining graph, $N = 2{,}000$) and leave-one-gene-out influence, strongest cell\\*
{\scriptsize\begin{tabular*}{\textwidth}{@{\extracolsep{\fill}}llrrrrrrrrrrr@{}}
& & & \multicolumn{4}{c}{Top 5\% by degree removed} & \multicolumn{4}{c}{Top 10\% by degree removed} & \multicolumn{2}{c}{Leave one gene out} \\
\cmidrule(lr){4-7}\cmidrule(lr){8-11}\cmidrule(l){12-13}
Extraction & Cell & $r$ & Genes & Edges kept & $r$ & $z$ & Genes & Edges kept & $r$ & $z$ & Max.\ change & Jackknife SE \\
\midrule
lung & $L_{3}\,\SV{1}$ & 0.435 & 283 & 0.42 & 0.515 & 40.4 & 565 & 0.25 & 0.511 & 30.7 & 0.0045 & 0.025 \\
lung (ext.) & $L_{8}\,\SV{2}$ & 0.406 & 281 & 0.42 & 0.518 & 39.5 & 562 & 0.25 & 0.487 & 28.5 & 0.0042 & 0.028 \\
kidney~B & $L_{9}\,\SV{3}$ & 0.404 & 239 & 0.44 & 0.499 & 36.4 & 478 & 0.26 & 0.488 & 27.9 & 0.0036 & 0.026 \\
immune & $L_{4}\,\SV{6}$ & 0.566 & 102 & 0.29 & 0.537 & 10.0 & 204 & 0.12 & 0.528 & 6.7 & 0.0141 & 0.052 \\
kidney~A & $L_{5}\,\SV{2}$ & 0.702 & 99 & 0.32 & 0.569 & 11.9 & 198 & 0.12 & 0.512 & 6.9 & 0.0075 & 0.050 \\
\end{tabular*}}\\
\midrule
\textbf{D}\enspace STRING edges against co-expression-matched non-edges (stratified AUROC), and co-complex pairs as edge sets\\*
{\scriptsize\begin{tabular*}{\textwidth}{@{\extracolsep{\fill}}llcccccccc@{}}
& & & \multicolumn{3}{c}{Deciles of $|$co-expression$|$, extraction cells} & \multicolumn{2}{c}{50 bins} & \multicolumn{2}{c}{3,000 cells} \\
\cmidrule(lr){4-6}\cmidrule(lr){7-8}\cmidrule(l){9-10}
Extraction & Cell & Mean $|$coexpr$|$ & Cosine & $-|\Delta x|$ & $|$coexpr$|$ & Cosine & $|$coexpr$|$ & Cosine & $|$coexpr$|$ \\
\midrule
lung & $L_{3}\,\SV{1}$ & 0.125/0.108 & 0.725 (0.706--0.743) & 0.631 (0.612--0.647) & 0.560 (0.539--0.579) & 0.740 & 0.555 & 0.725 & 0.561 \\
lung (ext.) & $L_{8}\,\SV{2}$ & 0.085/0.077 & 0.679 (0.662--0.697) & 0.617 (0.600--0.634) & 0.534 (0.519--0.549) & 0.690 & 0.527 & 0.690 & 0.551 \\
kidney~B & $L_{9}\,\SV{3}$ & 0.121/0.111 & 0.681 (0.660--0.699) & 0.625 (0.610--0.641) & 0.538 (0.521--0.554) & 0.695 & 0.530 & 0.683 & 0.542 \\
immune & $L_{4}\,\SV{6}$ & 0.216/0.158 & 0.810 (0.770--0.855) & 0.577 (0.525--0.625) & 0.663 (0.617--0.715) & 0.800 & 0.553 & 0.803 & 0.658 \\
kidney~A & $L_{5}\,\SV{2}$ & 0.282/0.269 & 0.873 (0.838--0.912) & 0.733 (0.687--0.784) & 0.567 (0.529--0.604) & 0.860 & 0.528 & 0.872 & 0.554 \\
\end{tabular*}}\\
\addlinespace[4pt]
{\scriptsize\begin{tabular*}{\textwidth}{@{\extracolsep{\fill}}lrrrrrrrr@{}}
& \multicolumn{4}{c}{lung, $L_3\,\SV{1}$} & \multicolumn{4}{c}{lung (ext.), $L_8\,\SV{2}$} \\
\cmidrule(lr){2-5}\cmidrule(l){6-9}
Co-complex edge set (MS, $N = 2{,}000$) & Pairs & $r$ & $z$ & Sig. & Pairs & $r$ & $z$ & Sig. \\
\midrule
ComplexPortal & 1,934 & 0.588 & 26.1 & 119 & 1,863 & 0.524 & 23.3 & 120 \\
ComplexPortal, not STRING $\geq 700$ & 246 & 0.423 & 7.0 & 79 & 220 & 0.278 & 4.6 & 38 \\
hu.MAP 2.0 & 6,782 & 0.265 & 21.5 & 113 & 6,639 & 0.246 & 20.0 & 119 \\
hu.MAP 2.0, not STRING $\geq 700$ & 5,159 & 0.196 & 14.1 & 95 & 5,039 & 0.172 & 12.0 & 95 \\
\end{tabular*}}\\
\midrule
\textbf{E}\enspace Co-pole rate as a function of pole size, and label shuffle against rewiring at $K/n = 0.27$\\*
{\scriptsize\begin{tabular*}{\textwidth}{@{\extracolsep{\fill}}llccccc@{}}
& & \multicolumn{5}{c}{Observed / MS-null co-pole rate (95\% CI) at pole size $K/n$} \\
\cmidrule(l){3-7}
Extraction & Cell & 0.10 & 0.20 & 0.27 & 0.33 & 0.50 \\
\midrule
lung & $L_{3}\,\SV{1}$ & 2.27 (1.84--2.58) & 1.88 (1.69--2.05) & 1.70 (1.57--1.82) & 1.60 (1.50--1.69) & 1.32 (1.25--1.37) \\
lung (ext.) & $L_{8}\,\SV{2}$ & 2.10 (1.71--2.37) & 1.85 (1.66--2.01) & 1.70 (1.58--1.82) & 1.56 (1.45--1.66) & 1.30 (1.23--1.35) \\
kidney~B & $L_{9}\,\SV{3}$ & 2.06 (1.72--2.35) & 1.86 (1.68--2.03) & 1.66 (1.53--1.77) & 1.55 (1.44--1.65) & 1.26 (1.21--1.31) \\
\end{tabular*}}\\
\addlinespace[4pt]
{\scriptsize\begin{tabular*}{\textwidth}{@{\extracolsep{\fill}}llrrrrrrrrr@{}}
& & & \multicolumn{4}{c}{Label shuffle} & \multicolumn{4}{c}{MS rewiring} \\
\cmidrule(lr){4-7}\cmidrule(l){8-11}
Extraction & Cell & Observed & Null mean & Null SD & Ratio & $z$ & Null mean & Null SD & Ratio & $z$ \\
\midrule
lung & $L_{3}\,\SV{1}$ & 0.256 & 0.146 & 0.0074 & 1.76 & 14.9 & 0.150 & 0.0024 & 1.70 & 44.2 \\
lung (ext.) & $L_{8}\,\SV{2}$ & 0.258 & 0.146 & 0.0076 & 1.77 & 14.9 & 0.152 & 0.0024 & 1.70 & 44.1 \\
kidney~B & $L_{9}\,\SV{3}$ & 0.248 & 0.146 & 0.0077 & 1.71 & 13.4 & 0.150 & 0.0025 & 1.66 & 39.3 \\
kidney~B & $L_{11}\,\SV{1}$ & 0.116 & 0.146 & 0.0078 & 0.79 & $-$3.9 & 0.082 & 0.0021 & 1.40 & 15.8 \\
immune & $L_{4}\,\SV{6}$ & 0.344 & 0.145 & 0.0176 & 2.37 & 11.3 & 0.182 & 0.0088 & 1.88 & 18.4 \\
kidney~A & $L_{5}\,\SV{2}$ & 0.574 & 0.146 & 0.0184 & 3.93 & 23.2 & 0.284 & 0.0099 & 2.02 & 29.3 \\
kidney~A & $L_{2}\,\SV{2}$ & 0.614 & 0.145 & 0.0180 & 4.23 & 26.1 & 0.294 & 0.0101 & 2.09 & 31.8 \\
kidney~A & $L_{4}\,\SV{2}$ & 0.563 & 0.146 & 0.0181 & 3.87 & 23.1 & 0.275 & 0.0097 & 2.05 & 29.6 \\
kidney~A & $L_{0}\,\SV{1}$ & 0.212 & 0.145 & 0.0190 & 1.46 & 3.5 & 0.148 & 0.0063 & 1.43 & 10.1 \\
kidney~A & $L_{11}\,\SV{1}$ & 0.403 & 0.145 & 0.0187 & 2.77 & 13.8 & 0.255 & 0.0070 & 1.58 & 21.2 \\
\end{tabular*}}\\
\midrule
\textbf{F}\enspace STRING confidence gradient (physical links, combined score $\geq 150$, within $V$)\\*
{\scriptsize\begin{tabular*}{\textwidth}{@{\extracolsep{\fill}}lrlccccccr@{}}
& & & & & \multicolumn{4}{c}{AUROC by STRING combined-score band} & \\
\cmidrule(lr){6-9}
Extraction & Links & Layer & $\rho$ (95\% CI) & $\rho_{\mathrm{adj}}$ (95\% CI) & 150--399 & 400--699 & 700--899 & 900--999 & Monotone \\
\midrule
lung & 132,077 & $L_{0}$ & 0.098 (0.081--0.116) & 0.091 (0.075--0.108) & 0.569 & 0.597 & 0.642 & 0.692 &  \\
 &  & $L_{11}$ & 0.050 (0.032--0.069) & 0.051 (0.036--0.067) & 0.545 & 0.558 & 0.584 & 0.630 &  \\
 &  & all 12 & 0.050--0.102 & 0.051--0.096 [0.036] & 0.545--0.569 & 0.558--0.600 & 0.584--0.643 & 0.630--0.698 & 12/12 \\
\addlinespace[2pt]
lung (ext.) & 127,663 & $L_{0}$ & 0.092 (0.073--0.110) & 0.096 (0.078--0.112) & 0.568 & 0.590 & 0.643 & 0.683 &  \\
 &  & $L_{11}$ & 0.051 (0.032--0.069) & 0.060 (0.044--0.076) & 0.550 & 0.561 & 0.596 & 0.629 &  \\
 &  & all 12 & 0.051--0.095 & 0.060--0.099 [0.044] & 0.550--0.568 & 0.561--0.592 & 0.596--0.644 & 0.629--0.686 & 12/12 \\
\addlinespace[2pt]
kidney~B & 112,823 & $L_{0}$ & 0.097 (0.078--0.117) & 0.097 (0.079--0.115) & 0.556 & 0.584 & 0.639 & 0.694 &  \\
 &  & $L_{11}$ & 0.043 (0.023--0.063) & 0.046 (0.028--0.064) & 0.544 & 0.551 & 0.587 & 0.633 &  \\
 &  & all 12 & 0.043--0.099 & 0.046--0.098 [0.028] & 0.542--0.557 & 0.551--0.586 & 0.587--0.640 & 0.633--0.696 & 12/12 \\
\addlinespace[2pt]
immune & 7,803 & $L_{0}$ & 0.147 (0.101--0.195) & 0.111 (0.066--0.157) & 0.624 & 0.637 & 0.736 & 0.767 &  \\
 &  & $L_{11}$ & 0.092 (0.038--0.145) & 0.055 (0.006--0.104) & 0.590 & 0.585 & 0.636 & 0.691 &  \\
 &  & all 12 & 0.092--0.153 & 0.055--0.115 [0.006] & 0.589--0.624 & 0.584--0.641 & 0.636--0.742 & 0.691--0.768 & 8/12 \\
\addlinespace[2pt]
kidney~A & 7,938 & $L_{0}$ & 0.188 (0.120--0.264) & 0.190 (0.128--0.244) & 0.661 & 0.659 & 0.749 & 0.820 &  \\
 &  & $L_{11}$ & 0.085 (0.007--0.169) & 0.115 (0.050--0.178) & 0.647 & 0.613 & 0.699 & 0.761 &  \\
 &  & all 12 & 0.085--0.188 & 0.115--0.190 [0.050] & 0.636--0.663 & 0.606--0.659 & 0.691--0.750 & 0.758--0.822 & 0/12 \\
\end{tabular*}}\\
\end{longtable}
\endgroup

\setcounter{table}{3}
\begin{table}[!htbp]
\centering\footnotesize
\caption{\textbf{Step-by-step bridge on kidney~A from the screen's interaction co-pole statistic to the H1 statistic.} Each step changes one factor of the previous step (cumulatively, except the off-path check S1$'$, shown in italics, which branches from S1). Every step evaluates 36 cells, $\SV{2}$--$\SV{4}$ at $L_0$--$L_{11}$ of that step's decomposition (at S0 the zero-row SVD, whose $\SV{2}$--$\SV{4}$ correspond to observed $\SV{1}$--$\SV{3}$; from S1 the observed-gene SVD), with one null draw applied jointly to all 36. Co-pole rate, fraction of pairs whose two genes fall in the same pole (the top or bottom $K = \lfloor 0.27\,n \rfloor$ genes of the coordinate, $n$ the number of genes in the universe); label shuffle, permutation of coordinates (or pole labels) across the universe; MS, Maslov--Sneppen degree-preserving rewiring; $r$, scalar assortativity of the coordinate over the pairs. Zero-row SVD, SVD of the full 4,803-gene matrix including the all-zero rows of unobserved genes. Genes/pairs, size of the gene universe and number of gene pairs; $N$, null draws. Cells, number of the 36 with raw $p < 0.05$ / raw $p < 0.05/183$ (the screen's Bonferroni threshold over its 183 hypotheses; n.r., not resolvable because $1/(N+1) > 0.05/183$) / max-$T$ adjusted $p < 0.05$ over the 36 cells (descriptive). The minimum adjusted $p$ is at the permutation floor $1/(N+1)$ at every step except S1$'$ (0.017); at S5, 36 of the 36 cells are at the floor. Tracked direction: $L_5\,\SV{3}$ of the zero-row SVD at S0, which is $L_5\,\SV{2}$ of the observed-gene SVD (the zero-row SVD's leading direction separates unobserved from observed genes and shifts the low indices by one; observed-gene $L_5\,\SV{3}$ at S1 has a co-pole rate of 0.135, $z = -0.25$). Value, statistic of that direction; null, null mean; ratio, value / null mean (co-pole steps); $z = (\text{value} - \text{null mean})/\text{null SD}$; best cell, the cell of the 36 with the largest $z$ (largest $|z|$ for two-sided steps). S0 does not use the screen's own decomposition, of a 209-gene subset that includes 22 zero rows and maps only weakly onto the observed-gene basis (Table~S14); the screen's values on that subset are in Table~S13 (claim C21). At S3 the co-pole null mean of the tracked direction under MS rewiring is about twice its label-shuffle value, so the label-shuffle ratio overstates the enrichment about twofold (Table~S3, panel E); S4 and S5 use assortativity, whose null mean is close to zero under both nulls. Full values: \texttt{Table\_S4\_bridge.csv}.}\label{tab:S4}
\scriptsize\setlength{\tabcolsep}{2pt}
\begin{tabular*}{\textwidth}{@{\extracolsep{\fill}}l>{\raggedright\arraybackslash}p{3.65cm}r>{\raggedright\arraybackslash}p{1.6cm}rclrrrrlrr@{}}
\toprule
& & & & & & \multicolumn{5}{c}{Tracked direction} & \multicolumn{3}{c}{Best cell of 36} \\
\cmidrule(lr){7-11}\cmidrule(l){12-14}
Step & Change & Genes/pairs & Statistic; null & $N$ & Cells & Cell & Value & Null & Ratio & $z$ & Cell & Value & $z$ \\
\midrule
S0 & The screen's ingredients (the 187 observed genes of its curated list; functional STRING pairs, score $\geq 0.7$, among them), with coordinates from the zero-row SVD & 187/809 & co-pole rate; label shuffle, one-sided & 10,000 & 23/13/15 & $L_{5}\,\SV{3}$ & 0.333 & 0.141 & 2.36 & 7.5 & $L_{5}\,\SV{3}$ & 0.333 & 7.5 \\
S1 & SVD of the observed-gene matrix ($V$, 1,972 genes); coordinates still restricted to the 187 genes & 187/809 & co-pole rate; label shuffle, one-sided & 10,000 & 13/10/11 & $L_{5}\,\SV{2}$ & 0.333 & 0.141 & 2.36 & 7.4 & $L_{5}\,\SV{2}$ & 0.333 & 7.4 \\
\itshape S1$'$ & \itshape Off-path check: S1 with physical STRING $\geq 700$ pairs among the 187 genes & \itshape 187/163 & \itshape co-pole rate; label shuffle, one-sided & \itshape 10,000 & \itshape 9/0/3 & \itshape $L_{5}\,\SV{2}$ & \itshape 0.166 & \itshape 0.141 & \itshape 1.18 & \itshape 0.7 & \itshape $L_{11}\,\SV{2}$ & \itshape 0.276 & \itshape 3.9 \\
S2 & Universe $V$: poles and shuffle over all 1,972 observed genes; same 809 functional pairs & 1,972/809 & co-pole rate; label shuffle, one-sided & 10,000 & 19/13/15 & $L_{5}\,\SV{2}$ & 0.600 & 0.146 & 4.12 & 13.9 & $L_{6}\,\SV{2}$ & 0.656 & 15.7 \\
S3 & Physical STRING $\geq 700$ edges within $V$ instead of functional pairs & 1,972/1,390 & co-pole rate; label shuffle, one-sided & 10,000 & 29/23/26 & $L_{5}\,\SV{2}$ & 0.574 & 0.145 & 3.95 & 23.3 & $L_{2}\,\SV{2}$ & 0.614 & 25.3 \\
S4 & Scalar assortativity $r$ instead of the co-pole rate & 1,972/1,390 & $r$; label shuffle, two-sided & 10,000 & 36/36/36 & $L_{5}\,\SV{2}$ & 0.702 & $-$0.004 & -- & 27.7 & $L_{5}\,\SV{2}$ & 0.702 & 27.7 \\
S5 & Maslov--Sneppen rewiring instead of the label shuffle (the primary H1 run of kidney~A) & 1,972/1,390 & $r$; MS rewiring, two-sided & 2,000 & 36/n.r./36 & $L_{5}\,\SV{2}$ & 0.702 & $-$0.011 & -- & 28.2 & $L_{5}\,\SV{2}$ & 0.702 & 28.2 \\
\bottomrule
\end{tabular*}
\end{table}

\setcounter{table}{4}
\begingroup\footnotesize
\setlength{\tabcolsep}{3pt}
\begin{longtable}{@{}p{\textwidth}@{}}
\caption{\textbf{H1 conditioned on compartment co-membership and on the detection threshold.} Lung, lung (ext.) and kidney~B; $r$, scalar assortativity of a spectral coordinate over the STRING physical edges (combined score $\geq 700$) among observed genes, as in Table~S3; MS, Maslov--Sneppen degree-preserving rewiring; Sig., number of the 120 cells ($\SV{1}$--$\SV{10}$ $\times$ $L_0$--$L_{11}$) with two-sided max-$T$ adjusted $p < 0.05$ (the family minimum is at the permutation floor $1/(N+1)$ in every family; every significant cell has $r > 0$). 95\% CI: vertex subsampling ($n_V/2$ genes, $B = 500$) with delete-$d$ scaling, coordinates and residuals held fixed. \textbf{(A)} Each gene is assigned one class from its direct GO cellular-component annotations, in the priority order listed (extracellular space GO:0005615, plasma membrane GO:0005886, endoplasmic reticulum, ER lumen or Golgi, mitochondrion, nucleus, cytosol, any other direct GO annotation, none). Edge share, fraction of all STRING edges with both genes in that class; expected, its stub-matching (configuration-model) expectation; MS null, the same-class share in $N = 2{,}000$ rewired graphs. \textbf{(B)} Cells, the three cells with the largest $|z|$ under MS in Table~S3 ($^{*}$lung $L_7\,\SV{6}$, a cell with little compartment content, chosen post hoc). MS null, null mean of $r$ in an MS replica with the same $N$; permutation within classes, one permutation of coordinates within each class applied to all 120 cells (preserves class, not degree); class-pair-preserving MS, double-edge swaps accepted only if every edge keeps its pair of classes (preserves every degree and the multiset of edge class pairs; acceptance 0.360--0.363, lag-1 autocorrelation of the strongest cell's null statistic 0.18--0.23); retained, $(r - \text{null mean})/(r - \text{MS null mean})$, the share of the excess over MS that remains (no interval). Rows `120 cells': largest null mean over cells; Sig.; median (minimum) retained share. \textbf{(C)} Ordinary least squares of each coordinate on the regressors listed; entries are residual $r$ / raw $r$ at the strongest cell (with 95\% CI for the two tested designs) and the median over the 120 cells; Sig. for the residual coordinates. Class factor, seven class indicators and an intercept; 30 GO-CC/BP/MF terms, indicators of the 30 most frequent direct cellular-component, biological-process or molecular-function annotations among the genes of $V$; 8 class-defining terms, the eight class definitions of panel A as overlapping indicators; random gene sets, 30 sets of the same sizes as the CC terms; for these, ranges over the five seeds (two decimals). The 30-CC design explains 20--26\% of the strongest coordinate's variance. Share of STRING edges whose two genes share at least one term: CC 0.89--0.90 (random pairs 0.41--0.45), BP 0.39--0.41 (0.06), MF 0.92 (0.54--0.59; its most frequent term, `protein binding', is largely interaction-derived). \textbf{(D)} STRING edges split into those whose two genes share a class and the rest, each evaluated with its own MS chain (which does not keep edges within or between classes); class-centred, coordinates centred within each class before computing $r$. \textbf{(E)} $V$ restricted to genes with $n_{\mathrm{obs}}$ (number of sampled cells contributing to the gene) at or above the threshold, SVD recomputed on the restricted genes and edges among them; rows `all ($V$)', the unrestricted analysis (Table~S3; MS $N = 10{,}000$ for lung and lung (ext.), interval from the draws of this analysis); strongest direction of the full-$V$ analysis, the restricted direction at the same layer with the largest $|$correlation$|$ of coordinates on the kept genes. Full values: \texttt{Table\_S5\_compartment\_conditioning.csv}.}\label{tab:S5}\\
\toprule
\endfirsthead
\multicolumn{1}{@{}l}{\footnotesize\itshape Table~\thetable\ (continued)}\\
\toprule
\endhead
\midrule
\multicolumn{1}{@{}r}{\footnotesize\itshape continued on next page}\\
\endfoot
\bottomrule
\endlastfoot
\textbf{A}\enspace Compartment classes and the share of STRING edges whose two genes fall in the same class\\*
{\scriptsize\begin{tabular*}{\textwidth}{@{\extracolsep{\fill}}lrrrrrrrrr@{}}
& \multicolumn{3}{c}{lung} & \multicolumn{3}{c}{lung (ext.)} & \multicolumn{3}{c}{kidney~B} \\
\cmidrule(lr){2-4}\cmidrule(lr){5-7}\cmidrule(l){8-10}
Class & Genes & Edge share & Expected & Genes & Edge share & Expected & Genes & Edge share & Expected \\
\midrule
Extracellular space & 591 & 0.043 & 0.016 & 554 & 0.041 & 0.014 & 354 & 0.023 & 0.009 \\
Plasma membrane & 1,392 & 0.164 & 0.088 & 1,370 & 0.161 & 0.086 & 1,137 & 0.146 & 0.072 \\
ER or Golgi & 368 & 0.011 & 0.004 & 359 & 0.011 & 0.004 & 326 & 0.012 & 0.004 \\
Mitochondrion & 284 & 0.010 & 0.002 & 283 & 0.010 & 0.002 & 312 & 0.018 & 0.004 \\
Nucleus & 1,362 & 0.200 & 0.112 & 1,378 & 0.203 & 0.117 & 1,253 & 0.222 & 0.136 \\
Cytosol & 509 & 0.019 & 0.007 & 516 & 0.020 & 0.007 & 493 & 0.020 & 0.007 \\
Other annotated & 837 & 0.008 & 0.003 & 779 & 0.007 & 0.003 & 669 & 0.007 & 0.003 \\
Unannotated & 305 & 0.000 & 0.000 & 377 & 0.000 & 0.000 & 236 & 0.000 & 0.000 \\
\midrule
All classes (same-class edges) & 5,648 & 0.455 & 0.231 & 5,616 & 0.451 & 0.233 & 4,780 & 0.448 & 0.235 \\
MS null: mean (SD); $z$ &  & 0.231 (0.0031) & $z = 72.7$ &  & 0.232 (0.0032) & $z = 69.4$ &  & 0.234 (0.0033) & $z = 63.8$ \\
\end{tabular*}}\\
\midrule
\textbf{B}\enspace Compartment-aware nulls ($N = 2{,}000$) at the three strongest cells of each extraction\\*
{\scriptsize\begin{tabular*}{\textwidth}{@{\extracolsep{\fill}}llrrrrrrrrrr@{}}
& & & & \multicolumn{4}{c}{Permutation within classes} & \multicolumn{4}{c}{Class-pair-preserving MS} \\
\cmidrule(lr){5-8}\cmidrule(l){9-12}
Extraction & Cell & $r$ & MS null & Null mean & Null SD & $z$ & Retained & Null mean & Null SD & $z$ & Retained \\
\midrule
lung & $L_{3}\,\SV{1}$ & 0.435 & $-$0.001 & 0.075 & 0.0139 & 25.9 & 0.825 & 0.083 & 0.0077 & 46.0 & 0.807 \\
 & $L_{6}\,\SV{2}$ & 0.418 & $-$0.001 & 0.066 & 0.0136 & 25.8 & 0.839 & 0.074 & 0.0078 & 44.0 & 0.819 \\
 & $L_{8}\,\SV{3}$ & 0.418 & $-$0.001 & 0.067 & 0.0135 & 26.1 & 0.839 & 0.094 & 0.0077 & 42.1 & 0.774 \\
 & $L_{7}\,\SV{6}$$^{*}$ & 0.311 & $-$0.001 & 0.005 & 0.0090 & 33.7 & 0.979 & 0.013 & 0.0082 & 36.3 & 0.954 \\
 & 120 cells & -- & -- & max 0.075 &  & 119/120 & 0.970 (0.755) & max 0.094 &  & 120/120 & 0.957 (0.753) \\
\addlinespace[2pt]
lung (ext.) & $L_{8}\,\SV{2}$ & 0.406 & $-$0.001 & 0.064 & 0.0134 & 25.6 & 0.840 & 0.078 & 0.0079 & 41.5 & 0.806 \\
 & $L_{3}\,\SV{1}$ & 0.407 & $-$0.001 & 0.065 & 0.0136 & 25.1 & 0.837 & 0.069 & 0.0080 & 42.4 & 0.829 \\
 & $L_{7}\,\SV{2}$ & 0.393 & $-$0.001 & 0.063 & 0.0134 & 24.5 & 0.837 & 0.069 & 0.0080 & 40.6 & 0.823 \\
 & 120 cells & -- & -- & max 0.065 &  & 120/120 & 0.958 (0.837) & max 0.081 &  & 120/120 & 0.956 (0.785) \\
\addlinespace[2pt]
kidney~B & $L_{9}\,\SV{3}$ & 0.404 & $-$0.001 & 0.056 & 0.0136 & 25.6 & 0.860 & 0.075 & 0.0084 & 39.3 & 0.811 \\
 & $L_{8}\,\SV{3}$ & 0.385 & $-$0.001 & 0.052 & 0.0138 & 24.1 & 0.863 & 0.059 & 0.0086 & 38.0 & 0.844 \\
 & $L_{7}\,\SV{2}$ & 0.376 & $-$0.001 & 0.048 & 0.0136 & 24.2 & 0.870 & 0.051 & 0.0087 & 37.5 & 0.861 \\
 & 120 cells & -- & -- & max 0.056 &  & 120/120 & 0.972 (0.860) & max 0.075 &  & 120/120 & 0.963 (0.801) \\
\end{tabular*}}\\
\midrule
\textbf{C}\enspace Residualisation: residual $r$ / raw $r$ after regressing each coordinate on gene annotations (MS null, $N = 2{,}000$)\\*
{\scriptsize\setlength{\tabcolsep}{2.5pt}\begin{tabular*}{\textwidth}{@{\extracolsep{\fill}}>{\raggedright\arraybackslash}p{3.5cm}crrcrrcrr@{}}
& \multicolumn{3}{c}{lung, $L_3\,\SV{1}$} & \multicolumn{3}{c}{lung (ext.), $L_8\,\SV{2}$} & \multicolumn{3}{c}{kidney~B, $L_9\,\SV{3}$} \\
\cmidrule(lr){2-4}\cmidrule(lr){5-7}\cmidrule(l){8-10}
Design (regressors) & Strongest cell & Median & Sig. & Strongest cell & Median & Sig. & Strongest cell & Median & Sig. \\
\midrule
Class factor (8 levels) $+ \log n_{\mathrm{obs}}$ & 0.591 (0.535--0.649) & 0.873 & 119 & 0.602 (0.527--0.677) & 0.848 & 120 & 0.629 (0.569--0.688) & 0.848 & 120 \\
30 GO-CC terms $+ \log n_{\mathrm{obs}}$ (pre-specified) & 0.404 (0.327--0.477) & 0.702 & 116 & 0.412 (0.337--0.498) & 0.684 & 120 & 0.471 (0.395--0.552) & 0.723 & 120 \\
\multicolumn{10}{@{}l}{\itshape Post hoc controls (no test)}\\
$\log n_{\mathrm{obs}}$ alone & 1.006 & 0.999 & -- & 0.980 & 0.994 & -- & 0.967 & 0.980 & -- \\
30 GO-CC terms, without $\log n_{\mathrm{obs}}$ & 0.390 & 0.714 & -- & 0.405 & 0.706 & -- & 0.471 & 0.750 & -- \\
8 class-defining terms, multi-label $+ \log n_{\mathrm{obs}}$ & 0.587 & 0.849 & -- & 0.612 & 0.834 & -- & 0.658 & 0.820 & -- \\
30 GO-BP terms $+ \log n_{\mathrm{obs}}$ & 0.583 & 0.704 & -- & 0.539 & 0.713 & -- & 0.609 & 0.688 & -- \\
30 GO-MF terms $+ \log n_{\mathrm{obs}}$ & 0.711 & 0.785 & -- & 0.707 & 0.784 & -- & 0.708 & 0.787 & -- \\
30 random gene sets $+ \log n_{\mathrm{obs}}$ (5 seeds) & 0.99--1.01 & 0.99--1.00 & -- & 0.97--0.98 & 0.98--0.99 & -- & 0.96--0.97 & 0.97--0.98 & -- \\
\end{tabular*}}\\
\midrule
\textbf{D}\enspace Edges within one class and between classes, each against its own MS null ($N = 1{,}000$)\\*
{\scriptsize\begin{tabular*}{\textwidth}{@{\extracolsep{\fill}}lllrccrrrr@{}}
& & & & \multicolumn{2}{c}{Strongest cell, $r$ (95\% CI)} & \multicolumn{2}{c}{Median $r$ over 120 cells} & \multicolumn{2}{c}{Sig.} \\
\cmidrule(lr){5-6}\cmidrule(lr){7-8}\cmidrule(l){9-10}
Extraction & Cell & Edges & $m$ & Raw & Class-centred & Raw & Class-centred & Raw & Class-centred \\
\midrule
lung & $L_{3}\,\SV{1}$ & Same class & 6,582 & 0.497 (0.432--0.567) & 0.309 (0.241--0.372) & 0.197 & 0.175 & 117 & 115 \\
 &  & Different classes & 7,892 & 0.377 (0.315--0.438) & 0.204 (0.159--0.258) & 0.155 & 0.145 & 120 & 117 \\
lung (ext.) & $L_{8}\,\SV{2}$ & Same class & 6,380 & 0.477 (0.404--0.548) & 0.295 (0.227--0.363) & 0.186 & 0.166 & 120 & 120 \\
 &  & Different classes & 7,752 & 0.340 (0.277--0.397) & 0.198 (0.145--0.250) & 0.145 & 0.134 & 119 & 119 \\
kidney~B & $L_{9}\,\SV{3}$ & Same class & 5,565 & 0.481 (0.416--0.549) & 0.312 (0.248--0.381) & 0.225 & 0.199 & 120 & 120 \\
 &  & Different classes & 6,869 & 0.331 (0.262--0.397) & 0.204 (0.147--0.256) & 0.181 & 0.166 & 116 & 116 \\
\end{tabular*}}\\
\midrule
\textbf{E}\enspace Gene set restricted to genes observed in at least 3, 5 or 10 sampled cells (SVD recomputed; MS null, $N = 2{,}000$)\\*
{\scriptsize\begin{tabular*}{\textwidth}{@{\extracolsep{\fill}}llrrrlcrlrcr@{}}
& & & & & \multicolumn{3}{c}{Strongest cell} & \multicolumn{4}{c}{Strongest direction of the full-$V$ analysis} \\
\cmidrule(lr){6-8}\cmidrule(l){9-12}
Extraction & $n_{\mathrm{obs}}$ & $n_V$ & $m$ & Sig. & Cell & $r$ (95\% CI) & Median $r$ & Cell & $|$corr$|$ & $r$ (95\% CI) & $z$ \\
\midrule
lung & all ($V$) & 5,648 & 14,474 & 120 & $L_{3}\,\SV{1}$ & 0.435 (0.384--0.486) & 0.181 & $L_{3}\,\SV{1}$ & 1 & 0.435 (0.384--0.486) & 53.1 \\
 & $\geq 3$ & 4,221 & 10,685 & 120 & $L_{4}\,\SV{1}$ & 0.418 (0.365--0.470) & 0.189 & $L_{3}\,\SV{1}$ & 0.999 & 0.406 (0.353--0.463) & 41.8 \\
 & $\geq 5$ & 3,558 & 8,244 & 120 & $L_{4}\,\SV{1}$ & 0.397 (0.339--0.457) & 0.191 & $L_{3}\,\SV{1}$ & 0.996 & 0.393 (0.336--0.452) & 36.5 \\
 & $\geq 10$ & 2,580 & 5,228 & 119 & $L_{5}\,\SV{3}$ & 0.401 (0.330--0.477) & 0.198 & $L_{3}\,\SV{1}$ & 0.987 & 0.366 (0.295--0.436) & 27.8 \\
\addlinespace[2pt]
lung (ext.) & all ($V$) & 5,616 & 14,132 & 120 & $L_{8}\,\SV{2}$ & 0.406 (0.351--0.459) & 0.170 & $L_{8}\,\SV{2}$ & 1 & 0.406 (0.351--0.459) & 49.2 \\
 & $\geq 3$ & 4,529 & 11,604 & 120 & $L_{3}\,\SV{1}$ & 0.401 (0.344--0.458) & 0.163 & $L_{8}\,\SV{2}$ & 0.986 & 0.389 (0.336--0.443) & 43.0 \\
 & $\geq 5$ & 3,976 & 9,981 & 118 & $L_{3}\,\SV{1}$ & 0.396 (0.338--0.455) & 0.160 & $L_{8}\,\SV{2}$ & 0.975 & 0.381 (0.323--0.441) & 38.4 \\
 & $\geq 10$ & 3,060 & 7,341 & 117 & $L_{10}\,\SV{2}$ & 0.398 (0.336--0.463) & 0.174 & $L_{8}\,\SV{2}$ & 0.962 & 0.379 (0.305--0.442) & 32.7 \\
\addlinespace[2pt]
kidney~B & all ($V$) & 4,780 & 12,434 & 120 & $L_{9}\,\SV{3}$ & 0.404 (0.351--0.457) & 0.209 & $L_{9}\,\SV{3}$ & 1 & 0.404 (0.351--0.457) & 45.9 \\
 & $\geq 3$ & 3,528 & 8,812 & 120 & $L_{5}\,\SV{4}$ & 0.369 (0.294--0.445) & 0.186 & $L_{9}\,\SV{3}$ & 0.990 & 0.350 (0.289--0.406) & 33.8 \\
 & $\geq 5$ & 2,914 & 6,676 & 120 & $L_{5}\,\SV{4}$ & 0.397 (0.319--0.488) & 0.198 & $L_{9}\,\SV{3}$ & 0.976 & 0.334 (0.263--0.402) & 27.2 \\
 & $\geq 10$ & 2,023 & 4,134 & 120 & $L_{5}\,\SV{3}$ & 0.452 (0.373--0.532) & 0.216 & $L_{9}\,\SV{3}$ & 0.787 & 0.304 (0.229--0.377) & 20.0 \\
\end{tabular*}}\\
\end{longtable}
\endgroup

\setcounter{table}{5}
\begingroup
\footnotesize
\setlength{\tabcolsep}{3pt}
\renewcommand{\arraystretch}{1.05}
\begin{longtable}{@{}>{\scriptsize\raggedright\arraybackslash}p{3.1cm}*{6}{>{\scriptsize\raggedright\arraybackslash}p{2.1cm}}@{}}
\caption{\textbf{Transcription-factor identity (H2): primary statistic, baselines, adjustment, spectral bands and the reported 0.744.} Score: leave-one-out (LOO) mean cosine of a gene's $V$-centred vector to the census TFs of Lambert et al.\ (excluding itself); statistic: AUROC of TFs against all other observed genes $V$, stratified by deciles of $\log n_{\mathrm{obs}}$ ($n_{\mathrm{obs}}$, cells whose input contained the gene) unless marked unstratified. $|V|$, observed genes; $L_\ell$, layer; $z$, standardised distance from the within-stratum permutation null; max-$T$, Westfall--Young correction over the 12 layers (primary, $N = 10{,}000$ joint permutations, floor $1/10{,}001 = 1.0\times10^{-4}$) or over the 936 band cells (Panels D, E; $N = 2{,}000$, floor $5.0\times10^{-4}$). Intervals are 95\% gene-bootstrap percentile intervals ($B = 1{,}000$; LOO scores held fixed), written [lower, upper]; ranges over layers or extractions use an en-dash (or `to' when a bound is negative). Because the TF centroid is not re-estimated in the bootstrap, a half-sampling SE that re-estimates it is 0.83--0.84 times the bootstrap SE in lung and lung (ext.) but 1.15--1.32 times it in the two immune extractions checked (immune, CD4 T cells), so intervals in the immune extractions are about 15--30\% too narrow. \textbf{A}: all 16 extractions (the 11 robustness extractions below the rule; the three checkpoints are applied to the same 192 kidney cells); the decision uses lung and lung (ext.). \textbf{B}: the same statistic for single-gene scores (higher value scores higher) and for static embeddings scored by the same LOO cosine; unstratified AUROC in parentheses for detection; ESM2-15B is the 5{,}120-dimensional per-protein embedding (UCE file), matched for 94--96\% of $V$ in the lung extractions; paired differences (static minus residual-stream layer) use identical genes, labels and bootstrap draws, with the range over the 12 layers and the lowest lower bound; `all 16' gives the range over extractions (and, for the paired rows, how many of the 192 extraction--layer intervals include 0). \textbf{C}: standardised coefficient of the score, its smallest lower 95\% Wald bound and the largest likelihood-ratio $p$ over the 12 layers; restricted-$V$ rows re-centre and re-stratify on the retained genes ($N = 2{,}000$; genes / TFs below the range). \textbf{D}: each direction scored by its raw coordinate with a sign-free AUROC, at the layer where it is highest; $p$ is the 936-cell max-$T$ $p$ of that cell; $^{a}$primary analysis (Panel A): $z$ and 12-layer max-$T$ $p$ from $N = 10{,}000$ permutations. \textbf{E}: cos variant (LOO cosine within the band) unless stated; the centroid (LDA) variant is fitted in-sample and refitted in each null draw, so its maximum is optimistic. Screen bands are those of the automated screen, indexed on the observed-gene basis; in parentheses, the number of layers with $p<0.05$ on the band's own 12-layer max-$T$ and the smallest $p$ in the 936-cell max-$T$; $^{b}$kidney~A: the band read in the reference decomposition of the full kidney~A matrix, zero rows included, and mapped to the observed-gene basis, in which its $\SV{k}$ is observed $\SV{k-1}$ (Table~S14). The screen assembled $\SV{2}$--$\SV{4}$ on a 209-gene decomposition whose band also maps to observed $\SV{1}$--$\SV{3}$, and its later kidney~A iterations on mean embeddings used $\SV{5}$--$\SV{7}$ of the observed genes, which maps to itself: AUROC 0.429--0.696 over the 12 layers (10 layers with $p<0.05$; smallest $p$ of 936, 0.003). Cross-fit: the best cell chosen on a random half of the genes and scored on the other half, 200 halvings; the brackets are the 2.5--97.5\% quantiles over halvings of one data set, not sampling intervals. Immune and kidney~A are the extractions in which the screen's results were found. \textbf{F}: 361 genes named by the screen, 295 observed; per-layer values in the CSV. TRRUST task: TRRUST TFs against genes that are only ever TRRUST targets, same machinery as Panel A. Per-layer values of every analysis and all 936 band cells of the five main extractions are in the CSV. Full values: \texttt{Table\_S6\_tf\_identity.csv}.}\label{tab:S6}\\
\toprule
\endfirsthead
\multicolumn{7}{@{}l}{\footnotesize Table~\thetable\ (continued)}\\
\toprule
\endhead
\midrule
\multicolumn{7}{r@{}}{\footnotesize continued on next page}\\
\endfoot
\bottomrule
\endlastfoot
\multicolumn{7}{@{}p{16.4cm}@{}}{\footnotesize\textbf{A. Pre-registered H2 statistic (full 512-dimensional representation), per extraction}}\\
Extraction & $|V|$ / TFs & AUROC, 12 layers & Best layer: AUROC [95\% CI] & $z$, 12 layers & Smallest max-$T$ $p$ & Layers max-$T$ $p<0.05$\\
\cmidrule(lr){1-7}
lung & 5{,}648 / 583 & 0.625--0.665 & $L_{5}$: 0.665\newline [0.641, 0.688] & 8.8--11.5 & $1.0\times10^{-4}$ & 12/12\\
lung (ext.) & 5{,}616 / 623 & 0.627--0.662 & $L_{1}$: 0.662\newline [0.637, 0.686] & 8.9--11.6 & $1.0\times10^{-4}$ & 12/12\\
kidney~B & 4{,}780 / 511 & 0.615--0.654 & $L_{1}$: 0.654\newline [0.628, 0.681] & 7.3--10.6 & $1.0\times10^{-4}$ & 12/12\\
immune & 2{,}039 / 113 & 0.566--0.720 & $L_{2}$: 0.720\newline [0.657, 0.772] & 2.3--6.5 & $1.0\times10^{-4}$ & 12/12\\
kidney~A & 1{,}972 / 90 & 0.678--0.783 & $L_{3}$: 0.783\newline [0.732, 0.829] & 5.4--7.6 & $1.0\times10^{-4}$ & 12/12\\
\cmidrule(lr){1-7}
immune (seed 43) & 2{,}080 / 105 & 0.549--0.657 & $L_{2}$: 0.657\newline [0.596, 0.719] & 1.8--4.8 & $1.0\times10^{-4}$ & 10/12\\
immune (seed 44) & 1{,}972 / 105 & 0.604--0.720 & $L_{2}$: 0.720\newline [0.667, 0.771] & 3.4--6.6 & $1.0\times10^{-4}$ & 12/12\\
immune (donor TSP14) & 1{,}964 / 103 & 0.584--0.704 & $L_{2}$: 0.704\newline [0.640, 0.759] & 2.8--6.0 & $1.0\times10^{-4}$ & 12/12\\
immune (donor TSP21) & 1{,}765 / 95 & 0.575--0.693 & $L_{2}$: 0.693\newline [0.631, 0.748] & 2.4--5.5 & $1.0\times10^{-4}$ & 12/12\\
immune (donor TSP25) & 1{,}791 / 98 & 0.547--0.649 & $L_{0}$: 0.649\newline [0.584, 0.715] & 1.7--4.5 & $1.0\times10^{-4}$ & 9/12\\
immune (B cells) & 1{,}398 / 85 & 0.615--0.694 & $L_{2}$: 0.694\newline [0.613, 0.772] & 3.3--5.0 & $1.0\times10^{-4}$ & 12/12\\
immune (CD4 T cells) & 1{,}465 / 78 & 0.576--0.663 & $L_{2}$: 0.663\newline [0.586, 0.739] & 2.1--4.0 & $2.0\times10^{-4}$ & 12/12\\
immune (CD8 T cells) & 1{,}595 / 96 & 0.599--0.707 & $L_{0}$: 0.707\newline [0.640, 0.768] & 3.0--5.5 & $1.0\times10^{-4}$ & 12/12\\
kidney (whole-human ckpt) & 4{,}923 / 545 & 0.612--0.652 & $L_{1}$: 0.652\newline [0.626, 0.676] & 7.3--10.5 & $1.0\times10^{-4}$ & 12/12\\
kidney (kidney ckpt) & 4{,}923 / 545 & 0.565--0.595 & $L_{3}$: 0.595\newline [0.567, 0.622] & 4.7--6.9 & $1.0\times10^{-4}$ & 12/12\\
kidney (brain ckpt) & 4{,}923 / 545 & 0.562--0.607 & $L_{0}$: 0.607\newline [0.579, 0.633] & 4.2--6.7 & $1.0\times10^{-4}$ & 12/12\\
\midrule
\multicolumn{7}{@{}p{16.4cm}@{}}{\footnotesize\textbf{B. Baselines: detection-stratified AUROC of single-gene scores and static embeddings on the same genes and strata}}\\
 & lung & lung (ext.) & kidney~B & immune & kidney~A & all 16\\
\cmidrule(lr){1-7}
Detection alone, $\log n_{\mathrm{obs}}$ (unstratified AUROC) & 0.503 (0.458)\newline [0.484, 0.522] & 0.531 (0.454)\newline [0.509, 0.552] & 0.510 (0.461)\newline [0.490, 0.529] & 0.522 (0.650)\newline [0.485, 0.560] & 0.540 (0.555)\newline [0.490, 0.590] & 0.466--0.540\\
Mean expression, extraction cells & 0.588\newline [0.565, 0.611] & 0.595\newline [0.571, 0.620] & 0.532\newline [0.504, 0.560] & 0.486\newline [0.424, 0.544] & 0.588\newline [0.518, 0.658] & 0.435--0.595\\
Mean expression, population & 0.579\newline [0.556, 0.602] & 0.576\newline [0.552, 0.601] & 0.524\newline [0.496, 0.550] & 0.512\newline [0.459, 0.569] & 0.606\newline [0.546, 0.665] & 0.506--0.606\\
scGPT input gene embedding, $V$-centred & 0.760\newline [0.740, 0.781] & 0.744\newline [0.722, 0.765] & 0.742\newline [0.718, 0.765] & 0.750\newline [0.684, 0.806] & 0.770\newline [0.724, 0.812] & 0.678--0.770\\
scGPT input gene embedding, uncentred & 0.668 & 0.638 & 0.649 & 0.738 & 0.747 & 0.638--0.762\\
ESM2-15B per-protein embedding, $V$-centred & 0.881\newline [0.865, 0.897] & 0.883\newline [0.868, 0.898] & 0.862\newline [0.844, 0.880] & 0.938\newline [0.909, 0.963] & 0.917\newline [0.877, 0.948] & 0.862--0.943\\
ESM2-15B per-protein embedding, uncentred & 0.761 & 0.761 & 0.725 & 0.865 & 0.844 & 0.725--0.866\\
Genes with an ESM2-15B vector & 5{,}413 & 5{,}306 & 4{,}602 & 1{,}554 & 1{,}399 & 1{,}075--5{,}413\\
Input embedding minus layer (paired; 12 layers) & 0.095--0.135\newline lower $\geq$ 0.074 & 0.082--0.118\newline lower $\geq$ 0.063 & 0.089--0.127\newline lower $\geq$ 0.068 & 0.030--0.184\newline lower $\geq$ $-$0.005 & $-$0.013 to $+$0.092\newline lower $\geq$ $-$0.044 & $-$0.013 to $+$0.184\newline 19/192 CIs include 0\\
ESM2-15B minus layer (paired; 12 layers) & 0.216--0.263\newline lower $\geq$ 0.188 & 0.209--0.245\newline lower $\geq$ 0.183 & 0.208--0.248\newline lower $\geq$ 0.175 & 0.238--0.400\newline lower $\geq$ 0.188 & 0.195--0.296\newline lower $\geq$ 0.139 & 0.195--0.402\newline 0/192 CIs include 0\\
\midrule
\multicolumn{7}{@{}p{16.4cm}@{}}{\footnotesize\textbf{C. Adjustment: logistic regression TF $\sim$ score + $\log n_{\mathrm{obs}}$ + $\log$ mean expression (standardised), per layer; and $V$ restricted by detection}}\\
 & lung & lung (ext.) & kidney~B & immune & kidney~A & all 16\\
\cmidrule(lr){1-7}
$\beta$ per SD of score, extraction-cell expression & 0.367--0.560 & 0.454--0.565 & 0.462--0.598 & 0.464--1.158 & 1.036--2.029 & 0.177--2.029\\
Smallest lower 95\% bound of $\beta$ & 0.265 & 0.365 & 0.363 & 0.172 & 0.762 & 0.086\\
Largest likelihood-ratio $p$, extraction-cell expression & $6.1\times10^{-13}$ & $1.8\times10^{-23}$ & $8.0\times10^{-20}$ & 0.002 & $1.5\times10^{-16}$ & 0.008\\
$\beta$ per SD of score, population expression & 0.365--0.549 & 0.477--0.579 & 0.456--0.589 & 0.484--1.168 & 1.000--2.083 & 0.176--2.083\\
Largest likelihood-ratio $p$, population expression & $2.5\times10^{-13}$ & $1.2\times10^{-26}$ & $3.6\times10^{-20}$ & 0.001 & $1.4\times10^{-13}$ & 0.008\\
H2 AUROC, genes with $n_{\mathrm{obs}} \geq 3$ (12 layers) & 0.646--0.658\newline 4{,}221 / 407 & 0.627--0.662\newline 4{,}529 / 502 & 0.611--0.649\newline 3{,}528 / 367 & 0.625--0.723\newline 1{,}341 / 90 & 0.660--0.764\newline 1{,}346 / 79 & 0.517--0.764\\
H2 AUROC, genes with $n_{\mathrm{obs}} \geq 5$ (12 layers) & 0.649--0.670\newline 3{,}558 / 329 & 0.635--0.663\newline 3{,}976 / 420 & 0.607--0.654\newline 2{,}914 / 296 & 0.632--0.716\newline 1{,}149 / 85 & 0.613--0.729\newline 1{,}208 / 74 & 0.508--0.730\\
H2 AUROC, genes with $n_{\mathrm{obs}} \geq 10$ (12 layers) & 0.668--0.694\newline 2{,}580 / 223 & 0.666--0.682\newline 3{,}060 / 297 & 0.598--0.661\newline 2{,}023 / 200 & 0.660--0.730\newline 965 / 77 & 0.645--0.736\newline 1{,}041 / 64 & 0.510--0.736\\
\midrule
\multicolumn{7}{@{}p{16.4cm}@{}}{\footnotesize\textbf{D. Single directions $\SV{1}$--$\SV{12}$ (raw coordinate, sign-free AUROC), best of the 12 layers}}\\
Direction & lung: AUROC (layer) & lung: $z$ & lung: $p$ (936 cells) & lung (ext.): AUROC (layer) & lung (ext.): $z$ & lung (ext.): $p$ (936 cells)\\
\cmidrule(lr){1-7}
$\SV{1}$ & 0.572 ($L_{3}$) & 8.0 & $5.0\times10^{-4}$ & 0.561 ($L_{4}$) & 6.4 & $5.0\times10^{-4}$\\
$\SV{2}$ & 0.567 ($L_{5}$) & 7.5 & $5.0\times10^{-4}$ & 0.558 ($L_{1}$) & 6.1 & $5.0\times10^{-4}$\\
$\SV{3}$ & 0.559 ($L_{8}$) & 6.4 & $1.0\times10^{-3}$ & 0.570 ($L_{1}$) & 7.8 & $5.0\times10^{-4}$\\
$\SV{4}$ & 0.564 ($L_{0}$) & 7.1 & $5.0\times10^{-4}$ & 0.572 ($L_{0}$) & 8.1 & $5.0\times10^{-4}$\\
$\SV{5}$ & 0.586 ($L_{10}$) & 9.9 & $5.0\times10^{-4}$ & 0.582 ($L_{4}$) & 9.6 & $5.0\times10^{-4}$\\
$\SV{6}$ & 0.598 ($L_{4}$) & 11.2 & $5.0\times10^{-4}$ & 0.654 ($L_{3}$) & 18.7 & $5.0\times10^{-4}$\\
$\SV{7}$ & 0.559 ($L_{8}$) & 6.4 & $1.0\times10^{-3}$ & 0.637 ($L_{8}$) & 15.9 & $5.0\times10^{-4}$\\
$\SV{8}$ & 0.631 ($L_{0}$) & 15.2 & $5.0\times10^{-4}$ & 0.624 ($L_{2}$) & 14.7 & $5.0\times10^{-4}$\\
$\SV{9}$ & 0.616 ($L_{4}$) & 13.7 & $5.0\times10^{-4}$ & 0.593 ($L_{4}$) & 10.7 & $5.0\times10^{-4}$\\
$\SV{10}$ & 0.601 ($L_{8}$) & 11.4 & $5.0\times10^{-4}$ & 0.600 ($L_{10}$) & 11.8 & $5.0\times10^{-4}$\\
$\SV{11}$ & 0.569 ($L_{10}$) & 7.0 & $5.0\times10^{-4}$ & 0.631 ($L_{11}$) & 15.4 & $5.0\times10^{-4}$\\
$\SV{12}$ & 0.560 ($L_{1}$) & 6.3 & $1.0\times10^{-3}$ & 0.617 ($L_{7}$) & 13.5 & $5.0\times10^{-4}$\\
Full 512, best layer$^{a}$ & 0.665 ($L_{5}$) & 11.3 & $1.0\times10^{-4}$ & 0.662 ($L_{1}$) & 11.6 & $1.0\times10^{-4}$\\
\midrule
\multicolumn{7}{@{}p{16.4cm}@{}}{\footnotesize\textbf{E. All 78 contiguous bands of $\SV{1}$--$\SV{12}$ at 12 layers (936 cells), screen bands, and cross-fitted band choice}}\\
 & lung & lung (ext.) & kidney~B & immune & kidney~A & all 16\\
\cmidrule(lr){1-7}
Cells with max-$T$ $p<0.05$ (of 936) & 633 & 690 & 591 & 71 & 359 & 4--690\\
Best cell by AUROC & $\SV{3}$--$\SV{12}$, $L_{2}$\newline 0.676 & $\SV{3}$--$\SV{12}$, $L_{2}$\newline 0.711 & $\SV{7}$--$\SV{11}$, $L_{2}$\newline 0.662 & $\SV{3}$--$\SV{12}$, $L_{0}$\newline 0.730 & $\SV{1}$--$\SV{9}$, $L_{2}$\newline 0.738 & 0.592--0.738\\
Cell with the largest $z$: AUROC ($z$) & $\SV{8}$, $L_{0}$\newline 0.631 (15.2) & $\SV{6}$, $L_{3}$\newline 0.654 (18.7) & $\SV{7}$, $L_{2}$\newline 0.646 (16.9) & $\SV{11}$, $L_{9}$\newline 0.674 (8.1) & $\SV{5}$, $L_{3}$\newline 0.699 (8.4) & --\\
Null 97.5\% quantile of the maximum AUROC & 0.548 & 0.547 & 0.551 & 0.613 & 0.623 & 0.547--0.641\\
Centroid (LDA) variant, best cell (in-sample) & $\SV{1}$--$\SV{12}$, $L_{2}$\newline 0.700 & $\SV{1}$--$\SV{12}$, $L_{2}$\newline 0.734 & $\SV{1}$--$\SV{12}$, $L_{2}$\newline 0.693 & $\SV{1}$--$\SV{12}$, $L_{0}$\newline 0.763 & $\SV{1}$--$\SV{12}$, $L_{0}$\newline 0.805 & 0.619--0.805\\
Screen band $\SV{2}$--$\SV{4}$: AUROC, 12 layers (layers $p<0.05$; smallest $p$ of 936) & 0.522--0.590\newline (11; 0.003) & 0.530--0.588\newline (12; $5.0\times10^{-4}$) & 0.519--0.587\newline (10; 0.002) & 0.503--0.593\newline (3; 0.40) & $\SV{1}$--$\SV{3}$$^{b}$: 0.595--0.638\newline (12; 0.040) & --\\
Screen band $\SV{5}$--$\SV{7}$: AUROC, 12 layers (layers $p<0.05$; smallest $p$ of 936) & 0.497--0.597\newline (11; $5.0\times10^{-4}$) & 0.550--0.654\newline (12; $5.0\times10^{-4}$) & 0.546--0.641\newline (12; $5.0\times10^{-4}$) & 0.256--0.612\newline (4; 0.35) & $\SV{4}$--$\SV{6}$$^{b}$: 0.566--0.669\newline (11; 0.030) & --\\
Screen band $\SV{2}$--$\SV{7}$: AUROC, 12 layers (layers $p<0.05$; smallest $p$ of 936) & 0.574--0.611\newline (12; $5.0\times10^{-4}$) & 0.580--0.653\newline (12; $5.0\times10^{-4}$) & 0.555--0.645\newline (12; $5.0\times10^{-4}$) & 0.515--0.622\newline (4; 0.20) & $\SV{1}$--$\SV{6}$$^{b}$: 0.628--0.715\newline (12; $5.0\times10^{-4}$) & --\\
Cross-fitted band, held-out AUROC (median [2.5--97.5\% of 200 halvings]) & 0.669\newline [0.644, 0.692] & 0.702\newline [0.678, 0.725] & 0.648\newline [0.618, 0.676] & 0.681\newline [0.578, 0.745] & 0.695\newline [0.572, 0.760] & 0.571--0.702\\
Full 512, held-out AUROC (median [2.5--97.5\% of 200 halvings]) & 0.660\newline [0.643, 0.685] & 0.659\newline [0.636, 0.680] & 0.651\newline [0.630, 0.674] & 0.687\newline [0.600, 0.773] & 0.761\newline [0.699, 0.814] & 0.589--0.761\\
Band minus full 512, held-out (median [2.5--97.5\% of 200 halvings]) & $+$0.007\newline [$-$0.014, 0.029] & $+$0.042\newline [0.019, 0.065] & $-$0.004\newline [$-$0.032, 0.024] & $-$0.002\newline [$-$0.143, 0.058] & $-$0.058\newline [$-$0.232, $-$0.012] & $-$0.058 to $+$0.042\\
Optimism (in-sample minus held-out), median & 0.007 & 0.005 & 0.019 & 0.040 & 0.055 & 0.005--0.101\\
Most frequent choice (of 200) & $\SV{4}$--$\SV{12}$ at $L_{2}$ (57) & $\SV{3}$--$\SV{12}$ at $L_{2}$ (133) & $\SV{7}$--$\SV{11}$ at $L_{2}$ (64) & $\SV{3}$--$\SV{12}$ at $L_{0}$ (77) & $\SV{2}$--$\SV{5}$ at $L_{3}$ (35) & --\\
\midrule
\multicolumn{7}{@{}p{16.4cm}@{}}{\footnotesize\textbf{F. The reported 0.744 (immune; screen labels: 73 TFs against 222 target-only genes) and the TRRUST task}}\\
Step & Layer-mean AUROC & \multicolumn{5}{>{\scriptsize}l}{Procedure and remarks}\\
\cmidrule(lr){1-7}
Value reported by the screen & 0.744 & \multicolumn{5}{>{\scriptsize\raggedright\arraybackslash}p{11.2cm}@{}}{As recorded in the screen's ledger.}\\
Screen procedure, reproduced & 0.747 & \multicolumn{5}{>{\scriptsize\raggedright\arraybackslash}p{11.2cm}@{}}{Logistic regression ($C = 1$) on standardised $\SV{2}$--$\SV{7}$ coordinates, 5-fold stratified CV (seed 42), fold-mean AUROC averaged over 12 layers; over CV seeds 0--49 median 0.744, 2.5--97.5\% 0.727--0.759; permutation $p = 1/1{,}001$ ($N = 1{,}000$, the floor); best layer $L_{2}$ 0.792.}\\
Same out-of-fold predictions, pooled & 0.745 & \multicolumn{5}{>{\scriptsize\raggedright\arraybackslash}p{11.2cm}@{}}{Unstratified AUROC of the pooled out-of-fold predictions.}\\
Same out-of-fold predictions, stratified & 0.683 & \multicolumn{5}{>{\scriptsize\raggedright\arraybackslash}p{11.2cm}@{}}{Detection-stratified AUROC (deciles of $\log n_{\mathrm{obs}}$).}\\
LOO cosine in $\SV{2}$--$\SV{7}$, unstratified & 0.644 & \multicolumn{5}{>{\scriptsize\raggedright\arraybackslash}p{11.2cm}@{}}{Unsupervised score; layer-mean permutation $p$ $1.0\times10^{-4}$ ($N = 10{,}000$).}\\
LOO cosine in $\SV{2}$--$\SV{7}$, stratified & 0.553 & \multicolumn{5}{>{\scriptsize\raggedright\arraybackslash}p{11.2cm}@{}}{Layer-mean permutation $p$ 0.042; per layer highest at $L_{2}$ (0.666), lowest at $L_{9}$ (0.398).}\\
LOO cosine, full 512, unstratified & 0.702 & \multicolumn{5}{>{\scriptsize\raggedright\arraybackslash}p{11.2cm}@{}}{Layer-mean permutation $p$ $1.0\times10^{-4}$.}\\
LOO cosine, full 512, stratified & 0.629 & \multicolumn{5}{>{\scriptsize\raggedright\arraybackslash}p{11.2cm}@{}}{Layer-mean permutation $p$ $5.0\times10^{-4}$.}\\
Attribution of the 0.747 & -- & \multicolumn{5}{>{\scriptsize\raggedright\arraybackslash}p{11.2cm}@{}}{Order-dependent: removing supervision first costs 0.10 and then stratifying 0.09 (0.747 $\to$ 0.644 $\to$ 0.553); stratifying first costs 0.06 and then removing supervision 0.13 (0.747 / 0.745 $\to$ 0.683 $\to$ 0.553). Supervision accounts for 0.10--0.13, detection for 0.06--0.09.}\\
\cmidrule(lr){1-7}
TRRUST task: TRRUST TFs against target-only genes & lung & lung (ext.) & kidney~B & immune & kidney~A & all 16\\
\cmidrule(lr){1-7}
Genes / TFs & 1{,}956 / 555 & 1{,}971 / 579 & 1{,}577 / 480 & 283 / 68 & 180 / 55 & 173--1{,}971 / 50--579\\
Full 512, AUROC over 12 layers (layers max-$T$ $p<0.05$) & 0.646--0.713\newline (12/12) & 0.643--0.710\newline (12/12) & 0.618--0.692\newline (12/12) & 0.519--0.671\newline (8/12) & 0.602--0.718\newline (12/12) & 0.487--0.718\\
Mean expression alone & 0.628 & 0.609 & 0.567 & 0.458 & 0.608 & 0.445--0.628\\
Detection alone & 0.502 & 0.520 & 0.510 & 0.551 & 0.615 & 0.444--0.615\\
Cross-fitted band minus full 512 (median [2.5--97.5\%]) & $-$0.047\newline [$-$0.095, $-$0.032] & $-$0.029\newline [$-$0.073, $-$0.011] & $-$0.061\newline [$-$0.126, $-$0.034] & $+$0.039\newline [$-$0.139, 0.185] & $-$0.036\newline [$-$0.265, 0.118] & $-$0.071 to $+$0.056\\
\end{longtable}
\endgroup

\setcounter{table}{6}
\begingroup
% The caption is longer than one line of 16,383pt, the limit of longtable's centring test; within this group the
% caption is always set as a paragraph (the branch longtable itself takes for any caption wider than the line).
\makeatletter
\def\LT@makecaption#1#2#3{\LT@mcol\LT@cols c{\hbox to\z@{\hss\parbox[t]\LTcapwidth{\reset@font
  #1{#2: }#3\endgraf\vskip\baselineskip}\hss}}}
\makeatother
\footnotesize
\setlength{\tabcolsep}{1.5pt}
\renewcommand{\arraystretch}{1.05}
\begin{longtable}{@{}>{\scriptsize\raggedright\arraybackslash}p{1.9cm}>{\scriptsize\raggedright\arraybackslash}p{1.5cm}>{\scriptsize\raggedright\arraybackslash}p{1.572cm}>{\scriptsize\raggedright\arraybackslash}p{1.572cm}>{\scriptsize\raggedright\arraybackslash}p{1.572cm}>{\scriptsize\raggedright\arraybackslash}p{1.572cm}>{\scriptsize\raggedright\arraybackslash}p{1.572cm}>{\scriptsize\raggedright\arraybackslash}p{1.572cm}>{\scriptsize\raggedright\arraybackslash}p{1.572cm}>{\scriptsize\raggedright\arraybackslash}p{1.572cm}@{}}
\caption{\textbf{Regulatory edges (H3): the benchmark in every extraction and every edge-level sensitivity analysis.} Positives: TRRUST pairs with both genes observed; candidates: up to ten per positive, pairs of the same factor whose target is a documented target of another factor and that are absent, in either direction, from TRRUST and from DoRothEA at any confidence level, and matched where possible to the positive on the target's decile of TRRUST in-degree and of $n_{\mathrm{obs}}$ (on both for 75--90\% of candidates, otherwise on in-degree or not at all; CSV column \texttt{frac\_cand\_fully\_matched}); factors with $\geq 3$ of each. Statistic: within-factor stratified AUROC of the signed cosine of the $V$-centred 512-dimensional vectors at layer $L_\ell$. CI, 95\% interval; UB, one-sided 95\% upper bound; both from the two-way (factor $\times$ target) pigeonhole bootstrap, $B = 2{,}000$ unless stated (1{,}000 for some strata, CSV column \texttt{B}); Max.\ UB, the largest UB over the twelve layers. $|r|$, absolute Pearson co-expression on the same pairs, on the extraction's own cells (extr.; 256, or 192 for kidney~B and the checkpoints) or on 3{,}000 cells of the same population (3k); $\Delta$, cosine minus $|r|$ AUROC on the same bootstrap draws. \textbf{A}: max-$T$ $p$, within-factor label permutation applied jointly to the twelve layers ($N = 10{,}000$; floor $1.0\times10^{-4}$); MDE, the AUROC (and $\Delta$) detected with 80\% power at two-sided $\alpha = 0.05$ in a simulation that resamples factors as clusters with the observed between-factor heterogeneity and tests with the two-way bootstrap. Verdict: D, detected (max-$T$ $p < 0.05$); U, every layer's UB $< 0.60$ (AUROC $\geq 0.60$ excluded); C, every layer's $\Delta$ UB $< 0.05$ (not above co-expression), against the extraction cells (C$_{\mathrm{3k}}$: against the 3k tier only). $^{*}$0.0502 at $L_{2}$; two independent bootstrap seeds give 0.0528 and 0.0537, so C fails by less than the Monte-Carlo error. Pre-declared replication (same sign, and a lung (ext.) CI excluding 0.5): the sign agrees at 12/12 layers and the CI excludes 0.5 at 0/12, so it is not met. \textbf{B}: Shortcut rows score a pair by its target's TRRUST in-degree (no layer). ChIP-seq rows: positives, DoRothEA ChIP-seq pairs (level D) absent from TRRUST and from DoRothEA levels A--C; candidates constructed as for TRRUST, excluding in either direction every pair documented in any regulatory resource at any level, the ChIP-seq positives included; about half of these candidates (49--55\%) could not be matched on in-degree, so the fully matched rows keep only candidates matched on both deciles (and the positives of factors that have at least one; CSV column \texttt{chip\_candidates}). $^{a}$Positives in strata that hold both classes. $^{b}$$|r|$ AUROC, and $\Delta$, computed with exactly tied $|r|$ values ordered by floating-point noise rather than counted as one half (effect on the $|r|$ AUROC, and hence on $\Delta$, $\leq 0.00012$ where checked). Contamination: AUROC and UB corrected for a fraction $\pi$ of undocumented edges among the candidates, $0.5 + (A - 0.5)/(1 - \pi)$; seventh column, layers whose corrected UB is below 0.60. $^{c}$Mean of the per-factor AUROCs with a factor-bootstrap CI; last column, fraction of factors above 0.5 (range over layers). $^{d}$Covariates: the pair's log $|r|$ and the target's log mean expression, log $n_{\mathrm{obs}}$, STRING physical degree and GO annotation count (factor-level covariates cancel within the factor stratum); two-way bootstrap with refitting, $B = 500$. Residualised: AUROC of the cosine after OLS residualisation on the covariates. Conditional logit: odds ratio (OR) per pooled SD of the cosine, adjusted (unadjusted in the eighth column); layer mean, mean of the twelve coefficients. $^{e}$Effective number of factors under each statistic's weights ($n_{\mathrm{pos}} n_{\mathrm{cand}}$ per factor; Fisher information) and the share of the five largest. \textbf{C}: post hoc subsets, no multiplicity adjustment; $\Delta_{\mathrm{3k}}$, difference from $|r|$ on 3{,}000 cells. PAX5$\to$CD79A (illustration): detected, fraction of the extraction's cells expressing the gene; rank 1 is the highest cosine among PAX5's documented and candidate pairs ($n$); candidates below, share of PAX5's candidates scored lower. \textbf{D}: each split holds out a random fifth of the 245 kidney~A TRRUST pairs (49 positives, three candidate non-edges each); nothing is trained, so the spread over splits is the sampling distribution of the test statistic, not a CI. Uncentred, cosine of raw vectors, as the screen scored pairs; zero-row SVD, the reference decomposition of the full kidney~A matrix including the unobserved all-zero gene rows (Table~S14). Reference split: one fixed 80/20 split of these pairs; its value (first row with its i.i.d.\ CI) and its percentile among the 500 splits. MDE as in A, for the pooled AUROC (positives against all held-out non-edges) and for the within-factor AUROC; size, rejection rate at a population AUROC of 0.5 (nominal 0.05); perm.\ SD, Wald test with the permutation SD; oracle SD, the simulated sampling SD. \textbf{E}: regulator importance, normalised per target, of each factor for each target: GENIE3, the extra-trees variant (300 trees, scikit-learn); GRNBoost2, GRNBoost2-style gradient boosting (scikit-learn, early stopping on a tenth of the cells); T1, five seeds on the extraction cells; T2, two seeds on 3{,}000 cells; AUROC of the seed-averaged score, and mean (SD) of the per-seed AUROCs; $\Delta$, the method's AUROC minus that of $|r|$ on the same cells or of the cosine at the stated layer, on the same bootstrap draws. Cosine $L_0$ (first column): scGPT on the same pairs. \textbf{F}: L2-penalised logistic probes on features of the factor's and the target's vectors $x_i$, $x_j$ (bilinear: outer product of 16 principal components fitted on training genes; $[x_j, \cdot]$: target block added); AUROC, mean over 25 test folds; CI, corrected resampled $t$ (Nadeau--Bengio) with the test-to-training ratio counted in pairs (about 0.0625) or in genes (0.25, the cautious reading). $-$ target-only: minus, fold by fold, the better of the target-only and the factor-only probe (the latter scores 0.5 within factor). Rewired labels: degree-preserving swaps of targets between factors' positives, 100 draws, each retrained on one repetition's folds; $p = (b + 1)/(N + 1)$, floor 0.0099; max-$T$ over layers. Difference in differences: bilinear minus concatenation AUROC on the true labels against the same difference on each rewired draw. Input gene-token embedding: scGPT whole-human checkpoint; ESM2-15B: 5{,}120-dimensional per-protein embedding (UCE); both without layers. Non-linear: histogram gradient boosting and a two-layer perceptron, one repetition; rows, probe $\times$ layer $\times$ feature rows whose interval (pair ratio) lies above or below the model's own target-only control. \textbf{G}: fractions among untied pairs; sign test, exact two-sided; regulator-cluster CI, bootstrap over source factors; probe CIs, binomial (pairs independent) or gene-cluster bootstrap ($B = 2{,}000$); strict subset, reverse absent from every regulatory set; out-/in-degree, the same fraction for TRRUST degrees (no model). Full values: \texttt{Table\_S7\_edges.csv}.}\label{tab:S7}\\
\toprule
\endfirsthead
\multicolumn{10}{@{}l}{\hbox to 0pt{\footnotesize Table~\thetable\ (continued)\hss}}\\
\toprule
\endhead
\midrule
\multicolumn{10}{r@{}}{\hbox to 0pt{\hss\footnotesize continued on next page}}\\
\endfoot
\bottomrule
\endlastfoot
\multicolumn{10}{@{}p{16.923cm}@{}}{\footnotesize\textbf{A. Primary benchmark: within-factor AUROC of the full-dimensional signed cosine in all 16 extractions}}\\*
Extraction & Factors; positives; candidates & AUROC, range $L_0$--$L_{11}$ & Best layer: AUROC [95\% CI] & Max.\ UB & Min.\ max-$T$ $p$ & MDE: AUROC; $\Delta$ & $|r|$ AUROC: extr.; 3k & Max.\ $\Delta$ UB: extr.; 3k & Verdict\\*
\noalign{\nobreak\vskip 1.5pt\hrule height \cmidrulewidth\nobreak\vskip 2pt\nobreak}
lung & 312; 5{,}266; 47{,}160 & 0.516--0.542 & $L_{1}$ 0.542 [0.496,\,0.589] & 0.583 & $1.0\times10^{-4}$ & 0.564; 0.071 & 0.508; 0.515 & 0.074; 0.068 & D, U\\
lung (ext.) & 315; 5{,}093; 45{,}854 & 0.518--0.538 & $L_{1}$ 0.538 [0.493,\,0.588] & 0.581 & $1.0\times10^{-4}$ & 0.562; 0.072 & 0.524; 0.551 & 0.058; 0.033 & D, U; C$_{\mathrm{3k}}$\\
kidney~B & 254; 3{,}668; 32{,}712 & 0.524--0.554 & $L_{0}$ 0.554 [0.499,\,0.603] & 0.597 & $1.0\times10^{-4}$ & 0.576; 0.085 & 0.515; 0.522 & 0.090; 0.084 & D, U\\
kidney (whole-human ckpt) & 274; 4{,}111; 36{,}798 & 0.519--0.542 & $L_{0}$ 0.542 [0.491,\,0.593] & 0.587 & $1.0\times10^{-4}$ & 0.575; 0.091 & 0.510; 0.522 & 0.078; 0.067 & D, U\\
kidney (kidney ckpt) & 274; 4{,}111; 36{,}798 & 0.513--0.521 & $L_{2}$ 0.521 [0.475,\,0.568] & 0.561 & 0.0020 & 0.567; 0.084 & 0.510; 0.522 & 0.050$^{*}$; 0.040 & D, U; C$_{\mathrm{3k}}$\\
kidney (brain ckpt) & 274; 4{,}111; 36{,}798 & 0.516--0.523 & $L_{4}$ 0.523 [0.464,\,0.573] & 0.566 & $4.0\times10^{-4}$ & 0.564; 0.090 & 0.510; 0.522 & 0.069; 0.055 & D, U\\
immune & 36; 547; 3{,}144 & 0.577--0.638 & $L_{0}$ 0.638 [0.574,\,0.715] & 0.700 & $1.0\times10^{-4}$ & 0.601; 0.162 & 0.559; 0.562 & 0.151; 0.141 & D\\
immune (seed 43) & 34; 515; 2{,}877 & 0.584--0.651 & $L_{0}$ 0.651 [0.575,\,0.726] & 0.712 & $1.0\times10^{-4}$ & 0.583; 0.113 & 0.565; 0.582 & 0.156; 0.140 & D\\
immune (seed 44) & 36; 538; 3{,}074 & 0.570--0.644 & $L_{0}$ 0.644 [0.572,\,0.712] & 0.701 & $1.0\times10^{-4}$ & 0.590; 0.122 & 0.545; 0.574 & 0.166; 0.140 & D\\
immune (donor TSP14) & 34; 489; 2{,}711 & 0.588--0.650 & $L_{1}$ 0.650 [0.575,\,0.722] & 0.709 & $1.0\times10^{-4}$ & 0.597; 0.131 & 0.580; 0.572 & 0.139; 0.149 & D\\
immune (donor TSP21) & 34; 498; 2{,}796 & 0.571--0.646 & $L_{1}$ 0.646 [0.580,\,0.726] & 0.709 & $1.0\times10^{-4}$ & 0.595; 0.135 & 0.547; 0.561 & 0.197; 0.157 & D\\
immune (donor TSP25) & 35; 467; 2{,}498 & 0.562--0.627 & $L_{1}$ 0.627 [0.558,\,0.692] & 0.684 & $1.0\times10^{-4}$ & 0.589; 0.111 & 0.572; 0.574 & 0.143; 0.136 & D\\
immune (B cells) & 27; 330; 1{,}531 & 0.549--0.633 & $L_{0}$ 0.633 [0.555,\,0.711] & 0.697 & $1.0\times10^{-4}$ & 0.608; 0.163 & 0.531; 0.560 & 0.187; 0.159 & D\\
immune (CD4 T cells) & 29; 384; 1{,}791 & 0.561--0.618 & $L_{1}$ 0.618 [0.534,\,0.697] & 0.685 & $1.0\times10^{-4}$ & 0.603; 0.191 & 0.508; 0.528 & 0.193; 0.167 & D\\
immune (CD8 T cells) & 31; 440; 2{,}330 & 0.589--0.647 & $L_{1}$ 0.647 [0.577,\,0.731] & 0.717 & $1.0\times10^{-4}$ & 0.593; 0.169 & 0.553; 0.551 & 0.213; 0.209 & D\\
kidney~A & 29; 213; 1{,}446 & 0.601--0.652 & $L_{1}$ 0.652 [0.561,\,0.747] & 0.731 & $1.0\times10^{-4}$ & 0.632; 0.157 & 0.559; 0.577 & 0.195; 0.169 & D\\
\midrule
\multicolumn{10}{@{}p{16.923cm}@{}}{\footnotesize\textbf{B. Sensitivity analyses in the two lung extractions (cosine at $L_0$, and over the twelve layers)}}\\*
Analysis & Extraction & Factors; positives & $L_0$ [95\% CI] & Range $L_0$--$L_{11}$ & Max.\ UB & Layers with CI $>0.5$ & $|r|$ AUROC, same pairs & $\Delta$ vs $|r|$, range & Max.\ $\Delta$ UB\\*
\noalign{\nobreak\vskip 1.5pt\hrule height \cmidrulewidth\nobreak\vskip 2pt\nobreak}
\multirow[t]{4}{1.9cm}{\scriptsize\raggedright All pairs (primary)} & lung & 312; 5{,}266 & 0.541 [0.497,\,0.588] & 0.516--0.542 & 0.583 & 0/12 & 0.508 & $+$0.008 to $+$0.034 & 0.074\\*
 & lung (ext.) & 315; 5{,}093 & 0.538 [0.496,\,0.586] & 0.518--0.538 & 0.581 & 0/12 & 0.524 & $-$0.006 to $+$0.014 & 0.058\\
\multirow[t]{4}{1.9cm}{\scriptsize\raggedright Shortcut: target in-degree, incl.\ own edge} & lung & 312; 5{,}266 & 0.733 [0.658,\,0.799] & (no layer) & 0.787 & -- & -- & -- & --\\*
 & lung (ext.) & 315; 5{,}093 & 0.750 [0.675,\,0.818] & (no layer) & 0.803 & -- & -- & -- & --\\
\multirow[t]{4}{1.9cm}{\scriptsize\raggedright Shortcut: target's other-factor in-degree} & lung & 312; 5{,}266 & 0.626 [0.551,\,0.702] & (no layer) & 0.687 & -- & -- & -- & --\\*
 & lung (ext.) & 315; 5{,}093 & 0.633 [0.554,\,0.711] & (no layer) & 0.696 & -- & -- & -- & --\\
\multirow[t]{4}{1.9cm}{\scriptsize\raggedright Strata: factor $\times$ other-factor in-degree} & lung & --; 3{,}564$^{a}$ & 0.530 [0.471,\,0.607] & 0.510--0.531 & 0.597 & 0/12 & 0.504$^{b}$ & $+$0.005 to $+$0.027$^{b}$ & 0.093$^{b}$\\*
 & lung (ext.) & --; 3{,}195$^{a}$ & 0.542 [0.485,\,0.619] & 0.522--0.542 & 0.606 & 0/12 & 0.524$^{b}$ & $-$0.002 to $+$0.018$^{b}$ & 0.084$^{b}$\\
\multirow[t]{4}{1.9cm}{\scriptsize\raggedright ChIP-seq positives (not in TRRUST)} & lung & 225; 57{,}017 & 0.528 [0.506,\,0.551] & 0.528--0.538 & 0.559 & 12/12 & 0.542 & $-$0.014 to $-$0.004 & 0.018\\*
 & lung (ext.) & 218; 54{,}332 & 0.533 [0.510,\,0.557] & 0.533--0.541 & 0.561 & 12/12 & 0.515 & $+$0.018 to $+$0.026 & 0.047\\
\multirow[t]{4}{1.9cm}{\scriptsize\raggedright ChIP-seq positives, fully matched candidates} & lung & 224; 57{,}014 & 0.529 [0.511,\,0.545] & 0.521--0.531 & 0.544 & 12/12 & 0.516 & $+$0.005 to $+$0.014 & 0.030\\*
 & lung (ext.) & 218; 54{,}332 & 0.528 [0.510,\,0.545] & 0.526--0.531 & 0.546 & 12/12 & 0.507 & $+$0.019 to $+$0.024 & 0.040\\
\multirow[t]{4}{1.9cm}{\scriptsize\raggedright $\geq 2$ publications} & lung & 131; 579 & 0.591 [0.507,\,0.677] & 0.535--0.593 & 0.669 & 2/12 & 0.511$^{b}$ & $+$0.024 to $+$0.081$^{b}$ & 0.163$^{b}$\\*
 & lung (ext.) & 129; 553 & 0.588 [0.506,\,0.672] & 0.539--0.588 & 0.672 & 2/12 & 0.551$^{b}$ & $-$0.012 to $+$0.037$^{b}$ & 0.135$^{b}$\\
\multirow[t]{4}{1.9cm}{\scriptsize\raggedright 1 publication} & lung & 312; 4{,}687 & 0.531 [0.483,\,0.576] & 0.512--0.532 & 0.571 & 0/12 & 0.508$^{b}$ & $+$0.005 to $+$0.024$^{b}$ & 0.064$^{b}$\\*
 & lung (ext.) & 315; 4{,}540 & 0.528 [0.485,\,0.575] & 0.513--0.528 & 0.570 & 0/12 & 0.519$^{b}$ & $-$0.005 to $+$0.010$^{b}$ & 0.051$^{b}$\\
\multirow[t]{4}{1.9cm}{\scriptsize\raggedright Factor is input in $\geq 10\%$ of cells} & lung & 94; 1{,}753 & 0.566 [0.515,\,0.642] & 0.557--0.584 & 0.633 & 12/12 & 0.528 & $+$0.029 to $+$0.056 & 0.117\\*
 & lung (ext.) & 111; 1{,}921 & 0.560 [0.514,\,0.613] & 0.541--0.567 & 0.610 & 12/12 & 0.528 & $+$0.012 to $+$0.038 & 0.090\\
\multirow[t]{4}{1.9cm}{\scriptsize\raggedright Factor is input in $<10\%$ of cells} & lung & 218; 3{,}513 & 0.527 [0.479,\,0.584] & 0.493--0.527 & 0.578 & 0/12 & 0.497 & $-$0.004 to $+$0.030 & 0.075\\*
 & lung (ext.) & 204; 3{,}172 & 0.524 [0.474,\,0.590] & 0.503--0.524 & 0.582 & 0/12 & 0.521 & $-$0.019 to $+$0.003 & 0.056\\
\multirow[t]{4}{1.9cm}{\scriptsize\raggedright Factor expressed in $\geq 10\%$ of cells} & lung & 265; 4{,}669 & 0.539 [0.488,\,0.585] & 0.515--0.540 & 0.582 & 0/12 & 0.506 & $+$0.008 to $+$0.034 & 0.078\\*
 & lung (ext.) & 227; 3{,}955 & 0.530 [0.483,\,0.586] & 0.512--0.530 & 0.582 & 0/12 & 0.517 & $-$0.005 to $+$0.014 & 0.067\\
\multirow[t]{4}{1.9cm}{\scriptsize\raggedright Factor expressed in $<10\%$ of cells} & lung & 47; 597 & 0.577 [0.507,\,0.651] & 0.536--0.577 & 0.641 & 3/12 & 0.538 & $-$0.002 to $+$0.038 & 0.111\\*
 & lung (ext.) & 88; 1{,}138 & 0.581 [0.527,\,0.648] & 0.550--0.583 & 0.637 & 8/12 & 0.567 & $-$0.016 to $+$0.016 & 0.075\\
\multirow[t]{4}{1.9cm}{\scriptsize\raggedright Contamination $\pi = 0.05$} & lung & -- & 0.543 & 0.517--0.544 & 0.587 & UB $<0.60$: 12/12 & -- & -- & --\\*
 & lung (ext.) & -- & 0.540 & 0.518--0.540 & 0.585 & UB $<0.60$: 12/12 & -- & -- & --\\
\multirow[t]{4}{1.9cm}{\scriptsize\raggedright Contamination $\pi = 0.10$} & lung & -- & 0.546 & 0.518--0.547 & 0.592 & UB $<0.60$: 12/12 & -- & -- & --\\*
 & lung (ext.) & -- & 0.542 & 0.520--0.542 & 0.590 & UB $<0.60$: 12/12 & -- & -- & --\\
\multirow[t]{4}{1.9cm}{\scriptsize\raggedright Contamination $\pi = 0.20$} & lung & -- & 0.551 & 0.520--0.552 & 0.604 & UB $<0.60$: 9/12 & -- & -- & --\\*
 & lung (ext.) & -- & 0.547 & 0.522--0.548 & 0.601 & UB $<0.60$: 9/12 & -- & -- & --\\
\multirow[t]{4}{1.9cm}{\scriptsize\raggedright Macro AUROC$^{c}$} & lung & 312; -- & 0.580 [0.565,\,0.596] & 0.554--0.580 & 0.593 & 12/12 & 0.518 & \multicolumn{2}{>{\scriptsize\raggedright\arraybackslash}p{3.248cm}@{}}{factors $>0.5$: 68--75\%}\\*
 & lung (ext.) & 315; -- & 0.573 [0.558,\,0.588] & 0.553--0.573 & 0.586 & 12/12 & 0.520 & \multicolumn{2}{>{\scriptsize\raggedright\arraybackslash}p{3.248cm}@{}}{factors $>0.5$: 69--72\%}\\
\multirow[t]{4}{1.9cm}{\scriptsize\raggedright Residualised$^{d}$} & lung & 312; 5{,}266 & 0.536 [0.495,\,0.581] & 0.512--0.537 & 0.577 & 0/12 & -- & \multicolumn{2}{>{\scriptsize\raggedright\arraybackslash}p{3.248cm}@{}}{resid.\ $-$ raw: $-$0.005 to $-$0.004}\\*
 & lung (ext.) & 315; 5{,}093 & 0.533 [0.489,\,0.582] & 0.515--0.533 & 0.579 & 0/12 & -- & \multicolumn{2}{>{\scriptsize\raggedright\arraybackslash}p{3.248cm}@{}}{resid.\ $-$ raw: $-$0.005 to $-$0.002}\\
\multirow[t]{4}{1.9cm}{\scriptsize\raggedright Conditional logit OR$^{d}$} & lung & 312; 5{,}266 & 1.238 [1.148,\,1.348] & 1.123--1.239 & -- & CI $>1$: 12/12 & unadj.\ 1.132--1.247 & \multicolumn{2}{>{\scriptsize\raggedright\arraybackslash}p{3.248cm}@{}}{layer mean 1.189 [1.114,\,1.281]; unadj.\ 1.198}\\*
 & lung (ext.) & 315; 5{,}093 & 1.218 [1.119,\,1.336] & 1.132--1.219 & -- & CI $>1$: 12/12 & unadj.\ 1.140--1.230 & \multicolumn{2}{>{\scriptsize\raggedright\arraybackslash}p{3.248cm}@{}}{layer mean 1.184 [1.098,\,1.294]; unadj.\ 1.193}\\
\multirow[t]{4}{1.9cm}{\scriptsize\raggedright Effective no.\ of factors$^{e}$} & lung & 312; -- & \multicolumn{2}{>{\scriptsize\raggedright\arraybackslash}p{3.248cm}}{pooled AUROC 14.9 (top 5: 51\%)} & \multicolumn{2}{>{\scriptsize\raggedright\arraybackslash}p{3.248cm}}{cond.\ logit 75.8 (top 5: 19\%)} & \multicolumn{3}{>{\scriptsize\raggedright\arraybackslash}p{4.925cm}@{}}{}\\*
 & lung (ext.) & 315; -- & \multicolumn{2}{>{\scriptsize\raggedright\arraybackslash}p{3.248cm}}{pooled AUROC 14.6 (top 5: 51\%)} & \multicolumn{2}{>{\scriptsize\raggedright\arraybackslash}p{3.248cm}}{cond.\ logit 73.5 (top 5: 19\%)} & \multicolumn{3}{>{\scriptsize\raggedright\arraybackslash}p{4.925cm}@{}}{}\\
\midrule
\multicolumn{10}{@{}p{16.923cm}@{}}{\footnotesize\textbf{C. Cell-type extractions restricted to factors detected in $\geq 10\%$ of their cells, and the PAX5$\to$CD79A example}}\\*
Subset & Factors; positives; candidates & $L_0$ [95\% CI] & Best layer [95\% CI] & Range $L_0$--$L_{11}$ & Max.\ UB & $|r|$ extr.\ [95\% CI] & $\Delta$ at $L_0$ [95\% CI] & Layers with $\Delta$ CI $>0$ & $\Delta_{\mathrm{3k}}$ at $L_0$ [95\% CI]\\*
\noalign{\nobreak\vskip 1.5pt\hrule height \cmidrulewidth\nobreak\vskip 2pt\nobreak}
immune (B cells) & 21; 292; 1{,}200 & 0.629 [0.537,\,0.707] & $L_{0}$ 0.629 [0.537,\,0.707] & 0.543--0.629 & 0.694 & 0.519 [0.423,\,0.621] & $+$0.110 [$-$0.010, 0.221] & 0/12 & $+$0.080 [$-$0.040, 0.187]\\
immune (CD4 T cells) & 22; 319; 1{,}290 & 0.605 [0.520,\,0.685] & $L_{1}$ 0.606 [0.520,\,0.685] & 0.555--0.606 & 0.670 & 0.487 [0.414,\,0.569] & $+$0.118 [0.014,\,0.222] & 3/12 ($L_{0}$--$L_{2}$) & $+$0.092 [$-$0.023, 0.185]\\
immune (CD8 T cells) & 24; 353; 1{,}610 & 0.622 [0.550,\,0.689] & $L_{1}$ 0.627 [0.556,\,0.693] & 0.576--0.627 & 0.683 & 0.559 [0.462,\,0.663] & $+$0.064 [$-$0.053, 0.180] & 0/12 & $+$0.073 [$-$0.030, 0.171]\\
\multicolumn{10}{@{}p{16.923cm}@{}}{\scriptsize\textit{PAX5$\to$CD79A: rank of the documented pair among all of PAX5's pairs (documented and candidate), over $L_0$--$L_{11}$}}\\*
Extraction & Detected: PAX5; CD79A & Rank ($n$ pairs) & Candidates below & $|r|$ rank: extr.; 3k & Extraction & Detected: PAX5; CD79A & Rank ($n$ pairs) & Candidates below & $|r|$ rank: extr.; 3k\\*
\noalign{\nobreak\vskip 1.5pt\hrule height \cmidrulewidth\nobreak\vskip 2pt\nobreak}
immune (B cells) & 78.5\%; 91.4\% & 5--7 (94) & 95--98\% & 5; 53 & immune (CD4 T cells) & 0.8\%; 3.1\% & 3--5 (86) & 96--99\% & 57; 73\\
immune (CD8 T cells) & 0.4\%; 0.4\% & 1--4 (69) & 97--100\% & 61; 6 & immune & 13.3\%; 21.1\% & 5 (108) & 98\% & 3; 3\\
immune (seed 43) & 14.8\%; 24.6\% & 4 (110) & 99\% & 3; 3 & immune (seed 44) & 15.6\%; 21.1\% & 5 (110) & 98\% & 2; 3\\
immune (donor TSP14) & 9.0\%; 15.2\% & 4 (87) & 99\% & 3; 3 & immune (donor TSP21) & 25.4\%; 30.1\% & 4--5 (103) & 98--99\% & 3; 3\\
immune (donor TSP25) & 26.6\%; 34.0\% & 5--8 (108) & 95--98\% & 1; 3 & lung & 1.6\%; 0.4\% & 2--28 (176) & 86--100\% & 54; 1\\
lung (ext.) & 1.2\%; 2.7\% & 4--6 (198) & 99--100\% & 1; 1 & kidney~B & 6.2\%; 6.8\% & 4--7 (176) & 98--100\% & 2; 2\\
kidney~A & 3.5\%; 2.7\% & 5--19 (74) & 77--98\% & 3; 3 & kidney (whole-human ckpt) & 7.8\%; 6.8\% & 4--9 (187) & 96--99\% & 1; 2\\
kidney (kidney ckpt) & 7.8\%; 6.8\% & 2--14 (187) & 93--99\% & 1; 2 & kidney (brain ckpt) & 7.8\%; 6.8\% & 5--97 (187) & 49--99\% & 1; 2\\
\multicolumn{10}{@{}>{\scriptsize\raggedright\arraybackslash}p{16.923cm}@{}}{All nine immune extractions: rank 1--8 at every layer.}\\
\midrule
\multicolumn{10}{@{}p{16.923cm}@{}}{\footnotesize\textbf{D. An 80/20 held-out split of the kidney~A pairs (49 positives, 147 candidates, 27 factors, 117 targets): 500 random splits}}\\*
\multicolumn{2}{@{}>{\scriptsize\raggedright\arraybackslash}p{3.504cm}}{Scorer ($L_0$)} & Quantity & Mean & SD & 2.5--97.5\% & Min--max & Fraction $>0.5$ ($>0$ for $\Delta$) & \multicolumn{2}{>{\scriptsize\raggedright\arraybackslash}p{3.248cm}@{}}{Reference split (percentile)}\\*
\noalign{\nobreak\vskip 1.5pt\hrule height \cmidrulewidth\nobreak\vskip 2pt\nobreak}
\multicolumn{2}{@{}>{\scriptsize\raggedright\arraybackslash}p{3.504cm}}{Full 512-d, uncentred (screen convention)} & AUROC & 0.580 & 0.043 & 0.501--0.664 & 0.458--0.699 & 0.974 & \multicolumn{2}{>{\scriptsize\raggedright\arraybackslash}p{3.248cm}@{}}{0.591 [0.499,\,0.684] (60.8th)}\\
\multicolumn{2}{@{}>{\scriptsize\raggedright\arraybackslash}p{3.504cm}}{} & $\Delta$ vs $|r|$ & $+$0.081 & 0.047 & $-$0.014 to $+$0.169 & $-$0.075 to $+$0.230 & 0.954 & \multicolumn{2}{>{\scriptsize\raggedright\arraybackslash}p{3.248cm}@{}}{$+$0.035 (15.4th)}\\
\multicolumn{2}{@{}>{\scriptsize\raggedright\arraybackslash}p{3.504cm}}{Full 512-d, $V$-centred} & AUROC & 0.633 & 0.040 & 0.553--0.710 & 0.514--0.740 & 1.000 & \multicolumn{2}{>{\scriptsize\raggedright\arraybackslash}p{3.248cm}@{}}{--}\\
\multicolumn{2}{@{}>{\scriptsize\raggedright\arraybackslash}p{3.504cm}}{} & $\Delta$ vs $|r|$ & $+$0.134 & 0.049 & $+$0.036 to $+$0.230 & $-$0.023 to $+$0.296 & 0.996 & \multicolumn{2}{>{\scriptsize\raggedright\arraybackslash}p{3.248cm}@{}}{--}\\
\multicolumn{2}{@{}>{\scriptsize\raggedright\arraybackslash}p{3.504cm}}{$\SV{2}$--$\SV{7}$, zero-row SVD (reference basis)} & AUROC & 0.598 & 0.040 & 0.525--0.673 & 0.450--0.714 & 0.994 & \multicolumn{2}{>{\scriptsize\raggedright\arraybackslash}p{3.248cm}@{}}{0.589}\\
\multicolumn{2}{@{}>{\scriptsize\raggedright\arraybackslash}p{3.504cm}}{} & $\Delta$ vs $|r|$ & $+$0.098 & 0.054 & $-$0.006 to $+$0.202 & $-$0.058 to $+$0.270 & 0.962 & \multicolumn{2}{>{\scriptsize\raggedright\arraybackslash}p{3.248cm}@{}}{--}\\
\multicolumn{2}{@{}>{\scriptsize\raggedright\arraybackslash}p{3.504cm}}{$\SV{2}$--$\SV{7}$, observed-gene SVD} & AUROC & 0.607 & 0.040 & 0.532--0.682 & 0.461--0.729 & 0.998 & \multicolumn{2}{>{\scriptsize\raggedright\arraybackslash}p{3.248cm}@{}}{--}\\
\multicolumn{2}{@{}>{\scriptsize\raggedright\arraybackslash}p{3.504cm}}{} & $\Delta$ vs $|r|$ & $+$0.107 & 0.053 & $+$0.000 to $+$0.206 & $-$0.046 to $+$0.291 & 0.974 & \multicolumn{2}{>{\scriptsize\raggedright\arraybackslash}p{3.248cm}@{}}{--}\\
\multicolumn{2}{@{}>{\scriptsize\raggedright\arraybackslash}p{3.504cm}}{$\SV{5}$--$\SV{7}$, zero-row SVD (reference basis)} & AUROC & 0.545 & 0.043 & 0.463--0.627 & 0.424--0.663 & 0.854 & \multicolumn{2}{>{\scriptsize\raggedright\arraybackslash}p{3.248cm}@{}}{0.541}\\
\multicolumn{2}{@{}>{\scriptsize\raggedright\arraybackslash}p{3.504cm}}{} & $\Delta$ vs $|r|$ & $+$0.046 & 0.060 & $-$0.070 to $+$0.158 & $-$0.135 to $+$0.232 & 0.800 & \multicolumn{2}{>{\scriptsize\raggedright\arraybackslash}p{3.248cm}@{}}{--}\\
\multicolumn{2}{@{}>{\scriptsize\raggedright\arraybackslash}p{3.504cm}}{$\SV{5}$--$\SV{7}$, observed-gene SVD} & AUROC & 0.567 & 0.044 & 0.485--0.648 & 0.438--0.689 & 0.926 & \multicolumn{2}{>{\scriptsize\raggedright\arraybackslash}p{3.248cm}@{}}{--}\\
\multicolumn{2}{@{}>{\scriptsize\raggedright\arraybackslash}p{3.504cm}}{} & $\Delta$ vs $|r|$ & $+$0.068 & 0.057 & $-$0.042 to $+$0.189 & $-$0.103 to $+$0.234 & 0.872 & \multicolumn{2}{>{\scriptsize\raggedright\arraybackslash}p{3.248cm}@{}}{--}\\
\multicolumn{2}{@{}>{\scriptsize\raggedright\arraybackslash}p{3.504cm}}{$|r|$, 3{,}000 cells} & AUROC & 0.500 & 0.044 & 0.408--0.581 & 0.344--0.599 & 0.496 & \multicolumn{2}{>{\scriptsize\raggedright\arraybackslash}p{3.248cm}@{}}{0.556 (89.6th)}\\
\multicolumn{2}{@{}>{\scriptsize\raggedright\arraybackslash}p{3.504cm}}{Random scores} & AUROC & 0.500 & 0.049 & 0.407--0.600 & 0.349--0.642 & 0.486 & \multicolumn{2}{>{\scriptsize\raggedright\arraybackslash}p{3.248cm}@{}}{0.513}\\
\multicolumn{2}{@{}>{\scriptsize\raggedright\arraybackslash}p{3.504cm}}{} & $\Delta$ vs $|r|$ & $+$0.000 & 0.067 & $-$0.130 to $+$0.132 & $-$0.191 to $+$0.178 & 0.492 & \multicolumn{2}{>{\scriptsize\raggedright\arraybackslash}p{3.248cm}@{}}{--}\\
\multicolumn{2}{@{}>{\scriptsize\raggedright\arraybackslash}p{3.504cm}}{Full 512-d, uncentred (screen convention), within-factor AUROC} & AUROC & 0.579 & 0.053 & 0.484--0.682 & 0.367--0.741 & 0.944 & \multicolumn{2}{>{\scriptsize\raggedright\arraybackslash}p{3.248cm}@{}}{--}\\
\multicolumn{10}{@{}p{16.923cm}@{}}{\scriptsize\textit{Minimum detectable effect of that design (80\% power, two-sided $\alpha = 0.05$; target ICC 0.47, factor ICC 0.17; 400 simulations per grid point)}}\\*
\multicolumn{2}{@{}>{\scriptsize\raggedright\arraybackslash}p{3.504cm}}{Statistic} & Test & MDE AUROC & Size at null & MDE $\Delta$ & Size at null ($\Delta$) & \multicolumn{3}{>{\scriptsize\raggedright\arraybackslash}p{4.925cm}@{}}{}\\*
\noalign{\nobreak\vskip 1.5pt\hrule height \cmidrulewidth\nobreak\vskip 2pt\nobreak}
\multicolumn{2}{@{}>{\scriptsize\raggedright\arraybackslash}p{3.504cm}}{Pooled AUROC (positives against all held-out non-edges)} & two-way bootstrap & 0.716 & 0.022 & 0.254 & 0.007 & \multicolumn{3}{>{\scriptsize\raggedright\arraybackslash}p{4.925cm}@{}}{}\\
\multicolumn{2}{@{}>{\scriptsize\raggedright\arraybackslash}p{3.504cm}}{} & factor-only bootstrap & 0.651 & 0.163 & 0.173 & 0.152 & \multicolumn{3}{>{\scriptsize\raggedright\arraybackslash}p{4.925cm}@{}}{}\\
\multicolumn{2}{@{}>{\scriptsize\raggedright\arraybackslash}p{3.504cm}}{} & i.i.d.\ DeLong & 0.644 & 0.195 & 0.165 & 0.158 & \multicolumn{3}{>{\scriptsize\raggedright\arraybackslash}p{4.925cm}@{}}{}\\
\multicolumn{2}{@{}>{\scriptsize\raggedright\arraybackslash}p{3.504cm}}{} & Wald, perm.\ SD & 0.643 & 0.188 & -- & -- & \multicolumn{3}{>{\scriptsize\raggedright\arraybackslash}p{4.925cm}@{}}{}\\
\multicolumn{2}{@{}>{\scriptsize\raggedright\arraybackslash}p{3.504cm}}{} & oracle SD & 0.700 & -- & 0.216 & -- & \multicolumn{3}{>{\scriptsize\raggedright\arraybackslash}p{4.925cm}@{}}{}\\
\multicolumn{2}{@{}>{\scriptsize\raggedright\arraybackslash}p{3.504cm}}{Within-factor AUROC} & two-way bootstrap & 0.795 & 0.003 & 0.359 & 0.003 & \multicolumn{3}{>{\scriptsize\raggedright\arraybackslash}p{4.925cm}@{}}{}\\
\multicolumn{2}{@{}>{\scriptsize\raggedright\arraybackslash}p{3.504cm}}{} & factor-only bootstrap & 0.716 & 0.025 & 0.261 & 0.028 & \multicolumn{3}{>{\scriptsize\raggedright\arraybackslash}p{4.925cm}@{}}{}\\
\multicolumn{2}{@{}>{\scriptsize\raggedright\arraybackslash}p{3.504cm}}{} & Wald, perm.\ SD & 0.671 & 0.092 & -- & -- & \multicolumn{3}{>{\scriptsize\raggedright\arraybackslash}p{4.925cm}@{}}{}\\
\multicolumn{2}{@{}>{\scriptsize\raggedright\arraybackslash}p{3.504cm}}{} & oracle SD & 0.695 & -- & 0.253 & -- & \multicolumn{3}{>{\scriptsize\raggedright\arraybackslash}p{4.925cm}@{}}{}\\
\multicolumn{2}{@{}>{\scriptsize\raggedright\arraybackslash}p{3.504cm}}{Either statistic} & Hanley--McNeil & 0.634 & -- & 0.161 & -- & \multicolumn{3}{>{\scriptsize\raggedright\arraybackslash}p{4.925cm}@{}}{}\\
\midrule
\multicolumn{10}{@{}p{16.923cm}@{}}{\footnotesize\textbf{E. GENIE3 (extra-trees) and GRNBoost2-style boosting on the same benchmark (T1, the extraction's own cells; T2, 3{,}000 cells)}}\\*
Extraction (cosine $L_0$) & Method, tier (cells) & AUROC, seed ensemble [95\% CI] & UB & Seed mean (SD; seeds) & $|r|$, same cells & $\Delta$ vs $|r|$ [95\% CI] & $\Delta$ vs cosine $L_0$ [95\% CI] & $\Delta$ vs cosine, range $L_0$--$L_{11}$ & T2 $-$ T1 [95\% CI]\\*
\noalign{\nobreak\vskip 1.5pt\hrule height \cmidrulewidth\nobreak\vskip 2pt\nobreak}
lung (0.541) & GENIE3 T1 (256) & 0.516 [0.488,\,0.550] & 0.545 & 0.516 (0.0031; 5) & 0.508 & $+$0.008 [$-$0.017, 0.036] & $-$0.025 [$-$0.065, 0.013] & $-$0.026 to $+$0.000 & --\\
 & GENIE3 T2 (3{,}000) & 0.513 [0.479,\,0.550] & 0.544 & 0.513 (0.0027; 2) & 0.515 & $-$0.002 [$-$0.032, 0.032] & $-$0.028 [$-$0.068, 0.011] & $-$0.029 to $-$0.003 & $-$0.003 [$-$0.029, 0.022]\\
 & GRNBoost2 T1 (256) & 0.518 [0.494,\,0.546] & 0.541 & 0.518 (0.0059; 5) & 0.508 & $+$0.010 [$-$0.022, 0.038] & $-$0.023 [$-$0.071, 0.026] & $-$0.024 to $+$0.002 & --\\
 & GRNBoost2 T2 (3{,}000) & 0.512 [0.483,\,0.544] & 0.539 & 0.515 (0.0020; 2) & 0.515 & $-$0.003 [$-$0.031, 0.025] & $-$0.029 [$-$0.073, 0.014] & $-$0.030 to $-$0.004 & $-$0.006 [$-$0.038, 0.026]\\
lung (ext.) (0.538) & GENIE3 T1 (256) & 0.535 [0.508,\,0.565] & 0.559 & 0.533 (0.0028; 5) & 0.524 & $+$0.011 [$-$0.017, 0.036] & $-$0.002 [$-$0.051, 0.036] & $-$0.003 to $+$0.018 & --\\
 & GENIE3 T2 (3{,}000) & 0.538 [0.501,\,0.585] & 0.576 & 0.538 (0.0003; 2) & 0.551 & $-$0.013 [$-$0.044, 0.022] & $+$0.001 [$-$0.047, 0.043] & $+$0.000 to $+$0.021 & $+$0.003 [$-$0.028, 0.041]\\
 & GRNBoost2 T1 (256) & 0.521 [0.493,\,0.551] & 0.547 & 0.520 (0.0070; 5) & 0.524 & $-$0.003 [$-$0.030, 0.025] & $-$0.016 [$-$0.063, 0.020] & $-$0.017 to $+$0.004 & --\\
 & GRNBoost2 T2 (3{,}000) & 0.535 [0.504,\,0.570] & 0.564 & 0.533 (0.0014; 2) & 0.551 & $-$0.017 [$-$0.048, 0.019] & $-$0.003 [$-$0.053, 0.036] & $-$0.004 to $+$0.017 & $+$0.013 [$-$0.022, 0.049]\\
kidney~B (0.554) & GENIE3 T1 (192) & 0.529 [0.496,\,0.564] & 0.559 & 0.527 (0.0024; 5) & 0.515 & $+$0.014 [$-$0.021, 0.049] & $-$0.026 [$-$0.081, 0.023] & $-$0.026 to $+$0.005 & --\\
 & GENIE3 T2 (3{,}000) & 0.528 [0.497,\,0.563] & 0.558 & 0.527 (0.0003; 2) & 0.522 & $+$0.006 [$-$0.034, 0.052] & $-$0.027 [$-$0.077, 0.023] & $-$0.027 to $+$0.004 & $-$0.001 [$-$0.036, 0.038]\\
 & GRNBoost2 T1 (192) & 0.519 [0.492,\,0.548] & 0.542 & 0.513 (0.0027; 5) & 0.515 & $+$0.004 [$-$0.029, 0.035] & $-$0.035 [$-$0.089, 0.017] & $-$0.035 to $-$0.004 & --\\
 & GRNBoost2 T2 (3{,}000) & 0.528 [0.501,\,0.557] & 0.553 & 0.527 (0.0003; 2) & 0.522 & $+$0.006 [$-$0.028, 0.044] & $-$0.027 [$-$0.076, 0.022] & $-$0.027 to $+$0.004 & $+$0.008 [$-$0.025, 0.045]\\
immune (0.638) & GENIE3 T1 (256) & 0.578 [0.508,\,0.664] & 0.649 & 0.577 (0.0048; 5) & 0.559 & $+$0.019 [$-$0.045, 0.082] & $-$0.060 [$-$0.137, 0.020] & $-$0.060 to $+$0.000 & --\\
 & GENIE3 T2 (3{,}000) & 0.610 [0.532,\,0.693] & 0.677 & 0.609 (0.0015; 2) & 0.562 & $+$0.048 [$-$0.031, 0.121] & $-$0.028 [$-$0.115, 0.055] & $-$0.028 to $+$0.032 & $+$0.032 [$-$0.027, 0.083]\\
 & GRNBoost2 T1 (256) & 0.574 [0.502,\,0.667] & 0.651 & 0.572 (0.0067; 5) & 0.559 & $+$0.015 [$-$0.048, 0.091] & $-$0.064 [$-$0.140, 0.021] & $-$0.064 to $-$0.003 & --\\
 & GRNBoost2 T2 (3{,}000) & 0.605 [0.529,\,0.685] & 0.672 & 0.604 (0.0087; 2) & 0.562 & $+$0.044 [$-$0.025, 0.115] & $-$0.032 [$-$0.116, 0.051] & $-$0.032 to $+$0.028 & $+$0.031 [$-$0.039, 0.095]\\
kidney~A (0.647) & GENIE3 T1 (256) & 0.557 [0.460,\,0.677] & 0.649 & 0.554 (0.0042; 5) & 0.559 & $-$0.002 [$-$0.096, 0.097] & $-$0.091 [$-$0.205, 0.030] & $-$0.095 to $-$0.044 & --\\
 & GENIE3 T2 (3{,}000) & 0.596 [0.504,\,0.701] & 0.683 & 0.597 (0.0030; 2) & 0.577 & $+$0.019 [$-$0.096, 0.144] & $-$0.051 [$-$0.153, 0.056] & $-$0.056 to $-$0.005 & $+$0.040 [$-$0.053, 0.132]\\
 & GRNBoost2 T1 (256) & 0.554 [0.449,\,0.670] & 0.647 & 0.544 (0.0126; 5) & 0.559 & $-$0.005 [$-$0.104, 0.097] & $-$0.094 [$-$0.214, 0.038] & $-$0.098 to $-$0.047 & --\\
 & GRNBoost2 T2 (3{,}000) & 0.595 [0.502,\,0.700] & 0.678 & 0.589 (0.0108; 2) & 0.577 & $+$0.018 [$-$0.089, 0.120] & $-$0.052 [$-$0.151, 0.053] & $-$0.057 to $-$0.006 & $+$0.041 [$-$0.054, 0.128]\\
\multicolumn{10}{@{}>{\scriptsize\raggedright\arraybackslash}p{16.923cm}@{}}{Two-sided intervals excluding 0: $\Delta$ vs $|r|$, 0 of 20; $\Delta$ vs cosine, 0 of 240 (extraction $\times$ method $\times$ tier $\times$ layer); one one-sided bound lies below 0 (kidney~A GENIE3 T1 vs $L_{1}$, UB $-$0.0008; unadjusted). One-sided UB of $\Delta$ vs $|r|$ in lung, lung (ext.) and kidney~B: at most 0.044.}\\
\midrule
\multicolumn{10}{@{}p{16.923cm}@{}}{\footnotesize\textbf{F. Supervised probes on pair-disjoint folds (neither gene of a test pair seen in training; 5 $\times$ 5 folds)}}\\*
\multicolumn{4}{@{}>{\scriptsize\raggedright\arraybackslash}p{6.859cm}}{Quantity} & \multicolumn{2}{>{\scriptsize\raggedright\arraybackslash}p{3.248cm}}{lung} & \multicolumn{2}{>{\scriptsize\raggedright\arraybackslash}p{3.248cm}}{lung (ext.)} & \multicolumn{2}{>{\scriptsize\raggedright\arraybackslash}p{3.248cm}@{}}{kidney~B}\\*
\noalign{\nobreak\vskip 1.5pt\hrule height \cmidrulewidth\nobreak\vskip 2pt\nobreak}
\multicolumn{4}{@{}>{\scriptsize\raggedright\arraybackslash}p{6.859cm}}{Benchmark (each test fold holds about 4\% of the pairs)} & \multicolumn{2}{>{\scriptsize\raggedright\arraybackslash}p{3.248cm}}{52{,}426 pairs; 5{,}266 positives} & \multicolumn{2}{>{\scriptsize\raggedright\arraybackslash}p{3.248cm}}{50{,}947 pairs; 5{,}093 positives} & \multicolumn{2}{>{\scriptsize\raggedright\arraybackslash}p{3.248cm}@{}}{36{,}380 pairs; 3{,}668 positives}\\
\multicolumn{4}{@{}>{\scriptsize\raggedright\arraybackslash}p{6.859cm}}{Best scGPT probe: feature, layer} & \multicolumn{2}{>{\scriptsize\raggedright\arraybackslash}p{3.248cm}}{bilinear, $L_{2}$} & \multicolumn{2}{>{\scriptsize\raggedright\arraybackslash}p{3.248cm}}{bilinear, $L_{8}$} & \multicolumn{2}{>{\scriptsize\raggedright\arraybackslash}p{3.248cm}@{}}{$[x_j$, bilinear$]$, $L_{0}$}\\
\multicolumn{4}{@{}>{\scriptsize\raggedright\arraybackslash}p{6.859cm}}{\quad AUROC [95\% CI, pair ratio]; upper end of the 95\% CI, gene ratio} & \multicolumn{2}{>{\scriptsize\raggedright\arraybackslash}p{3.248cm}}{0.597 [0.573,\,0.621]; 0.638} & \multicolumn{2}{>{\scriptsize\raggedright\arraybackslash}p{3.248cm}}{0.591 [0.557,\,0.624]; 0.647} & \multicolumn{2}{>{\scriptsize\raggedright\arraybackslash}p{3.248cm}@{}}{0.589 [0.555,\,0.624]; 0.647}\\
\multicolumn{4}{@{}>{\scriptsize\raggedright\arraybackslash}p{6.859cm}}{Target-only probe: AUROC, range over layers (factor-only: 0.500 by construction)} & \multicolumn{2}{>{\scriptsize\raggedright\arraybackslash}p{3.248cm}}{0.540--0.554} & \multicolumn{2}{>{\scriptsize\raggedright\arraybackslash}p{3.248cm}}{0.558--0.573} & \multicolumn{2}{>{\scriptsize\raggedright\arraybackslash}p{3.248cm}@{}}{0.553--0.573}\\
\multicolumn{4}{@{}>{\scriptsize\raggedright\arraybackslash}p{6.859cm}}{Concatenation (additive) $-$ target-only: range} & \multicolumn{2}{>{\scriptsize\raggedright\arraybackslash}p{3.248cm}}{$-$0.001 to $+$0.004} & \multicolumn{2}{>{\scriptsize\raggedright\arraybackslash}p{3.248cm}}{$-$0.001 to $+$0.003} & \multicolumn{2}{>{\scriptsize\raggedright\arraybackslash}p{3.248cm}@{}}{$-$0.004 to $+$0.004}\\
\multicolumn{4}{@{}>{\scriptsize\raggedright\arraybackslash}p{6.859cm}}{Hadamard $-$ target-only: range; layers with CI $<0$ (pair; gene ratio)} & \multicolumn{2}{>{\scriptsize\raggedright\arraybackslash}p{3.248cm}}{$-$0.009 to $+$0.002; 0; 0} & \multicolumn{2}{>{\scriptsize\raggedright\arraybackslash}p{3.248cm}}{$-$0.036 to $-$0.016; 3; 0} & \multicolumn{2}{>{\scriptsize\raggedright\arraybackslash}p{3.248cm}@{}}{$-$0.041 to $-$0.017; 0; 0}\\
\multicolumn{4}{@{}>{\scriptsize\raggedright\arraybackslash}p{6.859cm}}{$|x_i - x_j|$ $-$ target-only: range; layers with CI $<0$ (pair; gene ratio)} & \multicolumn{2}{>{\scriptsize\raggedright\arraybackslash}p{3.248cm}}{$-$0.044 to $-$0.019; 7; 0} & \multicolumn{2}{>{\scriptsize\raggedright\arraybackslash}p{3.248cm}}{$-$0.068 to $-$0.034; 12; 8} & \multicolumn{2}{>{\scriptsize\raggedright\arraybackslash}p{3.248cm}@{}}{$-$0.074 to $-$0.033; 9; 5}\\
\multicolumn{4}{@{}>{\scriptsize\raggedright\arraybackslash}p{6.859cm}}{Bilinear $-$ target-only: range over layers} & \multicolumn{2}{>{\scriptsize\raggedright\arraybackslash}p{3.248cm}}{$-$0.004 to $+$0.051} & \multicolumn{2}{>{\scriptsize\raggedright\arraybackslash}p{3.248cm}}{$+$0.002 to $+$0.029} & \multicolumn{2}{>{\scriptsize\raggedright\arraybackslash}p{3.248cm}@{}}{$-$0.031 to $+$0.006}\\
\multicolumn{4}{@{}>{\scriptsize\raggedright\arraybackslash}p{6.859cm}}{\quad largest gain [95\% CI, pair ratio] [gene ratio]} & \multicolumn{2}{>{\scriptsize\raggedright\arraybackslash}p{3.248cm}}{$L_{2}$: $+$0.051 [0.026,\,0.076] [0.009,\,0.093]} & \multicolumn{2}{>{\scriptsize\raggedright\arraybackslash}p{3.248cm}}{$L_{8}$: $+$0.029 [0.008,\,0.051] [$-$0.007, 0.065]} & \multicolumn{2}{>{\scriptsize\raggedright\arraybackslash}p{3.248cm}@{}}{$L_{2}$: $+$0.006 [$-$0.019, 0.031] [$-$0.037, 0.049]}\\
\multicolumn{4}{@{}>{\scriptsize\raggedright\arraybackslash}p{6.859cm}}{\quad layers with CI $>0$: pair; gene ratio} & \multicolumn{2}{>{\scriptsize\raggedright\arraybackslash}p{3.248cm}}{7; 1} & \multicolumn{2}{>{\scriptsize\raggedright\arraybackslash}p{3.248cm}}{2; 0} & \multicolumn{2}{>{\scriptsize\raggedright\arraybackslash}p{3.248cm}@{}}{0; 0}\\
\multicolumn{4}{@{}>{\scriptsize\raggedright\arraybackslash}p{6.859cm}}{$[x_j$, bilinear$]$ $-$ target-only: range; layers with CI $>0$ (pair; gene ratio)} & \multicolumn{2}{>{\scriptsize\raggedright\arraybackslash}p{3.248cm}}{$+$0.001 to $+$0.025; 8; 3} & \multicolumn{2}{>{\scriptsize\raggedright\arraybackslash}p{3.248cm}}{$+$0.009 to $+$0.022; 9; 3} & \multicolumn{2}{>{\scriptsize\raggedright\arraybackslash}p{3.248cm}@{}}{$-$0.004 to $+$0.016; 0; 0}\\
\multicolumn{4}{@{}>{\scriptsize\raggedright\arraybackslash}p{6.859cm}}{Bilinear vs the same probe on rewired labels ($N = 100$): layers $p<0.05$; max-$T<0.05$} & \multicolumn{2}{>{\scriptsize\raggedright\arraybackslash}p{3.248cm}}{12/12; 11/12} & \multicolumn{2}{>{\scriptsize\raggedright\arraybackslash}p{3.248cm}}{8/12; 4/12} & \multicolumn{2}{>{\scriptsize\raggedright\arraybackslash}p{3.248cm}@{}}{4/12; 1/12}\\
\multicolumn{4}{@{}>{\scriptsize\raggedright\arraybackslash}p{6.859cm}}{Concatenation vs rewired labels: layers $p<0.05$; max-$T<0.05$; AUROC true vs rewired (mean)} & \multicolumn{2}{>{\scriptsize\raggedright\arraybackslash}p{3.248cm}}{9/12; 1/12; true 0.560--0.571 vs 0.534--0.542} & \multicolumn{2}{>{\scriptsize\raggedright\arraybackslash}p{3.248cm}}{0/12; 0/12; true 0.547--0.563 vs 0.562--0.574} & \multicolumn{2}{>{\scriptsize\raggedright\arraybackslash}p{3.248cm}@{}}{0/12; 0/12; true 0.547--0.573 vs 0.553--0.563}\\
\multicolumn{4}{@{}>{\scriptsize\raggedright\arraybackslash}p{6.859cm}}{Difference in differences, bilinear $-$ concatenation: true labels, range (rewired-label mean); layers max-$T<0.05$} & \multicolumn{2}{>{\scriptsize\raggedright\arraybackslash}p{3.248cm}}{$-$0.011 to $+$0.046 (rewired $-$0.022 to $-$0.008); 7/12 ($L_{1}$--$L_{7}$)} & \multicolumn{2}{>{\scriptsize\raggedright\arraybackslash}p{3.248cm}}{$+$0.000 to $+$0.028 (rewired $-$0.035 to $-$0.024); 7/12 ($L_{4}$--$L_{10}$)} & \multicolumn{2}{>{\scriptsize\raggedright\arraybackslash}p{3.248cm}@{}}{$-$0.053 to $+$0.008 (rewired $-$0.036 to $-$0.026); 0/12}\\
\multicolumn{4}{@{}>{\scriptsize\raggedright\arraybackslash}p{6.859cm}}{Relational signal (probe $-$ max[rewired, target-only]): layers with CI $>0$, pair; gene ratio} & \multicolumn{2}{>{\scriptsize\raggedright\arraybackslash}p{3.248cm}}{1; 0} & \multicolumn{2}{>{\scriptsize\raggedright\arraybackslash}p{3.248cm}}{0; 0} & \multicolumn{2}{>{\scriptsize\raggedright\arraybackslash}p{3.248cm}@{}}{0; 0}\\
\multicolumn{4}{@{}>{\scriptsize\raggedright\arraybackslash}p{6.859cm}}{Bilinear $-$ signed cosine on the same folds: range; layers with CI $>0$ (pair; gene ratio)} & \multicolumn{2}{>{\scriptsize\raggedright\arraybackslash}p{3.248cm}}{$+$0.020 to $+$0.059; 10; 3} & \multicolumn{2}{>{\scriptsize\raggedright\arraybackslash}p{3.248cm}}{$+$0.035 to $+$0.062; 12; 0} & \multicolumn{2}{>{\scriptsize\raggedright\arraybackslash}p{3.248cm}@{}}{$-$0.009 to $+$0.016; 0; 0}\\
\multicolumn{4}{@{}>{\scriptsize\raggedright\arraybackslash}p{6.859cm}}{scGPT input gene-token embedding (no cell context), bilinear: AUROC; $-$ target-only [pair ratio] [gene ratio]} & \multicolumn{2}{>{\scriptsize\raggedright\arraybackslash}p{3.248cm}}{0.581; $+$0.026 [0.003,\,0.050] [$-$0.013, 0.066]} & \multicolumn{2}{>{\scriptsize\raggedright\arraybackslash}p{3.248cm}}{0.584; $+$0.016 [$-$0.014, 0.046] [$-$0.034, 0.066]} & \multicolumn{2}{>{\scriptsize\raggedright\arraybackslash}p{3.248cm}@{}}{0.577; $+$0.005 [$-$0.022, 0.031] [$-$0.040, 0.049]}\\
\multicolumn{4}{@{}>{\scriptsize\raggedright\arraybackslash}p{6.859cm}}{\quad $p$ vs rewired labels; difference-in-differences $p$} & \multicolumn{2}{>{\scriptsize\raggedright\arraybackslash}p{3.248cm}}{0.0099; 0.0693} & \multicolumn{2}{>{\scriptsize\raggedright\arraybackslash}p{3.248cm}}{0.32; 0.0990} & \multicolumn{2}{>{\scriptsize\raggedright\arraybackslash}p{3.248cm}@{}}{0.0099; 0.21}\\
\multicolumn{4}{@{}>{\scriptsize\raggedright\arraybackslash}p{6.859cm}}{ESM2-15B per-protein embedding (UCE), bilinear: AUROC; $-$ target-only [pair ratio] [gene ratio]} & \multicolumn{2}{>{\scriptsize\raggedright\arraybackslash}p{3.248cm}}{0.567; $+$0.005 [$-$0.021, 0.031] [$-$0.039, 0.049]} & \multicolumn{2}{>{\scriptsize\raggedright\arraybackslash}p{3.248cm}}{0.563; $+$0.012 [$-$0.012, 0.035] [$-$0.027, 0.051]} & \multicolumn{2}{>{\scriptsize\raggedright\arraybackslash}p{3.248cm}@{}}{0.548; $-$0.013 [$-$0.040, 0.014] [$-$0.058, 0.032]}\\
\multicolumn{4}{@{}>{\scriptsize\raggedright\arraybackslash}p{6.859cm}}{\quad $p$ vs rewired labels; difference-in-differences $p$} & \multicolumn{2}{>{\scriptsize\raggedright\arraybackslash}p{3.248cm}}{0.0099; 0.0495} & \multicolumn{2}{>{\scriptsize\raggedright\arraybackslash}p{3.248cm}}{0.0198; 0.0198} & \multicolumn{2}{>{\scriptsize\raggedright\arraybackslash}p{3.248cm}@{}}{0.12; 0.0792}\\
\multicolumn{4}{@{}>{\scriptsize\raggedright\arraybackslash}p{6.859cm}}{scGPT bilinear $-$ input-embedding bilinear (paired): layers with CI $>0$; $<0$} & \multicolumn{2}{>{\scriptsize\raggedright\arraybackslash}p{3.248cm}}{0; 4} & \multicolumn{2}{>{\scriptsize\raggedright\arraybackslash}p{3.248cm}}{0; 0} & \multicolumn{2}{>{\scriptsize\raggedright\arraybackslash}p{3.248cm}@{}}{0; 5}\\
\multicolumn{4}{@{}>{\scriptsize\raggedright\arraybackslash}p{6.859cm}}{scGPT gain over its target-only control $-$ the same gain of the input embedding; of ESM2-15B: layers with CI $>0$} & \multicolumn{2}{>{\scriptsize\raggedright\arraybackslash}p{3.248cm}}{1; 1} & \multicolumn{2}{>{\scriptsize\raggedright\arraybackslash}p{3.248cm}}{0; 0} & \multicolumn{2}{>{\scriptsize\raggedright\arraybackslash}p{3.248cm}@{}}{0; 0}\\
\multicolumn{4}{@{}>{\scriptsize\raggedright\arraybackslash}p{6.859cm}}{Gradient boosting and MLP (fixed and tuned; $L_0$, $L_5$, $L_{11}$; concatenation, Hadamard): AUROC range; rows above; below own target-only} & \multicolumn{2}{>{\scriptsize\raggedright\arraybackslash}p{3.248cm}}{0.496--0.565; 0; 2 of 24} & \multicolumn{2}{>{\scriptsize\raggedright\arraybackslash}p{3.248cm}}{0.491--0.555; 0; 4 of 24} & \multicolumn{2}{>{\scriptsize\raggedright\arraybackslash}p{3.248cm}@{}}{0.493--0.548; 0; 0 of 24}\\
\multicolumn{4}{@{}>{\scriptsize\raggedright\arraybackslash}p{6.859cm}}{\quad held-out factors (targets seen): rows above own target-only} & \multicolumn{2}{>{\scriptsize\raggedright\arraybackslash}p{3.248cm}}{1 of 24} & \multicolumn{2}{>{\scriptsize\raggedright\arraybackslash}p{3.248cm}}{3 of 24} & \multicolumn{2}{>{\scriptsize\raggedright\arraybackslash}p{3.248cm}@{}}{1 of 24}\\
\midrule
\multicolumn{10}{@{}p{16.923cm}@{}}{\footnotesize\textbf{G. Direction: documented pairs $i \to j$ of two factors whose reverse $j \to i$ is in no regulatory resource}}\\*
\multicolumn{10}{@{}p{16.923cm}@{}}{\scriptsize\textit{GENIE3 and GRNBoost2 importances, seed ensemble: fraction of untied pairs with importance($i\to j$) $>$ importance($j\to i$)}}\\*
Extraction & Method & Pairs; untied T1; T2 & T1 fraction & T1 sign-test $p$ & T1 regulator-cluster CI & T2 fraction & T2 sign-test $p$ & T2 regulator-cluster CI & T1 per-seed range\\*
\noalign{\nobreak\vskip 1.5pt\hrule height \cmidrulewidth\nobreak\vskip 2pt\nobreak}
lung & GENIE3 & 855; 834; 845 & 0.511 & 0.56 & [0.464,\,0.559] & 0.509 & 0.63 & [0.461,\,0.557] & 0.480--0.513\\
 & GRNBoost2 & 855; 762; 723 & 0.528 & 0.14 & [0.486,\,0.566] & 0.523 & 0.23 & [0.473,\,0.571] & 0.488--0.510\\
lung (ext.) & GENIE3 & 886; 867; 882 & 0.525 & 0.15 & [0.479,\,0.567] & 0.505 & 0.81 & [0.453,\,0.559] & 0.491--0.521\\
 & GRNBoost2 & 886; 710; 705 & 0.501 & 0.97 & [0.462,\,0.538] & 0.487 & 0.50 & [0.438,\,0.534] & 0.469--0.533\\
kidney~B & GENIE3 & 669; 661; 668 & 0.523 & 0.24 & [0.470,\,0.574] & 0.551 & 0.0095 & [0.492,\,0.609] & 0.530--0.551\\
 & GRNBoost2 & 669; 592; 586 & 0.502 & 0.97 & [0.462,\,0.541] & 0.505 & 0.84 & [0.458,\,0.550] & 0.457--0.548\\
immune & GENIE3 & 62; 60; 60 & 0.517 & 0.90 & [0.387,\,0.633] & 0.467 & 0.70 & [0.347,\,0.607] & 0.448--0.542\\
 & GRNBoost2 & 62; 57; 58 & 0.544 & 0.60 & [0.407,\,0.673] & 0.483 & 0.90 & [0.356,\,0.632] & 0.463--0.589\\
kidney~A & GENIE3 & 21; 21; 21 & 0.429 & 0.66 & [0.222,\,0.654] & 0.524 & 1.00 & [0.263,\,0.789] & 0.381--0.524\\
 & GRNBoost2 & 21; 21; 21 & 0.571 & 0.66 & [0.350,\,0.800] & 0.619 & 0.38 & [0.364,\,0.857] & 0.421--0.600\\
\multicolumn{10}{@{}>{\scriptsize\raggedright\arraybackslash}p{16.923cm}@{}}{Sign tests with $p<0.05$: 1 of 20 (none below the Bonferroni threshold 0.05/20 = 0.0025). $|r|$ and the cosine are symmetric and cannot score direction.}\\
\multicolumn{10}{@{}p{16.923cm}@{}}{\scriptsize\textit{Logistic probe on the ordered concatenation, pair-disjoint folds: fraction of pairs scored forward $>$ reverse}}\\*
Extraction & Pairs; genes & Fraction, range $L_0$--$L_{11}$ & Layers, binomial CI excl.\ 0.5 & Layers, gene-cluster CI excl.\ 0.5 & \multicolumn{2}{>{\scriptsize\raggedright\arraybackslash}p{3.248cm}}{Largest fraction [gene-cluster CI]} & Strict subset: pairs; range & Out-degree $i > j$ & In-degree $j > i$\\*
\noalign{\nobreak\vskip 1.5pt\hrule height \cmidrulewidth\nobreak\vskip 2pt\nobreak}
lung & 435; 241 & 0.494--0.543 & 0/12 & 0/12 & \multicolumn{2}{>{\scriptsize\raggedright\arraybackslash}p{3.248cm}}{$L_{6}$ 0.543 [0.460,\,0.626]} & 343; 0.499--0.563 & 0.559 & 0.547\\
lung (ext.) & 453; 251 & 0.536--0.594 & 11/12 & 2/12 ($L_{10}$--$L_{11}$) & \multicolumn{2}{>{\scriptsize\raggedright\arraybackslash}p{3.248cm}}{$L_{11}$ 0.594 [0.506,\,0.677]} & 357; 0.529--0.599 & 0.532 & 0.586\\
kidney~B & 325; 197 & 0.563--0.628 & 12/12 & 1/12 ($L_{7}$) & \multicolumn{2}{>{\scriptsize\raggedright\arraybackslash}p{3.248cm}}{$L_{7}$ 0.628 [0.512,\,0.740]} & 260; 0.565--0.631 & 0.550 & 0.549\\
\multicolumn{10}{@{}>{\scriptsize\raggedright\arraybackslash}p{16.923cm}@{}}{Gene-cluster intervals excluding 0.5: 3 of 36 extraction $\times$ layer combinations.}\\
\end{longtable}
\endgroup

\setcounter{table}{7}
\begingroup
\footnotesize
\setlength{\tabcolsep}{3pt}
\renewcommand{\arraystretch}{1.05}
\begin{longtable}{@{}>{\scriptsize\raggedright\arraybackslash}p{3.1cm}*{6}{>{\scriptsize\raggedright\arraybackslash}p{2.1cm}}@{}}
\caption{\textbf{Rank reduction (H4): truncation to the leading spectral directions against the full representation.} $m$, layer-mean task metric (mean over the 12 layers) of the representation restricted to the first $k$ directions $U_k S_k$ of the observed-gene SVD at each layer ($m(512)$, the full representation). Tasks: interaction, AUROC of the cosine of STRING physical pairs (combined score $\geq 700$) against one degree-preserving rewired set of non-pairs; TF identity, the H2 statistic (detection-stratified LOO-cosine AUROC, Lambert TFs against all other observed genes); regulatory edges, the H3 statistic (within-TF stratified AUROC of the TF--target cosine, TRRUST pairs against matched candidates). \textbf{A}: $k^{*}$ maximises $m(k)$ over $k \in \{1, 2, 3, 5, 7, 10, 16, 32, 64, 128, 256, 512\}$ in the selection draw (ties to the smallest $k$); selection gap $= m(k^{*}) - m(512)$ there. \textbf{B}: $\Delta$ in the evaluation extractions with paired 95\% intervals ($B = 2{,}000$; the same resamples for $k^{*}$ and 512): interaction, vertex subsampling ($m = n/2$) with delete-$d$ scaling; TF identity, gene bootstrap; edges, two-way TF $\times$ target bootstrap. For TF identity the selection draw chose $k^{*} = 512$, so $\Delta = 0$ by construction. Last column: number of the 12 per-layer intervals entirely below / above 0. \textbf{C}--\textbf{E}: descriptive, not pre-registered; kidney~A and immune (seed 43, the selection draw) are shown for completeness; the last column has no interval. \textbf{D}: removing the leading directions instead of keeping them; intervals as in B (interaction delete-$d$). \textbf{E}: screen bands $\SV{2}$--$\SV{4}$, $\SV{5}$--$\SV{7}$, $\SV{2}$--$\SV{7}$ on the observed-gene basis against the first 3 or 6 directions (the screen used $\SV{5}$--$\SV{7}$ of the observed genes in its later kidney~A iterations; Table~S14); $^{b}$$\SV{5}$--$\SV{7}$ shifted down by one index, as when it is read in the reference decomposition of the full kidney~A matrix, zero rows included (there $\SV{2}$--$\SV{4}$ and $\SV{2}$--$\SV{7}$ become the first 3 and 6 directions, a difference of 0 by construction), shown in every extraction. \textbf{F}: one vector per gene: input gene-token embeddings (scGPT, Geneformer, tGPT); gene outputs of one forward pass on a neutral input (AIDO.Cell-100M, scFoundation); ESM2-15B, the 5{,}120-dimensional per-protein embedding (UCE file), and residue-mean-pooled ESM2-650M; 32 sequence descriptors; and expression across 154 Human Protein Atlas cell types (128 principal components). Each is mean-centred on the 5{,}226 lung genes present in all eleven vocabularies and scored on identical genes, pairs and labels; $d$, dimension; best $k$ over the same grid, gain $=$ best minus full (no interval). $^{c}$Reference row: the scGPT lung residual stream on the same genes, layer mean. No truncation exceeds the full embedding by more than 0.022 for any static representation and task. Per-layer values, all $k$, the descriptive any-$k$ permutation test and, for interaction, the intervals with the pre-registered scaling factor $\sqrt{m/n}$, which omits the finite-population term and so gives half-widths smaller by a factor $\sqrt{2}$ (columns \texttt{ci\_lo}/\texttt{ci\_hi}, beside the reported \texttt{ci\_lo\_ms}/\texttt{ci\_hi\_ms}), are in the CSV. Full values: \texttt{Table\_S8\_rank.csv}.}\label{tab:S8}\\
\toprule
\endfirsthead
\multicolumn{7}{@{}l}{\footnotesize Table~\thetable\ (continued)}\\
\toprule
\endhead
\midrule
\multicolumn{7}{r@{}}{\footnotesize continued on next page}\\
\endfoot
\bottomrule
\endlastfoot
\multicolumn{7}{@{}p{16.4cm}@{}}{\footnotesize\textbf{A. Truncation $k^{*}$ chosen in the selection draw (immune, seed 43), before any evaluation extraction}}\\
Task & $k^{*}$ & $m(k^{*})$ & $m(512)$ & Selection gap & $m(3) - m(512)$ & $m(10) - m(512)$\\
\cmidrule(lr){1-7}
Interaction & 64 & 0.7047 & 0.7021 & $+$0.0026 & $-$0.099 & $-$0.018\\
TF identity & 512 & 0.6088 & 0.6088 & $+$0.0000 & $-$0.093 & $-$0.055\\
Regulatory edges & 128 & 0.6199 & 0.6198 & $+$0.0001 & $-$0.087 & $-$0.037\\
\midrule
\multicolumn{7}{@{}p{16.4cm}@{}}{\footnotesize\textbf{B. Pre-registered estimand $\Delta = m(k^{*}) - m(512)$, layer mean, in the four evaluation extractions}}\\
Task, extraction & $m(k^{*})$ & $m(512)$ & $\Delta$ & \multicolumn{2}{>{\scriptsize}l}{95\% CI} & Layers with CI $<0$ / $>0$\\
\cmidrule(lr){1-7}
\multicolumn{7}{@{}l}{\scriptsize\textit{Interaction, $k^{*} = 64$}}\\
\quad lung & 0.6539 & 0.6549 & $-$0.0011 & \multicolumn{2}{>{\scriptsize}l}{[$-$0.0024, 0.0003]} & 6 / 0\\
\quad lung (ext.) & 0.6573 & 0.6579 & $-$0.0006 & \multicolumn{2}{>{\scriptsize}l}{[$-$0.0019, 0.0007]} & 0 / 0\\
\quad kidney~B & 0.6474 & 0.6486 & $-$0.0013 & \multicolumn{2}{>{\scriptsize}l}{[$-$0.0026, $-$0.0000]} & 7 / 0\\
\quad immune & 0.7020 & 0.6998 & $+$0.0022 & \multicolumn{2}{>{\scriptsize}l}{[$-$0.0018, 0.0064]} & 0 / 0\\
\multicolumn{7}{@{}l}{\scriptsize\textit{TF identity, $k^{*} = 512$}}\\
\quad lung & 0.6542 & 0.6542 & 0 (identical) & \multicolumn{2}{>{\scriptsize}l}{--} & --\\
\quad lung (ext.) & 0.6436 & 0.6436 & 0 (identical) & \multicolumn{2}{>{\scriptsize}l}{--} & --\\
\quad kidney~B & 0.6343 & 0.6343 & 0 (identical) & \multicolumn{2}{>{\scriptsize}l}{--} & --\\
\quad immune & 0.6553 & 0.6553 & 0 (identical) & \multicolumn{2}{>{\scriptsize}l}{--} & --\\
\multicolumn{7}{@{}l}{\scriptsize\textit{Regulatory edges, $k^{*} = 128$}}\\
\quad lung & 0.5304 & 0.5300 & $+$0.0004 & \multicolumn{2}{>{\scriptsize}l}{[$-$0.0022, 0.0026]} & 0 / 0\\
\quad lung (ext.) & 0.5303 & 0.5300 & $+$0.0003 & \multicolumn{2}{>{\scriptsize}l}{[$-$0.0021, 0.0026]} & 0 / 0\\
\quad kidney~B & 0.5411 & 0.5413 & $-$0.0002 & \multicolumn{2}{>{\scriptsize}l}{[$-$0.0027, 0.0023]} & 0 / 0\\
\quad immune & 0.6092 & 0.6096 & $-$0.0004 & \multicolumn{2}{>{\scriptsize}l}{[$-$0.0070, 0.0061]} & 0 / 0\\
\midrule
\multicolumn{7}{@{}p{16.4cm}@{}}{\footnotesize\textbf{C. The leading $k$ directions only: layer-mean value at $k = 3$ and $k = 10$ (difference from $m(512)$ in parentheses)}}\\
Task, $k$ & lung & lung (ext.) & kidney~B & immune & kidney~A & immune (seed 43)\\
\cmidrule(lr){1-7}
Interaction, $k = 512$ (full) & 0.655 & 0.658 & 0.649 & 0.700 & 0.767 & 0.702\\
Interaction, $k = 3$ & 0.614 ($-$0.041) & 0.626 ($-$0.032) & 0.610 ($-$0.038) & 0.609 ($-$0.091) & 0.718 ($-$0.049) & 0.604 ($-$0.099)\\
Interaction, $k = 10$ & 0.641 ($-$0.014) & 0.647 ($-$0.011) & 0.634 ($-$0.014) & 0.682 ($-$0.018) & 0.760 ($-$0.008) & 0.684 ($-$0.018)\\
TF identity, $k = 512$ (full) & 0.654 & 0.644 & 0.634 & 0.655 & 0.742 & 0.609\\
TF identity, $k = 3$ & 0.526 ($-$0.128) & 0.556 ($-$0.087) & 0.547 ($-$0.088) & 0.513 ($-$0.142) & 0.620 ($-$0.122) & 0.516 ($-$0.093)\\
TF identity, $k = 10$ & 0.612 ($-$0.042) & 0.625 ($-$0.018) & 0.601 ($-$0.034) & 0.571 ($-$0.084) & 0.700 ($-$0.041) & 0.554 ($-$0.055)\\
Regulatory edges, $k = 512$ (full) & 0.530 & 0.530 & 0.541 & 0.610 & 0.632 & 0.620\\
Regulatory edges, $k = 3$ & 0.496 ($-$0.034) & 0.492 ($-$0.038) & 0.504 ($-$0.037) & 0.527 ($-$0.083) & 0.568 ($-$0.063) & 0.533 ($-$0.087)\\
Regulatory edges, $k = 10$ & 0.516 ($-$0.014) & 0.519 ($-$0.011) & 0.527 ($-$0.014) & 0.566 ($-$0.044) & 0.603 ($-$0.028) & 0.583 ($-$0.037)\\
\midrule
\multicolumn{7}{@{}p{16.4cm}@{}}{\footnotesize\textbf{D. Removing the $j$ leading directions (keeping $\SV{j+1}$--$\SV{512}$): layer-mean difference from $m(512)$ [95\% CI]}}\\
Task, directions removed & lung & lung (ext.) & kidney~B & immune & kidney~A & immune (seed 43)\\
\cmidrule(lr){1-7}
Interaction, $j = 1$ & $-$0.007\newline [$-$0.012, $-$0.002] & $+$0.002\newline [$-$0.004, 0.009] & $+$0.001\newline [$-$0.005, 0.007] & $+$0.027\newline [0.011, 0.042] & $+$0.031\newline [0.015, 0.046] & $+$0.022\\
Interaction, $j = 2$ & $-$0.010\newline [$-$0.020, 0.000] & $-$0.023\newline [$-$0.036, $-$0.008] & $-$0.006\newline [$-$0.018, 0.005] & $+$0.018\newline [$-$0.004, 0.041] & $+$0.012\newline [$-$0.015, 0.037] & $+$0.012\\
Interaction, $j = 3$ & $-$0.024\newline [$-$0.037, $-$0.010] & $-$0.030\newline [$-$0.044, $-$0.014] & $-$0.019\newline [$-$0.033, $-$0.005] & $+$0.015\newline [$-$0.008, 0.038] & $+$0.004\newline [$-$0.026, 0.032] & $+$0.013\\
Interaction, $j = 5$ & $-$0.027\newline [$-$0.042, $-$0.013] & $-$0.036\newline [$-$0.051, $-$0.021] & $-$0.017\newline [$-$0.033, $-$0.002] & $-$0.000\newline [$-$0.027, 0.027] & $-$0.024\newline [$-$0.061, 0.012] & $+$0.004\\
Interaction, $j = 10$ & $-$0.041\newline [$-$0.056, $-$0.026] & $-$0.049\newline [$-$0.065, $-$0.033] & $-$0.033\newline [$-$0.048, $-$0.016] & $-$0.045\newline [$-$0.081, $-$0.011] & $-$0.068\newline [$-$0.110, $-$0.028] & $-$0.041\\
TF identity, $j = 1$ & $+$0.036\newline [0.027, 0.043] & $+$0.028\newline [0.019, 0.037] & $+$0.025\newline [0.015, 0.035] & $+$0.043\newline [0.006, 0.070] & $+$0.006\newline [$-$0.047, 0.027] & $+$0.034\\
TF identity, $j = 2$ & $+$0.060\newline [0.047, 0.070] & $+$0.098\newline [0.079, 0.108] & $+$0.070\newline [0.050, 0.079] & $+$0.082\newline [0.024, 0.112] & $-$0.016\newline [$-$0.067, 0.028] & $+$0.092\\
TF identity, $j = 3$ & $+$0.091\newline [0.074, 0.102] & $+$0.099\newline [0.080, 0.112] & $+$0.091\newline [0.065, 0.104] & $+$0.086\newline [0.021, 0.110] & $-$0.002\newline [$-$0.055, 0.044] & $+$0.101\\
TF identity, $j = 5$ & $+$0.099\newline [0.080, 0.111] & $+$0.107\newline [0.087, 0.122] & $+$0.108\newline [0.084, 0.122] & $+$0.092\newline [0.035, 0.122] & $-$0.007\newline [$-$0.090, 0.044] & $+$0.117\\
TF identity, $j = 10$ & $+$0.089\newline [0.063, 0.103] & $+$0.092\newline [0.060, 0.105] & $+$0.108\newline [0.077, 0.123] & $+$0.099\newline [0.037, 0.146] & $-$0.001\newline [$-$0.100, 0.059] & $+$0.130\\
Regulatory edges, $j = 1$ & $+$0.009\newline [$-$0.007, 0.021] & $+$0.011\newline [$-$0.005, 0.026] & $+$0.009\newline [$-$0.005, 0.024] & $+$0.005\newline [$-$0.019, 0.031] & $+$0.003\newline [$-$0.046, 0.046] & $-$0.000\\
Regulatory edges, $j = 2$ & $+$0.019\newline [$-$0.007, 0.045] & $+$0.031\newline [$-$0.006, 0.067] & $+$0.019\newline [$-$0.007, 0.043] & $+$0.020\newline [$-$0.025, 0.061] & $-$0.005\newline [$-$0.095, 0.065] & $+$0.010\\
Regulatory edges, $j = 3$ & $+$0.033\newline [$-$0.009, 0.071] & $+$0.037\newline [$-$0.007, 0.075] & $+$0.023\newline [$-$0.014, 0.060] & $+$0.025\newline [$-$0.017, 0.064] & $-$0.009\newline [$-$0.109, 0.072] & $+$0.014\\
Regulatory edges, $j = 5$ & $+$0.041\newline [$-$0.006, 0.081] & $+$0.039\newline [$-$0.007, 0.078] & $+$0.023\newline [$-$0.018, 0.067] & $+$0.026\newline [$-$0.021, 0.070] & $-$0.012\newline [$-$0.109, 0.076] & $+$0.018\\
Regulatory edges, $j = 10$ & $+$0.017\newline [$-$0.034, 0.060] & $+$0.014\newline [$-$0.033, 0.056] & $+$0.003\newline [$-$0.040, 0.047] & $+$0.022\newline [$-$0.058, 0.086] & $-$0.025\newline [$-$0.116, 0.067] & $+$0.005\\
\midrule
\multicolumn{7}{@{}p{16.4cm}@{}}{\footnotesize\textbf{E. The screen's bands against the leading directions of the same width: layer-mean difference [95\% CI]}}\\
Task, band vs partner & lung & lung (ext.) & kidney~B & immune & kidney~A & immune (seed 43)\\
\cmidrule(lr){1-7}
Interaction, $\SV{2}$--$\SV{4}$ vs top 3 & $-$0.018\newline [$-$0.024, $-$0.012] & $-$0.014\newline [$-$0.022, $-$0.006] & $-$0.009\newline [$-$0.017, $-$0.001] & $+$0.036\newline [0.009, 0.063] & $+$0.016\newline [$-$0.005, 0.038] & $+$0.032\\
Interaction, $\SV{5}$--$\SV{7}$ vs top 3 & $-$0.036\newline [$-$0.053, $-$0.019] & $-$0.049\newline [$-$0.067, $-$0.030] & $-$0.039\newline [$-$0.055, $-$0.022] & $+$0.037\newline [$-$0.002, 0.075] & $-$0.021\newline [$-$0.061, 0.016] & $+$0.042\\
Interaction, $\SV{2}$--$\SV{7}$ vs top 6 & $-$0.013\newline [$-$0.018, $-$0.008] & $-$0.007\newline [$-$0.014, $-$0.000] & $-$0.005\newline [$-$0.011, 0.001] & $+$0.031\newline [0.010, 0.051] & $+$0.023\newline [0.008, 0.038] & $+$0.033\\
Interaction, $\SV{4}$--$\SV{6}$$^{b}$ vs top 3 & $-$0.036\newline [$-$0.053, $-$0.018] & $-$0.051\newline [$-$0.069, $-$0.033] & $-$0.039\newline [$-$0.056, $-$0.021] & $+$0.026\newline [$-$0.009, 0.063] & $-$0.019\newline [$-$0.057, 0.019] & $+$0.027\\
TF identity, $\SV{2}$--$\SV{4}$ vs top 3 & $+$0.023\newline [0.008, 0.036] & $+$0.003\newline [$-$0.012, 0.020] & $+$0.004\newline [$-$0.010, 0.019] & $+$0.039\newline [$-$0.015, 0.103] & $-$0.060\newline [$-$0.124, $-$0.024] & $+$0.004\\
TF identity, $\SV{5}$--$\SV{7}$ vs top 3 & $+$0.045\newline [0.012, 0.083] & $+$0.063\newline [0.035, 0.103] & $+$0.044\newline [0.012, 0.083] & $-$0.037\newline [$-$0.094, 0.066] & $+$0.005\newline [$-$0.057, 0.078] & $+$0.026\\
TF identity, $\SV{2}$--$\SV{7}$ vs top 6 & $+$0.015\newline [0.006, 0.022] & $+$0.026\newline [0.014, 0.043] & $+$0.020\newline [0.005, 0.036] & $+$0.019\newline [$-$0.026, 0.065] & $-$0.014\newline [$-$0.060, 0.013] & $+$0.037\\
TF identity, $\SV{4}$--$\SV{6}$$^{b}$ vs top 3 & $+$0.051\newline [0.017, 0.090] & $+$0.008\newline [$-$0.024, 0.046] & $-$0.001\newline [$-$0.036, 0.041] & $-$0.012\newline [$-$0.073, 0.079] & $+$0.008\newline [$-$0.064, 0.082] & $-$0.023\\
Regulatory edges, $\SV{2}$--$\SV{4}$ vs top 3 & $+$0.000\newline [$-$0.020, 0.015] & $+$0.008\newline [$-$0.020, 0.040] & $+$0.006\newline [$-$0.016, 0.027] & $-$0.015\newline [$-$0.064, 0.045] & $-$0.030\newline [$-$0.099, 0.034] & $-$0.016\\
Regulatory edges, $\SV{5}$--$\SV{7}$ vs top 3 & $+$0.045\newline [$-$0.009, 0.095] & $+$0.050\newline [$-$0.003, 0.104] & $+$0.036\newline [$-$0.008, 0.078] & $+$0.016\newline [$-$0.062, 0.080] & $-$0.036\newline [$-$0.186, 0.095] & $+$0.006\\
Regulatory edges, $\SV{2}$--$\SV{7}$ vs top 6 & $+$0.007\newline [$-$0.011, 0.021] & $+$0.011\newline [$-$0.008, 0.029] & $+$0.013\newline [$-$0.004, 0.029] & $-$0.008\newline [$-$0.049, 0.030] & $-$0.015\newline [$-$0.068, 0.033] & $-$0.012\\
Regulatory edges, $\SV{4}$--$\SV{6}$$^{b}$ vs top 3 & $+$0.036\newline [$-$0.015, 0.084] & $+$0.042\newline [$-$0.012, 0.091] & $+$0.029\newline [$-$0.017, 0.082] & $+$0.009\newline [$-$0.062, 0.074] & $-$0.050\newline [$-$0.180, 0.060] & $+$0.000\\
\midrule
\multicolumn{7}{@{}p{16.4cm}@{}}{\footnotesize\textbf{F. Eleven static gene representations on the common lung gene set (5{,}226 genes): full embedding and best truncation}}\\
Representation ($d$) & Interaction: full & Interaction: best $k$ (gain) & TF identity: full & TF identity: best $k$ (gain) & Edges: full & Edges: best $k$ (gain)\\
\cmidrule(lr){1-7}
scGPT input embedding (512) & 0.661 & 64 ($+$0.001) & 0.751 & 512 ($+$0.000) & 0.566 & 128 ($+$0.004)\\
Geneformer V2-316M (1{,}152) & 0.650 & 5 ($+$0.004) & 0.821 & 1152 ($+$0.000) & 0.561 & 16 ($+$0.004)\\
Geneformer V1-10M (256) & 0.583 & 256 ($+$0.000) & 0.725 & 256 ($+$0.000) & 0.518 & 256 ($+$0.000)\\
Geneformer V2-104M (768) & 0.651 & 7 ($+$0.017) & 0.825 & 512 ($+$0.001) & 0.562 & 64 ($+$0.010)\\
AIDO.Cell-100M (640) & 0.605 & 256 ($+$0.000) & 0.627 & 64 ($+$0.022) & 0.546 & 640 ($+$0.000)\\
scFoundation (512) & 0.541 & 32 ($+$0.000) & 0.454 & 5 ($+$0.003) & 0.502 & 32 ($+$0.002)\\
tGPT (1{,}024) & 0.574 & 1024 ($+$0.000) & 0.612 & 512 ($+$0.001) & 0.519 & 512 ($+$0.001)\\
ESM2-15B (5{,}120) & 0.639 & 512 ($+$0.000) & 0.880 & 32 ($+$0.002) & 0.498 & 1 ($+$0.009)\\
ESM2-650M (1{,}280) & 0.615 & 512 ($+$0.000) & 0.846 & 64 ($+$0.001) & 0.497 & 1 ($+$0.007)\\
Sequence descriptors (32) & 0.667 & 32 ($+$0.000) & 0.867 & 32 ($+$0.000) & 0.499 & 5 ($+$0.001)\\
HPA cell types (128) & 0.595 & 16 ($+$0.002) & 0.573 & 128 ($+$0.000) & 0.515 & 16 ($+$0.001)\\
\cmidrule(lr){1-7}
scGPT lung, layer mean$^{c}$ (512) & 0.649 & 512 ($+$0.000) & 0.657 & 128 ($+$0.000) & 0.530 & 128 ($+$0.001)\\
\end{longtable}
\endgroup

\setcounter{table}{8}
\begingroup
\footnotesize
\setlength{\tabcolsep}{3pt}
\renewcommand{\arraystretch}{1.05}
\begin{longtable}{@{}>{\scriptsize\raggedright\arraybackslash}p{3.1cm}*{6}{>{\scriptsize\raggedright\arraybackslash}p{2.1cm}}@{}}
\caption{\textbf{Differences between extractions: variance components, transfer of a TF read-out, lineage TFs, detection, and the two kidney extractions.} Statistics (all AUROCs, V-centred embeddings, layer mean over the 12 layers unless stated): S1, cosine of STRING physical pairs against one degree-preserving rewired non-pair set; S2, the H2 statistic (detection-stratified LOO-cosine AUROC, Lambert TFs against all other observed genes); S2b, TRRUST TFs against target-only genes, in the full representation or in the band $\SV{2}$--$\SV{7}$; S3, the H3 statistic (within-TF stratified AUROC of the TF--target cosine). \textbf{A}: moment estimate of the between-extraction SD $\tau$ (AUROC units; 0 when the estimate of $\tau^2$ is negative) among the $k$ extractions of each group: cell draws (immune, seed 43, seed 44; kidney~B and the whole-human checkpoint run), donors TSP14/21/25, B/CD4 T/CD8 T cells, kidney~A vs B, the three checkpoints on the same 192 cells, and the four tissue-level extractions (immune, kidney~B, lung, lung (ext.)). Sampling covariances come from one Poisson gene bootstrap shared by all extractions ($B = 1{,}000$; S3 two-way TF $\times$ target; S1 node bootstrap, whose SE is within 5\% of the delete-$d$ vertex-subsampling SE). A tissue effect is claimed only if $D$, the tissue component minus the pooled within-tissue floor (cell draws, donors, gene axis), has a 95\% lower bound $>0$ under a parametric random-effects bootstrap ($R = 10{,}000$); the lower bound exceeds 0 neither for the layer mean nor at any single layer (0 of 12 layers for every statistic). Rows marked $\times 10^{3}$ are in units of $10^{-3}$ AUROC$^2$. Q-profile bounds treat the extractions as random draws; with 2--4 extractions per group they are wide. \textbf{B}: L2-regularised logistic regression (balanced classes; $C$ chosen by 5-fold CV from $\{0.01, 0.1, 1, 10\}$) on unit-normalised $V$-centred embeddings, fitted per layer on immune genes; novel genes are not observed in immune, shared genes are; the in-context reference is a 5-fold out-of-fold classifier within the test extraction scored on the same genes; intervals are 95\% gene-bootstrap intervals ($B = 2{,}000$); AUROCs without stratification, with the detection-stratified AUROC (within deciles of the test extraction's $\log n_{\mathrm{obs}}$, as for S2) in parentheses; below the range over layers, the smallest lower bound over layers; the difference, its intervals and, in parentheses, the layers whose interval of the difference lies below 0 are without stratification; per-layer values in the CSV. \textbf{C}: sets of the 30 Lambert TFs most specific to each tissue in raw expression (18 for lung, which must be expressed in both lung studies), selected from raw expression alone, without any embedding statistic; for example immune BACH2, ETS1, BCL11B, IKZF1, PAX5, EBF1, IRF4; kidney PAX2, PAX8, HNF1B, HNF4A, EMX2, SALL1; lung NKX2-1, FOXA2, FOXF1, GATA6, EPAS1 (full lists in the CSV). Score: LOO cosine to the observed set members; at least 5 observed members required; centred Poisson gene-bootstrap intervals ($B = 500$) clipped to $[0, 1]$; observed, members in $V$; random, mean AUROC of 50 random Lambert-TF subsets of the same size; $^{h}$the set's own tissue; $^{d}$interval unreliable (fewer than 10 observed members). Across sets and extractions the AUROC falls with the set's median $n_{\mathrm{obs}}$ (Spearman $\rho = -0.39$, $n = 29$; rows not independent). \textbf{D}: AUROC of $\log n_{\mathrm{obs}}$ alone (no stratification), gene-bootstrap 95\% intervals ($B = 2{,}000$); values above 0.5 mean TFs are observed in more cells. Cells whose input exceeded scGPT's 512-gene cap are truncated. \textbf{E}: $M$, genes observed in both extractions, with labels recomputed on $M$; paired gene bootstrap on $M$ ($B = 2{,}000$), $\Delta = $ kidney~A $-$ kidney~B. Curated-gene comparison: TRRUST TFs / targets of the curated gene list derived from the screen's kidney~A edge list, with the SVD fitted on the full matrix including genes no sampled cell contained (zero rows); the Hanley--McNeil SE is that of a single AUROC at the stated value and sample sizes. Full values: \texttt{Table\_S9\_cross\_extraction.csv}.}\label{tab:S9}\\
\toprule
\endfirsthead
\multicolumn{7}{@{}l}{\footnotesize Table~\thetable\ (continued)}\\
\toprule
\endhead
\midrule
\multicolumn{7}{r@{}}{\footnotesize continued on next page}\\
\endfoot
\bottomrule
\endlastfoot
\multicolumn{7}{@{}p{16.4cm}@{}}{\footnotesize\textbf{A. Variance components of the layer-mean statistic across the 16 extractions: between-extraction SD $\tau$ within each group, and the tissue decision}}\\*
Component & $k$ & S1 interaction & S2 TF identity & S2b TRRUST, full & S2b TRRUST, band & S3 edges\\*
\cmidrule(lr){1-7}
Cell draw, immune & 3 & 0.000 & 0.027 & 0.016 & 0.033 & 0.000\\
Cell draw, kidney & 2 & 0.000 & 0.000 & 0.000 & 0.000 & 0.000\\
Donor (immune) & 3 & 0.000 & 0.019 & 0.000 & 0.078 & 0.000\\
Cell type (immune) & 3 & 0.000 & 0.000 & 0.000 & 0.070 & 0.000\\
Gene axis (kidney~A vs B) & 2 & 0.075 & 0.073 & 0.000 & 0.000 & 0.052\\
Checkpoint (same cells) & 3 & 0.023 & 0.027 & 0.043 & 0.030 & 0.000\\
\cmidrule(lr){1-7}
Tissue & 4 & 0.020 & 0.000 & 0.000 & 0.000 & 0.033\\
Tissue, without lung (ext.) & 3 & 0.025 & 0.000 & 0.000 & 0.000 & 0.037\\
Study (lung vs lung (ext.)) & 2 & 0.000 & 0.005 & 0.000 & 0.012 & 0.000\\
Tissue $\tau$, upper Q-profile 95\% bound &  & 0.086 & 0.034 & 0.029 & 0.050 & 0.139\\
\cmidrule(lr){1-7}
$\tau^2_{\mathrm{tissue}}$ ($\times 10^{3}$) &  & 0.42 & $-$0.15 & $-$0.26 & $-$0.28 & 1.10\\
Pooled within-tissue floor $\tau^2_{\mathrm{within}}$ ($\times 10^{3}$) &  & 0.89 & 1.25 & $-$0.15 & 2.17 & 0.38\\
$D = \tau^2_{\mathrm{tissue}} - \tau^2_{\mathrm{within}}$ ($\times 10^{3}$) [95\% interval] &  & $-$0.47\newline [$-$2.33, 3.71] & $-$1.40\newline [$-$3.15, 2.68] & $-$0.12\newline [$-$1.51, 0.84] & $-$2.45\newline [$-$6.10, 5.00] & $+$0.72\newline [$-$2.69, 3.84]\\
$p$ ($H_0$: tissue $=$ floor); tissue effect claimed &  & 0.73; no & 0.83; no & 0.45; no & 0.79; no & 0.18; no\\
\midrule
\multicolumn{7}{@{}p{16.4cm}@{}}{\footnotesize\textbf{B. A TF classifier trained on immune genes, applied unchanged to the other extractions (layer mean)}}\\*
Test extraction & TFs / non-TFs (novel genes) & Novel genes: AUROC (detection-stratified) & Novel genes: range over layers & In-context classifier, same genes (detection-stratified) & Transfer minus in-context (layers CI $<0$) & Shared genes: AUROC\\*
\cmidrule(lr){1-7}
immune (training; 5-fold out-of-fold) & 113 / 1{,}926 & 0.835 (0.796) & -- & -- & -- & --\\
kidney~B & 438 / 3{,}632 & 0.652 (0.672) & 0.625--0.669\newline lower $\geq$ 0.598 & 0.784 (0.774) & $-$0.131 (12/12) & 0.955\\
lung & 495 / 4{,}297 & 0.674 (0.698) & 0.662--0.684\newline lower $\geq$ 0.636 & 0.786 (0.786) & $-$0.112 (12/12) & 0.966\\
lung (ext.) & 539 / 4{,}292 & 0.672 (0.693) & 0.656--0.680\newline lower $\geq$ 0.631 & 0.798 (0.796) & $-$0.126 (12/12) & 0.959\\
kidney~A & 45 / 1{,}414 & 0.871 (0.848) & 0.793--0.902\newline lower $\geq$ 0.728 & 0.876 (0.844) & $-$0.005 (0/12) & 0.964\\
\midrule
\multicolumn{7}{@{}p{16.4cm}@{}}{\footnotesize\textbf{C. Lineage-specific TF sets: detection-stratified AUROC of set members against non-TF genes (layer mean)}}\\*
Extraction & Immune set: AUROC [95\% CI] & Immune set: observed (random) & Kidney set: AUROC [95\% CI] & Kidney set: observed (random) & Lung set: AUROC [95\% CI] & Lung set: observed (random)\\*
\cmidrule(lr){1-7}
lung & 0.911\newline [0.859, 0.984] & 28 (0.595) & 0.795\newline [0.653, 1.000] & 13 (0.560) & 0.903$^{h}$\newline [0.836, 1.000] & 16 (0.551)\\
lung (ext.) & 0.908\newline [0.860, 0.973] & 27 (0.589) & 0.711\newline [0.545, 0.962] & 10 (0.539) & 0.916$^{h}$\newline [0.858, 1.000] & 16 (0.548)\\
kidney~B & 0.873\newline [0.816, 0.932] & 29 (0.585) & 0.962$^{h}$\newline [0.941, 0.991] & 21 (0.570) & 0.937\newline [0.860, 1.000]$^{d}$ & 8 (0.507)\\
immune & 0.709$^{h}$\newline [0.543, 0.911] & 20 (0.604) & -- & 3 & -- & 0\\
kidney~A & 0.763\newline [0.667, 0.877] & 21 (0.697) & -- & 2 & -- & 3\\
\cmidrule(lr){1-7}
kidney (whole-human ckpt) & 0.890\newline [0.844, 0.950] & 30 (0.570) & 0.961$^{h}$\newline [0.940, 0.998] & 21 (0.565) & 0.885\newline [0.705, 1.000]$^{d}$ & 6 (0.512)\\
kidney (kidney ckpt) & 0.776\newline [0.693, 0.872] & 30 (0.538) & 0.716$^{h}$\newline [0.582, 0.872] & 21 (0.537) & 0.883\newline [0.715, 1.000]$^{d}$ & 6 (0.448)\\
kidney (brain ckpt) & 0.796\newline [0.709, 0.906] & 30 (0.515) & 0.829$^{h}$\newline [0.751, 0.947] & 21 (0.495) & 0.614\newline [0.322, 1.000]$^{d}$ & 6 (0.422)\\
immune, 8 robustness extractions & 0.665--0.755 & 19--20 & 1 of 8 evaluable & 2--5 & -- & 0\\
\midrule
\multicolumn{7}{@{}p{16.4cm}@{}}{\footnotesize\textbf{D. Detection alone: AUROC of $\log n_{\mathrm{obs}}$ for separating TFs, and the extraction regime}}\\*
Extraction & $|V|$; cells (input $>512$ genes) & Median non-zero genes per cell & Median $n_{\mathrm{obs}}$, Lambert TFs / non-TFs & Lambert TFs vs non-TFs: AUROC [95\% CI] & TRRUST TFs vs target-only: AUROC [95\% CI] & Lambert TFs with $n_{\mathrm{obs}} \geq 10$\\*
\cmidrule(lr){1-7}
lung & 5{,}648; 256 (249) & 1{,}668 & 6 / 8 & 0.458\newline [0.434, 0.482] & 0.519\newline [0.492, 0.548] & 38\%\\
lung (ext.) & 5{,}616; 256 (220) & 935 & 9 / 12 & 0.454\newline [0.432, 0.476] & 0.534\newline [0.507, 0.560] & 48\%\\
kidney~B & 4{,}780; 192 (191) & 1{,}054.5 & 6 / 7 & 0.461\newline [0.436, 0.486] & 0.535\newline [0.506, 0.564] & 39\%\\
immune & 2{,}039; 256 (2) & 220 & 40 / 6 & 0.650\newline [0.590, 0.706] & 0.705\newline [0.623, 0.779] & 68\%\\
kidney~A & 1{,}972; 256 (65) & 436.5 & 20.5 / 11 & 0.555\newline [0.509, 0.599] & 0.545\newline [0.458, 0.632] & 71\%\\
\cmidrule(lr){1-7}
immune (seed 43) & 2{,}080; 256 (2) & 221 & 47 / 6 & 0.698\newline [0.643, 0.748] & 0.730\newline [0.651, 0.803] & 71\%\\
immune (seed 44) & 1{,}972; 256 (2) & 238 & 49 / 7 & 0.677\newline [0.621, 0.733] & 0.683\newline [0.603, 0.758] & 70\%\\
immune (donor TSP14) & 1{,}964; 256 (2) & 254.5 & 45 / 9 & 0.676\newline [0.619, 0.735] & 0.724\newline [0.639, 0.798] & 75\%\\
immune (donor TSP21) & 1{,}765; 256 (1) & 216 & 64 / 9 & 0.706\newline [0.651, 0.758] & 0.719\newline [0.638, 0.796] & 77\%\\
immune (donor TSP25) & 1{,}791; 256 (5) & 199 & 59.5 / 9 & 0.709\newline [0.651, 0.762] & 0.707\newline [0.621, 0.784] & 76\%\\
immune (B cells) & 1{,}398; 256 (0) & 197 & 54 / 4 & 0.708\newline [0.648, 0.763] & 0.703\newline [0.621, 0.789] & 72\%\\
immune (CD4 T cells) & 1{,}465; 256 (0) & 185 & 66 / 4 & 0.720\newline [0.659, 0.780] & 0.739\newline [0.657, 0.812] & 73\%\\
immune (CD8 T cells) & 1{,}595; 256 (0) & 201 & 43.5 / 4 & 0.650\newline [0.587, 0.708] & 0.682\newline [0.598, 0.763] & 59\%\\
kidney (whole-human ckpt) & 4{,}923; 192 (188) & 998 & 5 / 7 & 0.450\newline [0.427, 0.475] & 0.525\newline [0.495, 0.553] & 36\%\\
kidney (kidney ckpt) & 4{,}923; 192 (188) & 998 & 5 / 7 & 0.450\newline [0.427, 0.475] & 0.525\newline [0.497, 0.552] & 36\%\\
kidney (brain ckpt) & 4{,}923; 192 (188) & 998 & 5 / 7 & 0.450\newline [0.427, 0.474] & 0.525\newline [0.497, 0.552] & 36\%\\
\midrule
\multicolumn{7}{@{}p{16.4cm}@{}}{\footnotesize\textbf{E. The two kidney extractions (one donor, different gene axes): mechanism, matched re-evaluation and the curated-gene comparison}}\\*
Quantity & \multicolumn{2}{>{\scriptsize\raggedright\arraybackslash}p{4.2cm}}{kidney~A} & \multicolumn{2}{>{\scriptsize\raggedright\arraybackslash}p{4.2cm}}{kidney~B} & \multicolumn{2}{>{\scriptsize\raggedright\arraybackslash}p{4.2cm}}{Shared / difference A $-$ B [95\% CI]}\\
\cmidrule(lr){1-7}
Gene axis $G$ & \multicolumn{2}{>{\scriptsize\raggedright\arraybackslash}p{4.2cm}}{4{,}803} & \multicolumn{2}{>{\scriptsize\raggedright\arraybackslash}p{4.2cm}}{8{,}143} & \multicolumn{2}{>{\scriptsize\raggedright\arraybackslash}p{4.2cm}}{2{,}766 shared}\\
Observed genes $|V|$ & \multicolumn{2}{>{\scriptsize\raggedright\arraybackslash}p{4.2cm}}{1{,}972} & \multicolumn{2}{>{\scriptsize\raggedright\arraybackslash}p{4.2cm}}{4{,}780} & \multicolumn{2}{>{\scriptsize\raggedright\arraybackslash}p{4.2cm}}{1{,}126 observed in both ($M$)}\\
Lambert TFs in $V$ & \multicolumn{2}{>{\scriptsize\raggedright\arraybackslash}p{4.2cm}}{90} & \multicolumn{2}{>{\scriptsize\raggedright\arraybackslash}p{4.2cm}}{511} & \multicolumn{2}{>{\scriptsize\raggedright\arraybackslash}p{4.2cm}}{76 in $M$; 49 TRRUST TFs / 103 target-only}\\
Cells (input $>512$ genes) & \multicolumn{2}{>{\scriptsize\raggedright\arraybackslash}p{4.2cm}}{256 (65)} & \multicolumn{2}{>{\scriptsize\raggedright\arraybackslash}p{4.2cm}}{192 (191)} & \multicolumn{2}{>{\scriptsize\raggedright\arraybackslash}p{4.2cm}}{12 cells drawn in both}\\
Median non-zero genes per cell & \multicolumn{2}{>{\scriptsize\raggedright\arraybackslash}p{4.2cm}}{436.5} & \multicolumn{2}{>{\scriptsize\raggedright\arraybackslash}p{4.2cm}}{1{,}054.5} & \multicolumn{2}{>{\scriptsize\raggedright\arraybackslash}p{4.2cm}}{--}\\
Mean fraction of a cell's detected genes not on the other axis & \multicolumn{2}{>{\scriptsize\raggedright\arraybackslash}p{4.2cm}}{0.229} & \multicolumn{2}{>{\scriptsize\raggedright\arraybackslash}p{4.2cm}}{0.680} & \multicolumn{2}{>{\scriptsize\raggedright\arraybackslash}p{4.2cm}}{--}\\
Median $n_{\mathrm{obs}}$ on $M$: Lambert TFs / TRRUST TFs / non-TFs & \multicolumn{2}{>{\scriptsize\raggedright\arraybackslash}p{4.2cm}}{24 / 24 / 32} & \multicolumn{2}{>{\scriptsize\raggedright\arraybackslash}p{4.2cm}}{14 / 15 / 17} & \multicolumn{2}{>{\scriptsize\raggedright\arraybackslash}p{4.2cm}}{Spearman $\rho$ of $n_{\mathrm{obs}}$ on $M$ 0.970}\\
\cmidrule(lr){1-7}
Matched genes $M$: S2b, band $\SV{2}$--$\SV{7}$, each extraction's own SVD & \multicolumn{2}{>{\scriptsize\raggedright\arraybackslash}p{4.2cm}}{0.628} & \multicolumn{2}{>{\scriptsize\raggedright\arraybackslash}p{4.2cm}}{0.573} & \multicolumn{2}{>{\scriptsize\raggedright\arraybackslash}p{4.2cm}}{$+$0.055 [$-$0.033, 0.134]}\\
Matched genes $M$: S2b, band $\SV{2}$--$\SV{7}$, SVD fitted on $M$ & \multicolumn{2}{>{\scriptsize\raggedright\arraybackslash}p{4.2cm}}{0.664} & \multicolumn{2}{>{\scriptsize\raggedright\arraybackslash}p{4.2cm}}{0.679} & \multicolumn{2}{>{\scriptsize\raggedright\arraybackslash}p{4.2cm}}{$-$0.016 [$-$0.057, 0.024]}\\
Matched genes $M$: S2b, full dimension & \multicolumn{2}{>{\scriptsize\raggedright\arraybackslash}p{4.2cm}}{0.676} & \multicolumn{2}{>{\scriptsize\raggedright\arraybackslash}p{4.2cm}}{0.646} & \multicolumn{2}{>{\scriptsize\raggedright\arraybackslash}p{4.2cm}}{$+$0.030 [$-$0.023, 0.088]}\\
Matched genes $M$: S2, own detection strata & \multicolumn{2}{>{\scriptsize\raggedright\arraybackslash}p{4.2cm}}{0.647} & \multicolumn{2}{>{\scriptsize\raggedright\arraybackslash}p{4.2cm}}{0.637} & \multicolumn{2}{>{\scriptsize\raggedright\arraybackslash}p{4.2cm}}{$+$0.011 [$-$0.035, 0.056]}\\
Matched genes $M$: S2, shared detection strata & \multicolumn{2}{>{\scriptsize\raggedright\arraybackslash}p{4.2cm}}{0.637} & \multicolumn{2}{>{\scriptsize\raggedright\arraybackslash}p{4.2cm}}{0.640} & \multicolumn{2}{>{\scriptsize\raggedright\arraybackslash}p{4.2cm}}{$-$0.003 [$-$0.039, 0.035]}\\
Basis instability: $|\Delta_{\mathrm{band}}| - |\Delta_{\mathrm{full}}|$ (own SVD) & \multicolumn{2}{>{\scriptsize\raggedright\arraybackslash}p{4.2cm}}{--} & \multicolumn{2}{>{\scriptsize\raggedright\arraybackslash}p{4.2cm}}{--} & \multicolumn{2}{>{\scriptsize\raggedright\arraybackslash}p{4.2cm}}{$+$0.025 [$-$0.053, 0.092]}\\
\cmidrule(lr){1-7}
TRRUST comparison on the curated genes & \multicolumn{2}{>{\scriptsize\raggedright\arraybackslash}p{4.2cm}}{kidney~A, 187 genes (65 / 122)} & \multicolumn{2}{>{\scriptsize\raggedright\arraybackslash}p{4.2cm}}{kidney~A on kidney~B's 167 genes} & \multicolumn{2}{>{\scriptsize\raggedright\arraybackslash}p{4.2cm}}{kidney~B, 167 genes (59 / 108)}\\
\cmidrule(lr){1-7}
Band $\SV{2}$--$\SV{7}$, SVD including never-observed zero rows (the screen's convention), layer mean & \multicolumn{2}{>{\scriptsize\raggedright\arraybackslash}p{4.2cm}}{0.612} & \multicolumn{2}{>{\scriptsize\raggedright\arraybackslash}p{4.2cm}}{0.636} & \multicolumn{2}{>{\scriptsize\raggedright\arraybackslash}p{4.2cm}}{0.590\newline $L_0$--$L_3$: 0.493, 0.500, 0.529, 0.551; $L_4$--$L_{11}$: 0.589--0.664}\\
Same band, observed-gene SVD, layer mean & \multicolumn{2}{>{\scriptsize\raggedright\arraybackslash}p{4.2cm}}{0.609} & \multicolumn{2}{>{\scriptsize\raggedright\arraybackslash}p{4.2cm}}{0.637} & \multicolumn{2}{>{\scriptsize\raggedright\arraybackslash}p{4.2cm}}{0.583\newline $L_0$--$L_3$: 0.677, 0.665, 0.648, 0.665}\\
Full dimension, layer mean & \multicolumn{2}{>{\scriptsize\raggedright\arraybackslash}p{4.2cm}}{0.676} & \multicolumn{2}{>{\scriptsize\raggedright\arraybackslash}p{4.2cm}}{0.681} & \multicolumn{2}{>{\scriptsize\raggedright\arraybackslash}p{4.2cm}}{0.644}\\
Hanley--McNeil SE of one AUROC & \multicolumn{2}{>{\scriptsize\raggedright\arraybackslash}p{4.2cm}}{0.044 (at 0.61)} & \multicolumn{2}{>{\scriptsize\raggedright\arraybackslash}p{4.2cm}}{--} & \multicolumn{2}{>{\scriptsize\raggedright\arraybackslash}p{4.2cm}}{0.047 (at 0.55)}\\
\end{longtable}
\endgroup

\setcounter{table}{9}
\begingroup
\footnotesize
\setlength{\tabcolsep}{3pt}
\renewcommand{\arraystretch}{1.05}
\begin{longtable}{@{}>{\scriptsize\raggedright\arraybackslash}p{2.9cm}>{\scriptsize\raggedright\arraybackslash}p{2.0cm}*{2}{>{\scriptsize\raggedright\arraybackslash}p{2.1cm}}*{3}{>{\scriptsize\raggedright\arraybackslash}p{2.2cm}}@{}}
\caption{\textbf{Cross-model comparison: gene sets, eleven static representations on 17{,}693 genes, scGPT's residual stream against static representations, and the bias of a 186-gene curated set.} Unless stated otherwise, values in Panels B--D are AUROCs of full-dimensional vectors centred over the analysed gene set and compared by cosine. \emph{Lambert}: leave-one-out (LOO) mean cosine to the census transcription factors (TFs) of Lambert et al., TFs against all other genes; \emph{TRRUST}: the same score, TRRUST TFs against genes that are only targets; \emph{PPI}: STRING physical pairs (score $\geq 700$) against non-pairs from one degree-preserving rewiring; \emph{edge}: within-TF AUROC of the TF--target cosine, TRRUST edges against candidate non-edges (targets of other TFs, absent from TRRUST and DoRothEA). 95\% intervals: node tasks, gene bootstrap ($B = 1{,}000$); PPI, vertex subsampling ($m = n/2$, $B = 200$) with delete-$d$ scaling; edge, two-way TF $\times$ target bootstrap ($B = 2{,}000$); paired differences reuse the same genes and draws for both representations. $d$, dimension; LM, language model; HPA, Human Protein Atlas; $G$, gene axis; $V$, observed genes. \textbf{A}: Tier A, genes in all 11 vocabularies; Tier B, $V \cap {}$Tier A; the curated set starts from the automated screen's kidney~A TRRUST edge list. \textbf{B}: token table, the model's input gene-embedding table ($^{a}$with its LayerNorm applied); $^{b}$gene outputs of one forward pass on the HPA mean-expression profile (no gene-token table, so not like-for-like with the token tables); ESM2-15B, the 5{,}120-dimensional per-protein embedding distributed with UCE; ESM2-650M, mean-pooled over the canonical sequence; HPA, single-cell expression over 154 cell types, 128 principal components. Lambert is pooled within lung detection deciles (unstratified in parentheses); the 12{,}467 genes not observed in lung all sit in the median decile, so this stratification is largely nominal (Panel C controls detection); within-decile permutation $p = 1/2{,}001$ throughout. TRRUST depends on the centring set: all Tier A genes, or only the 2{,}753 labelled genes (775 TFs, 1{,}978 target-only; point estimates; the convention of Panel D). Labelled-gene centring raises the single-cell models other than scFoundation by 0.048--0.119 and moves the others by at most 0.010, roughly halving the gap between protein and expression-trained representations (last two rows: paired differences from scGPT [95\% CI], for Lambert of the stratified values). \textbf{C}: scGPT whole-human residual stream, one extraction per context (192--256 cells), centred over the Tier B set; Lambert is stratified within the extraction's deciles of $\log n_{\mathrm{obs}}$ (the H2 statistic), the other tasks are unstratified. Positives / comparison set: TFs / other genes (Lambert), TFs / target-only genes (TRRUST), pairs / rewired non-pairs (PPI), documented edges / candidate non-edges (edge). Best layer: highest AUROC, chosen post hoc. Static columns: the representation's AUROC on the same genes; range over the 12 layers of the paired difference, static minus layer; $k_{+}/k_{-}$, layers whose interval lies entirely above / below 0. \textbf{D}: TRRUST AUROC (centred over the subset) on the curated set and in 10{,}000 random subsets of the 2{,}753 labelled Tier A genes matched on size, TF:target-only ratio (64:122 for 186) and lung mean-expression decile (degree-matched: TRRUST-degree quartiles and lung mean-expression quintiles; degree, TRRUST out-degree of a TF or in-degree of a target-only gene; pool, the labelled genes). Percentile: share of null subsets below the curated value (ties half); two-sided $p$ with $+1$ correction (floor $2.0\times10^{-4}$); excess, curated value minus null mean; $^{c}$scored on the 178 genes in tGPT's vocabulary. Kendall $\tau$: agreement of the representation orders, with a paired gene-bootstrap interval ($B = 1{,}000$). Projection: Gaussian random projection (5 seeds) at a 5:1 ratio, scored as in Panel B (Lambert unstratified). All comparisons are sensitivity analyses outside the pre-registered hypotheses. For PPI, the vertex-subsampling intervals with the pre-registered scaling factor $\sqrt{m/n}$, which omits the finite-population term and so gives half-widths smaller by a factor $\sqrt{2}$, are in the CSV (columns \texttt{ci\_lo\_preregistered}/\texttt{ci\_hi\_preregistered}). Full values: \texttt{Table\_S10\_cross\_model.csv}.}\label{tab:S10}\\
\toprule
\endfirsthead
\multicolumn{7}{@{}l}{\footnotesize Table~\thetable\ (continued)}\\
\toprule
\endhead
\midrule
\multicolumn{7}{r@{}}{\footnotesize continued on next page}\\
\endfoot
\bottomrule
\endlastfoot
\multicolumn{7}{@{}p{16.9cm}@{}}{\footnotesize\textbf{A. Gene sets: vocabularies, intersections, the two tiers and the curated set}}\\*
Gene set & Genes & \multicolumn{5}{>{\scriptsize\raggedright\arraybackslash}p{11.6cm}@{}}{Composition, or genes lost at this step}\\*
\cmidrule(lr){1-7}
Vocabulary of each static representation (genes with a non-zero vector) & 19{,}264--60{,}694 & \multicolumn{5}{>{\scriptsize\raggedright\arraybackslash}p{11.6cm}@{}}{scGPT (whole-human) 60{,}694; Geneformer V1-10M 21{,}243; Geneformer V2-104M 20{,}271; Geneformer V2-316M 20{,}271; tGPT 21{,}046; AIDO.Cell-100M 19{,}264; scFoundation 19{,}264; ESM2-15B 19{,}790; ESM2-650M 20{,}394; 32 sequence descriptors 20{,}213 (181 all-zero rows dropped); HPA co-expression 20{,}151}\\
Geneformer V1-10M $\cap$ V2-104M $\cap$ ESM2-15B & 18{,}907 & \multicolumn{5}{>{\scriptsize\raggedright\arraybackslash}p{11.6cm}@{}}{the vocabularies used to define the 186-gene curated set (below)}\\
Tier A: all 11 vocabularies & 17{,}693 & \multicolumn{5}{>{\scriptsize\raggedright\arraybackslash}p{11.6cm}@{}}{Lambert TFs 1{,}587 vs 16{,}106 other genes; TRRUST TFs 775 vs 1{,}978 target-only genes; STRING physical pairs 73{,}430 vs 72{,}416 rewired non-pairs; TRRUST edges 7{,}593 (430 TFs) vs 70{,}882 candidate non-edges; 12{,}467 genes not observed in lung (assigned the median lung detection decile)}\\
Tier B, lung & 5{,}226 & \multicolumn{5}{>{\scriptsize\raggedright\arraybackslash}p{11.6cm}@{}}{gene axis $G$ 8{,}181; observed $V$ 5{,}648; Lambert TFs in Tier B 579; genes in a TRRUST edge within Tier B 1{,}912; 256 cells}\\
Tier B, lung (ext.) & 5{,}154 & \multicolumn{5}{>{\scriptsize\raggedright\arraybackslash}p{11.6cm}@{}}{gene axis $G$ 8{,}229; observed $V$ 5{,}616; Lambert TFs in Tier B 618; genes in a TRRUST edge within Tier B 1{,}948; 256 cells}\\
Tier B, kidney~B & 4{,}441 & \multicolumn{5}{>{\scriptsize\raggedright\arraybackslash}p{11.6cm}@{}}{gene axis $G$ 8{,}143; observed $V$ 4{,}780; Lambert TFs in Tier B 509; genes in a TRRUST edge within Tier B 1{,}542; 192 cells}\\
Tier B, immune & 1{,}338 & \multicolumn{5}{>{\scriptsize\raggedright\arraybackslash}p{11.6cm}@{}}{gene axis $G$ 4{,}941; observed $V$ 2{,}039; Lambert TFs in Tier B 111; genes in a TRRUST edge within Tier B 281; 256 cells}\\
Tier B, kidney~A & 1{,}298 & \multicolumn{5}{>{\scriptsize\raggedright\arraybackslash}p{11.6cm}@{}}{gene axis $G$ 4{,}803; observed $V$ 1{,}972; Lambert TFs in Tier B 90; genes in a TRRUST edge within Tier B 172; 256 cells}\\
Genes named in the curated kidney~A edge list & 209 & \multicolumn{5}{>{\scriptsize\raggedright\arraybackslash}p{11.6cm}@{}}{1{,}152 TRRUST edge and sampled non-edge rows, so the list includes genes that occur only in non-edges; all 209 are on the kidney~A gene axis}\\
\quad observed in kidney~A & 187 & \multicolumn{5}{>{\scriptsize\raggedright\arraybackslash}p{11.6cm}@{}}{22 never-observed genes dropped}\\
\quad $\cap$ Geneformer V1-10M, V2-104M, ESM2-15B: the 186-gene set & 186 & \multicolumn{5}{>{\scriptsize\raggedright\arraybackslash}p{11.6cm}@{}}{lost: NEAT1 (lncRNA); 64 TRRUST TFs, 122 target-only genes; 13 genes not observed in lung}\\
\quad $\cap$ Tier A: the 178-gene set & 178 & \multicolumn{5}{>{\scriptsize\raggedright\arraybackslash}p{11.6cm}@{}}{lost (absent from tGPT): HLA-A, HLA-B, HLA-C, HLA-DMB, HLA-DOB, HLA-DPB1, HLA-DRA, HLA-DRB1; 64 TRRUST TFs, 114 target-only genes}\\
Kidney~A $V \cap {}$Tier A, in a TRRUST edge & 172 & \multicolumn{5}{>{\scriptsize\raggedright\arraybackslash}p{11.6cm}@{}}{the most a kidney~A-only analysis of the TRRUST tasks could use (kidney~A observes 1{,}972 of its 4{,}803 genes)}\\
\midrule
\multicolumn{7}{@{}p{16.9cm}@{}}{\footnotesize\textbf{B. Tier A (17{,}693 genes): every static representation, full dimension, centred over the analysed genes}}\\*
Representation & Type; $d$ & Lambert: stratified (unstratified) [95\% CI] & TRRUST, centred over Tier A [95\% CI] & TRRUST, centred over labelled genes & PPI [95\% CI, delete-$d$] & Edge, within TF [95\% CI]\\*
\cmidrule(lr){1-7}
scGPT (whole-human) & token table; 512 & 0.747 (0.704)\newline [0.734, 0.760] & 0.684\newline [0.660, 0.708] & 0.789 & 0.669\newline [0.658, 0.679] & 0.575\newline [0.534, 0.610]\\
Geneformer V1-10M & token table$^{a}$; 256 & 0.776 (0.742)\newline [0.761, 0.789] & 0.688\newline [0.659, 0.715] & 0.755 & 0.588\newline [0.579, 0.597] & 0.529\newline [0.502, 0.562]\\
Geneformer V2-104M & token table$^{a}$; 768 & 0.864 (0.836)\newline [0.851, 0.876] & 0.774\newline [0.749, 0.796] & 0.821 & 0.662\newline [0.650, 0.671] & 0.580\newline [0.549, 0.610]\\
Geneformer V2-316M & token table$^{a}$; 1{,}152 & 0.855 (0.817)\newline [0.842, 0.867] & 0.743\newline [0.717, 0.765] & 0.802 & 0.659\newline [0.650, 0.670] & 0.578\newline [0.542, 0.611]\\
tGPT & token table; 1{,}024 & 0.690 (0.639)\newline [0.678, 0.703] & 0.520\newline [0.497, 0.543] & 0.639 & 0.571\newline [0.561, 0.582] & 0.544\newline [0.520, 0.570]\\
\cmidrule(lr){1-7}
AIDO.Cell-100M & forward pass$^{b}$; 640 & 0.623 (0.589)\newline [0.611, 0.636] & 0.579\newline [0.560, 0.600] & 0.690 & 0.603\newline [0.590, 0.614] & 0.552\newline [0.499, 0.591]\\
scFoundation & forward pass$^{b}$; 512 & 0.554 (0.532)\newline [0.537, 0.570] & 0.487\newline [0.458, 0.511] & 0.496 & 0.559\newline [0.544, 0.571] & 0.500\newline [0.473, 0.530]\\
\cmidrule(lr){1-7}
ESM2-15B & protein LM (UCE file); 5{,}120 & 0.902 (0.889)\newline [0.893, 0.911] & 0.870\newline [0.856, 0.884] & 0.877 & 0.647\newline [0.638, 0.657] & 0.491\newline [0.466, 0.523]\\
ESM2-650M & protein LM, mean-pooled; 1{,}280 & 0.861 (0.859)\newline [0.851, 0.870] & 0.826\newline [0.811, 0.843] & 0.826 & 0.610\newline [0.600, 0.618] & 0.492\newline [0.466, 0.530]\\
32 sequence descriptors & hand-computed; 32 & 0.916 (0.902)\newline [0.908, 0.924] & 0.869\newline [0.855, 0.883] & 0.876 & 0.671\newline [0.660, 0.680] & 0.497\newline [0.465, 0.538]\\
\cmidrule(lr){1-7}
HPA co-expression & 154 cell types, PCA; 128 & 0.617 (0.576)\newline [0.604, 0.630] & 0.615\newline [0.590, 0.640] & 0.618 & 0.598\newline [0.588, 0.608] & 0.518\newline [0.485, 0.560]\\
\cmidrule(lr){1-7}
ESM2-15B minus scGPT (paired) &  & $+$0.155\newline [0.141, 0.168] & $+$0.186\newline [0.159, 0.212] & $+$0.088\newline [0.066, 0.111] & $-$0.022\newline [$-$0.032, $-$0.011] & $-$0.084\newline [$-$0.124, $-$0.042]\\
Geneformer V2-104M minus scGPT (paired) &  & $+$0.117\newline [0.104, 0.130] & $+$0.089\newline [0.066, 0.111] & $+$0.032\newline [0.014, 0.051] & $-$0.007\newline [$-$0.012, $-$0.002] & $+$0.005\newline [$-$0.021, 0.029]\\
\midrule
\multicolumn{7}{@{}p{16.9cm}@{}}{\footnotesize\textbf{C. Tier B: scGPT's residual stream ($L_0$--$L_{11}$) against three static representations on each extraction's observed genes}}\\*
Extraction: task & Positives / comparison set & Residual stream: range over 12 layers & Best layer: AUROC [95\% CI] & ESM2-15B: AUROC; minus layer; $k_{+}/k_{-}$ & scGPT token table: AUROC; minus layer; $k_{+}/k_{-}$ & Geneformer V2-104M: AUROC; minus layer; $k_{+}/k_{-}$\\*
\cmidrule(lr){1-7}
lung: Lambert (stratified) & 579 / 4{,}647 & 0.638--0.668 & $L_{1}$: 0.668\newline [0.652, 0.687] & 0.880\newline 0.212--0.241\newline 12/0 & 0.751\newline 0.083--0.113\newline 12/0 & 0.825\newline 0.157--0.187\newline 12/0\\
lung: TRRUST & 597 / 1{,}454 & 0.658--0.718 & $L_{2}$: 0.718\newline [0.699, 0.737] & 0.857\newline 0.139--0.199\newline 12/0 & 0.779\newline 0.060--0.121\newline 12/0 & 0.797\newline 0.079--0.139\newline 12/0\\
lung: PPI & 13{,}965 / 13{,}822 & 0.616--0.674 & $L_{1}$: 0.674\newline [0.658, 0.690] & 0.639\newline $-$0.036 to $+$0.022\newline 1/5 & 0.661\newline $-$0.013 to $+$0.044\newline 5/4 & 0.651\newline $-$0.023 to $+$0.035\newline 4/4\\
lung: edge & 5{,}131 / 45{,}609 (304 TFs) & 0.515--0.542 & $L_{1}$: 0.542\newline [0.494, 0.589] & 0.498\newline $-$0.043 to $-$0.017\newline 0/3 & 0.566\newline 0.024--0.051\newline 4/0 & 0.562\newline 0.020--0.047\newline 3/0\\
\cmidrule(lr){1-7}
lung (ext.): Lambert (stratified) & 618 / 4{,}536 & 0.631--0.675 & $L_{1}$: 0.675\newline [0.654, 0.698] & 0.881\newline 0.206--0.250\newline 12/0 & 0.732\newline 0.056--0.100\newline 12/0 & 0.833\newline 0.158--0.201\newline 12/0\\
lung (ext.): TRRUST & 621 / 1{,}450 & 0.672--0.725 & $L_{1}$: 0.725\newline [0.707, 0.745] & 0.857\newline 0.132--0.185\newline 12/0 & 0.775\newline 0.050--0.103\newline 12/0 & 0.807\newline 0.082--0.135\newline 12/0\\
lung (ext.): PPI & 13{,}714 / 13{,}561 & 0.633--0.678 & $L_{2}$: 0.678\newline [0.660, 0.694] & 0.643\newline $-$0.035 to $+$0.010\newline 0/6 & 0.658\newline $-$0.020 to $+$0.026\newline 5/4 & 0.652\newline $-$0.026 to $+$0.020\newline 4/5\\
lung (ext.): edge & 5{,}035 / 44{,}994 (311 TFs) & 0.520--0.539 & $L_{1}$: 0.539\newline [0.492, 0.592] & 0.494\newline $-$0.045 to $-$0.025\newline 0/3 & 0.562\newline 0.023--0.042\newline 4/0 & 0.569\newline 0.029--0.049\newline 3/0\\
\cmidrule(lr){1-7}
kidney~B: Lambert (stratified) & 509 / 3{,}932 & 0.616--0.657 & $L_{1}$: 0.657\newline [0.635, 0.680] & 0.861\newline 0.204--0.245\newline 12/0 & 0.737\newline 0.080--0.120\newline 12/0 & 0.819\newline 0.162--0.202\newline 12/0\\
kidney~B: TRRUST & 548 / 1{,}188 & 0.637--0.711 & $L_{1}$: 0.711\newline [0.689, 0.732] & 0.843\newline 0.132--0.206\newline 12/0 & 0.754\newline 0.043--0.117\newline 12/0 & 0.791\newline 0.080--0.154\newline 12/0\\
kidney~B: PPI & 11{,}990 / 11{,}830 & 0.617--0.671 & $L_{1}$: 0.671\newline [0.649, 0.688] & 0.630\newline $-$0.041 to $+$0.012\newline 0/6 & 0.659\newline $-$0.012 to $+$0.042\newline 6/2 & 0.646\newline $-$0.025 to $+$0.028\newline 4/4\\
kidney~B: edge & 3{,}569 / 31{,}590 (248 TFs) & 0.522--0.553 & $L_{0}$: 0.553\newline [0.500, 0.602] & 0.500\newline $-$0.052 to $-$0.022\newline 0/3 & 0.568\newline 0.015--0.046\newline 4/0 & 0.569\newline 0.017--0.048\newline 4/0\\
\cmidrule(lr){1-7}
immune: Lambert (stratified) & 111 / 1{,}227 & 0.616--0.746 & $L_{2}$: 0.746\newline [0.677, 0.819] & 0.933\newline 0.187--0.317\newline 12/0 & 0.811\newline 0.065--0.195\newline 12/0 & 0.815\newline 0.069--0.199\newline 11/0\\
immune: TRRUST & 95 / 350 & 0.643--0.774 & $L_{2}$: 0.774\newline [0.720, 0.824] & 0.943\newline 0.169--0.300\newline 12/0 & 0.800\newline 0.026--0.156\newline 10/0 & 0.832\newline 0.059--0.189\newline 12/0\\
immune: PPI & 1{,}185 / 1{,}155 & 0.673--0.742 & $L_{2}$: 0.742\newline [0.701, 0.789] & 0.693\newline $-$0.050 to $+$0.020\newline 0/0 & 0.727\newline $-$0.015 to $+$0.054\newline 4/0 & 0.720\newline $-$0.022 to $+$0.047\newline 2/0\\
immune: edge & 544 / 2{,}976 (36 TFs) & 0.576--0.632 & $L_{0}$: 0.632\newline [0.570, 0.706] & 0.483\newline $-$0.149 to $-$0.093\newline 0/8 & 0.650\newline 0.018--0.075\newline 0/0 & 0.634\newline 0.002--0.058\newline 0/0\\
\cmidrule(lr){1-7}
kidney~A: Lambert (stratified) & 90 / 1{,}208 & 0.637--0.763 & $L_{3}$: 0.763\newline [0.691, 0.830] & 0.914\newline 0.151--0.277\newline 12/0 & 0.779\newline 0.016--0.142\newline 9/0 & 0.804\newline 0.042--0.167\newline 7/0\\
kidney~A: TRRUST & 91 / 288 & 0.707--0.828 & $L_{2}$: 0.828\newline [0.758, 0.885] & 0.940\newline 0.112--0.233\newline 12/0 & 0.831\newline 0.003--0.124\newline 7/0 & 0.839\newline 0.011--0.131\newline 6/0\\
kidney~A: PPI & 1{,}235 / 1{,}176 & 0.746--0.786 & $L_{1}$: 0.786\newline [0.741, 0.831] & 0.730\newline $-$0.055 to $-$0.016\newline 0/2 & 0.785\newline $-$0.001 to $+$0.039\newline 4/0 & 0.772\newline $-$0.014 to $+$0.025\newline 0/0\\
kidney~A: edge & 201 / 1{,}329 (28 TFs) & 0.618--0.664 & $L_{0}$: 0.664\newline [0.570, 0.761] & 0.529\newline $-$0.135 to $-$0.089\newline 0/2 & 0.650\newline $-$0.014 to $+$0.032\newline 0/0 & 0.636\newline $-$0.028 to $+$0.018\newline 0/0\\
\midrule
\multicolumn{7}{@{}p{16.9cm}@{}}{\footnotesize\textbf{D. The 186- and 178-gene curated sets against matched random subsets of Tier A (TRRUST task), and a projection control}}\\*
Representation & 186 genes: AUROC & 186: null mean [2.5\%, 97.5\%] & 186: percentile; two-sided $p$ & 186: excess over null mean & 178 genes: AUROC; percentile; excess & 186, degree-matched: null mean; percentile; excess\\*
\cmidrule(lr){1-7}
scGPT (whole-human) & 0.833 & 0.739\newline [0.661, 0.810] & 0.9950\newline $p = 0.010$ & $+$0.095 & 0.821; 0.9878\newline $+$0.085 & 0.770; 0.9787\newline $+$0.064\\
Geneformer V1-10M & 0.802 & 0.646\newline [0.536, 0.743] & 0.9996\newline $p = 1.0\times10^{-3}$ & $+$0.156 & 0.797; 0.9992\newline $+$0.149 & 0.667; 0.9987\newline $+$0.135\\
Geneformer V2-104M & 0.874 & 0.745\newline [0.655, 0.824] & 0.9992\newline $p = 1.8\times10^{-3}$ & $+$0.129 & 0.867; 0.9994\newline $+$0.122 & 0.771; 0.9996\newline $+$0.104\\
Geneformer V2-316M & 0.848 & 0.715\newline [0.624, 0.795] & 0.9996\newline $p = 1.0\times10^{-3}$ & $+$0.133 & 0.834; 0.9989\newline $+$0.118 & 0.747; 0.9980\newline $+$0.101\\
tGPT & 0.723$^{c}$ & 0.587\newline [0.518, 0.669] & 0.9996\newline $p = 1.0\times10^{-3}$ & $+$0.136 & 0.723; 0.9997\newline $+$0.139 & 0.622; 0.9947\newline $+$0.101\\
\cmidrule(lr){1-7}
AIDO.Cell-100M & 0.769 & 0.648\newline [0.582, 0.716] & 0.9997\newline $p = 8.0\times10^{-4}$ & $+$0.121 & 0.759; 0.9990\newline $+$0.109 & 0.688; 0.9946\newline $+$0.082\\
scFoundation & 0.549 & 0.300\newline [0.179, 0.434] & 0.9994\newline $p = 1.4\times10^{-3}$ & $+$0.250 & 0.533; 0.9972\newline $+$0.248 & 0.301; 0.9994\newline $+$0.249\\
\cmidrule(lr){1-7}
ESM2-15B & 0.954 & 0.862\newline [0.809, 0.911] & 1.0000\newline $p = 2.0\times10^{-4}$ & $+$0.092 & 0.949; 1.0000\newline $+$0.087 & 0.892; 0.9989\newline $+$0.062\\
ESM2-650M & 0.907 & 0.819\newline [0.763, 0.872] & 0.9999\newline $p = 4.0\times10^{-4}$ & $+$0.088 & 0.902; 0.9994\newline $+$0.083 & 0.847; 0.9942\newline $+$0.060\\
32 sequence descriptors & 0.966 & 0.865\newline [0.811, 0.913] & 1.0000\newline $p = 2.0\times10^{-4}$ & $+$0.101 & 0.962; 1.0000\newline $+$0.097 & 0.894; 1.0000\newline $+$0.071\\
\cmidrule(lr){1-7}
HPA co-expression & 0.688 & 0.594\newline [0.535, 0.649] & 0.9997\newline $p = 8.0\times10^{-4}$ & $+$0.094 & 0.685; 0.9970\newline $+$0.075 & 0.624; 0.9892\newline $+$0.063\\
\cmidrule(lr){1-7}
Kendall $\tau$, order of the representations, curated set vs Tier A &  & \multicolumn{5}{>{\scriptsize\raggedright\arraybackslash}p{11.6cm}@{}}{186 genes (10 representations, without tGPT): 0.867 [0.689, 0.956]; 178 genes (11): 0.855 [0.673, 0.964]; 178 genes vs Tier A centred over its TRRUST-labelled genes: 0.964 (point estimate)}\\
Mean TRRUST degree, 186-gene set vs pool &  & \multicolumn{5}{>{\scriptsize\raggedright\arraybackslash}p{11.6cm}@{}}{TFs 16.6 vs 10.8; target-only genes 7.0 vs 3.2}\\
\cmidrule(lr){1-7}
Projection control (Tier A): representation, task & $d$: full $\to$ projected & Full: AUROC [95\% CI] & Projected: mean (SD) over 5 seeds & Projected: min--max & Projected minus full & scGPT token table, full 512-$d$\\*
\cmidrule(lr){1-7}
ESM2-15B: Lambert (unstratified) & 5{,}120 $\to$ 1{,}024 & 0.889\newline [0.882, 0.897] & 0.887 (0.007) & 0.877--0.895 & $-$0.002 & 0.704\\
ESM2-15B: TRRUST & 5{,}120 $\to$ 1{,}024 & 0.870\newline [0.856, 0.884] & 0.870 (0.007) & 0.861--0.877 & $-$0.0005 & 0.684\\
ESM2-15B: PPI & 5{,}120 $\to$ 1{,}024 & 0.647\newline [0.638, 0.657] & 0.648 (0.004) & 0.641--0.650 & $+$0.001 & 0.669\\
ESM2-15B: edge & 5{,}120 $\to$ 1{,}024 & 0.491\newline [0.466, 0.523] & 0.490 (0.002) & 0.487--0.493 & $-$0.001 & 0.575\\
scGPT (whole-human): Lambert (unstratified) & 512 $\to$ 102 & 0.704\newline [0.694, 0.713] & 0.678 (0.005) & 0.672--0.685 & $-$0.027 & --\\
scGPT (whole-human): TRRUST & 512 $\to$ 102 & 0.684\newline [0.660, 0.708] & 0.635 (0.025) & 0.601--0.667 & $-$0.049 & --\\
scGPT (whole-human): PPI & 512 $\to$ 102 & 0.669\newline [0.658, 0.679] & 0.634 (0.002) & 0.631--0.636 & $-$0.036 & --\\
scGPT (whole-human): edge & 512 $\to$ 102 & 0.575\newline [0.534, 0.610] & 0.556 (0.012) & 0.540--0.572 & $-$0.019 & --\\
\end{longtable}
\endgroup

\setcounter{table}{10}
\begingroup
\footnotesize
\setlength{\tabcolsep}{3pt}
\renewcommand{\arraystretch}{1.05}
\begin{longtable}{@{}>{\scriptsize\raggedright\arraybackslash}p{3.1cm}*{6}{>{\scriptsize\raggedright\arraybackslash}p{2.1cm}}@{}}
\caption{\textbf{A co-occurrence-only instrument: Word2Vec trained on the token sets scGPT receives, compared with scGPT's residual stream on H1--H3.} Word2Vec: skip-gram, 256 dimensions, 10 negative samples, trained on every cell of each extraction's population (lung 15{,}000, lung (ext.) 9{,}000, kidney~B 8{,}000, immune 20{,}000, kidney~A 11{,}376 cells), each tokenised as scGPT tokenises it (non-zero genes, the 512 most expressed), with token and cell order reshuffled on every pass. Every representation is mean-centred over the observed genes $V$ and scored on identical genes, pairs, labels and resampling draws; scGPT is the whole-human checkpoint. H1b, AUROC of the cosine of STRING physical pairs (combined score $\geq 700$) against one degree-preserving rewired non-pair set, with vertex-subsampling intervals ($m = n/2$, $B = 2{,}000$) under delete-$d$ scaling; H2, detection-stratified AUROC of the LOO-cosine TF score (Lambert TFs against all other observed genes), within-stratum permutation null ($N = 2{,}000$; floor $5.0\times10^{-4}$) and gene-bootstrap intervals ($B = 2{,}000$, scores held fixed); H3, within-TF stratified AUROC of the TF--target cosine with two-way TF $\times$ target bootstrap intervals ($B = 2{,}000$); H1a, assortativity of $\SV{1}$--$\SV{10}$ of the observed-gene SVD on the STRING network with a Maslov--Sneppen null ($N = 1{,}000$; max-$T$ over 10 directions for Word2Vec and over the 120 direction--layer tests for scGPT). These $p$-values are not those of the primary H1/H2 tests. Differences are paired (Word2Vec minus scGPT on the same draws); the best layer is the scGPT layer with the highest value. Under the within-stratum null the LOO score is not centred at 0.5 (null means 0.27--0.50 for Word2Vec), so H2 is judged against its own null, and the null-centred difference subtracts each instrument's null mean (a descriptive adjustment). Seed variation is reported separately and is not propagated into the intervals. \textbf{B}: one seed; $^{a}$the 10-epoch paired differences are in Panel A (mean of 5 seeds). \textbf{C}: the five-token window used for cost against a window spanning the whole cell (512), and the five-token window with its number of pair updates matched to the 512-token window (571 epochs for lung, 365 for immune). Pair updates, the number of (centre, context) updates in training. Per-seed values, every layer, the extraction-cells-only models for all extractions and all CKA values are in the CSV. For H1b, the vertex-subsampling intervals with the pre-registered scaling factor $\sqrt{m/n}$, which omits the finite-population term and so gives half-widths smaller by a factor $\sqrt{2}$, are also in the CSV (columns without the suffix \texttt{\_deleted}, beside the reported \texttt{\_deleted} columns). Full values: \texttt{Table\_S11\_cooccurrence.csv}.}\label{tab:S11}\\
\toprule
\endfirsthead
\multicolumn{7}{@{}l}{\footnotesize Table~\thetable\ (continued)}\\
\toprule
\endhead
\midrule
\multicolumn{7}{r@{}}{\footnotesize continued on next page}\\
\endfoot
\bottomrule
\endlastfoot
\multicolumn{7}{@{}p{16.4cm}@{}}{\footnotesize\textbf{A. Word2Vec trained on the full population corpus (five-token window, 10 epochs, mean of 5 seeds) against scGPT's residual stream, per extraction}}\\
\multicolumn{2}{@{}>{\scriptsize\raggedright\arraybackslash}p{5.41cm}}{Task and quantity} & lung & lung (ext.) & kidney~B & immune & kidney~A\\
\cmidrule(lr){1-7}
\multicolumn{7}{@{}l}{\scriptsize\textit{H1b: STRING pair-cosine AUROC against rewired non-pairs (delete-$d$ vertex-subsampling intervals)}}\\
\multicolumn{2}{@{}>{\scriptsize\raggedright\arraybackslash}p{5.41cm}}{Word2Vec [95\% CI]} & 0.605\newline [0.589, 0.621] & 0.597\newline [0.581, 0.613] & 0.605\newline [0.586, 0.622] & 0.704\newline [0.656, 0.751] & 0.739\newline [0.690, 0.794]\\
\multicolumn{2}{@{}>{\scriptsize\raggedright\arraybackslash}p{5.41cm}}{scGPT, 12 layers (best layer)} & 0.619--0.680 ($L_{1}$) & 0.630--0.677 ($L_{1}$) & 0.621--0.675 ($L_{1}$) & 0.668--0.743 ($L_{2}$) & 0.749--0.814 ($L_{1}$)\\
\multicolumn{2}{@{}>{\scriptsize\raggedright\arraybackslash}p{5.41cm}}{Word2Vec minus scGPT at its best layer} & $-$0.075\newline [$-$0.090, $-$0.059] & $-$0.080\newline [$-$0.097, $-$0.062] & $-$0.071\newline [$-$0.089, $-$0.052] & $-$0.040\newline [$-$0.082, 0.005] & $-$0.074\newline [$-$0.107, $-$0.042]\\
\multicolumn{2}{@{}>{\scriptsize\raggedright\arraybackslash}p{5.41cm}}{Word2Vec minus scGPT at $L_{11}$} & $-$0.014\newline [$-$0.027, $-$0.000] & $-$0.032\newline [$-$0.049, $-$0.015] & $-$0.016\newline [$-$0.029, $-$0.002] & $+$0.036\newline [$-$0.005, 0.078] & $-$0.015\newline [$-$0.037, 0.008]\\
\multicolumn{2}{@{}>{\scriptsize\raggedright\arraybackslash}p{5.41cm}}{Layers with the difference's CI $<0$ (scGPT higher) / $>0$} & 12 / 0 & 12 / 0 & 12 / 0 & 0 / 0 & 8 / 0\\
\multicolumn{7}{@{}l}{\scriptsize\textit{H2: detection-stratified TF AUROC (LOO cosine; gene-bootstrap intervals)}}\\
\multicolumn{2}{@{}>{\scriptsize\raggedright\arraybackslash}p{5.41cm}}{Word2Vec [95\% CI]} & 0.632\newline [0.607, 0.656] & 0.602\newline [0.577, 0.628] & 0.601\newline [0.573, 0.629] & 0.511\newline [0.454, 0.569] & 0.469\newline [0.411, 0.529]\\
\multicolumn{2}{@{}>{\scriptsize\raggedright\arraybackslash}p{5.41cm}}{Word2Vec permutation $p$ (null mean)} & $5.0\times10^{-4}$\newline (0.494) & $5.0\times10^{-4}$\newline (0.497) & $5.0\times10^{-4}$\newline (0.494) & 0.22\newline (0.485) & 0.28\newline (0.444)\\
\multicolumn{2}{@{}>{\scriptsize\raggedright\arraybackslash}p{5.41cm}}{scGPT, 12 layers (best layer)} & 0.625--0.665 ($L_{5}$) & 0.627--0.662 ($L_{1}$) & 0.615--0.654 ($L_{1}$) & 0.566--0.720 ($L_{2}$) & 0.678--0.783 ($L_{3}$)\\
\multicolumn{2}{@{}>{\scriptsize\raggedright\arraybackslash}p{5.41cm}}{Word2Vec minus scGPT at its best layer, raw} & $-$0.033\newline [$-$0.060, $-$0.007] & $-$0.060\newline [$-$0.081, $-$0.038] & $-$0.053\newline [$-$0.078, $-$0.028] & $-$0.208\newline [$-$0.285, $-$0.124] & $-$0.314\newline [$-$0.384, $-$0.238]\\
\multicolumn{2}{@{}>{\scriptsize\raggedright\arraybackslash}p{5.41cm}}{Same, null-centred (each instrument's permutation-null mean subtracted)} & $-$0.033\newline [$-$0.060, $-$0.007] & $-$0.061\newline [$-$0.082, $-$0.039] & $-$0.051\newline [$-$0.075, $-$0.025] & $-$0.204\newline [$-$0.281, $-$0.120] & $-$0.284\newline [$-$0.354, $-$0.208]\\
\multicolumn{2}{@{}>{\scriptsize\raggedright\arraybackslash}p{5.41cm}}{Layers with the null-centred CI $<0$ (scGPT higher) / $>0$} & 7 / 0 & 12 / 0 & 8 / 0 & 10 / 0 & 12 / 0\\
\multicolumn{7}{@{}l}{\scriptsize\textit{H3: within-TF stratified edge AUROC (two-way TF $\times$ target bootstrap intervals)}}\\
\multicolumn{2}{@{}>{\scriptsize\raggedright\arraybackslash}p{5.41cm}}{Word2Vec [95\% CI]} & 0.512\newline [0.481, 0.549] & 0.534\newline [0.498, 0.570] & 0.528\newline [0.476, 0.582] & 0.603\newline [0.533, 0.696] & 0.577\newline [0.470, 0.695]\\
\multicolumn{2}{@{}>{\scriptsize\raggedright\arraybackslash}p{5.41cm}}{scGPT, 12 layers (best layer)} & 0.516--0.542 ($L_{1}$) & 0.518--0.538 ($L_{1}$) & 0.524--0.554 ($L_{0}$) & 0.577--0.638 ($L_{0}$) & 0.601--0.652 ($L_{1}$)\\
\multicolumn{2}{@{}>{\scriptsize\raggedright\arraybackslash}p{5.41cm}}{Word2Vec minus scGPT at its best layer} & $-$0.030\newline [$-$0.074, 0.011] & $-$0.005\newline [$-$0.051, 0.036] & $-$0.026\newline [$-$0.101, 0.029] & $-$0.035\newline [$-$0.100, 0.033] & $-$0.076\newline [$-$0.174, 0.021]\\
\multicolumn{2}{@{}>{\scriptsize\raggedright\arraybackslash}p{5.41cm}}{Layers with the difference's CI $<0$ / $>0$} & 0 / 0 & 0 / 0 & 0 / 0 & 0 / 0 & 0 / 0\\
\multicolumn{7}{@{}l}{\scriptsize\textit{H1a: assortativity of $\SV{1}$--$\SV{10}$ on the STRING network (descriptive; directions not index-matched)}}\\
\multicolumn{2}{@{}>{\scriptsize\raggedright\arraybackslash}p{5.41cm}}{Word2Vec: directions with max-$T$ $p<0.05$ in all 5 seeds (largest mean $r$)} & 10/10 (0.274) & 10/10 (0.295) & 10/10 (0.265) & 10/10 (0.620) & 9/10 (0.693)\\
\multicolumn{2}{@{}>{\scriptsize\raggedright\arraybackslash}p{5.41cm}}{scGPT: direction--layer tests with max-$T$ $p<0.05$ (largest $r$)} & 120/120 (0.435) & 120/120 (0.407) & 120/120 (0.404) & 118/120 (0.566) & 120/120 (0.702)\\
\multicolumn{2}{@{}>{\scriptsize\raggedright\arraybackslash}p{5.41cm}}{Largest seed SD of Word2Vec (H1b, H2, H3)} & 0.003 & 0.003 & 0.004 & 0.008 & 0.006\\
\midrule
\multicolumn{7}{@{}p{16.4cm}@{}}{\footnotesize\textbf{B. Saturation: full-corpus Word2Vec, seed 0, after 10 and 40 epochs}}\\
\multicolumn{3}{@{}>{\scriptsize\raggedright\arraybackslash}p{7.72cm}}{Quantity} & lung, 10 epochs & lung, 40 epochs & immune, 10 epochs & immune, 40 epochs\\
\cmidrule(lr){1-7}
\multicolumn{3}{@{}>{\scriptsize\raggedright\arraybackslash}p{7.72cm}}{Pair updates} & $4.5\times10^{8}$ & $1.8\times10^{9}$ & $2.9\times10^{8}$ & $1.2\times10^{9}$\\
\multicolumn{3}{@{}>{\scriptsize\raggedright\arraybackslash}p{7.72cm}}{H1b [95\% CI, delete-$d$]} & 0.606\newline [0.589, 0.622] & 0.612\newline [0.596, 0.628] & 0.705\newline [0.657, 0.752] & 0.702\newline [0.656, 0.748]\\
\multicolumn{3}{@{}>{\scriptsize\raggedright\arraybackslash}p{7.72cm}}{H2 [95\% CI]} & 0.630\newline [0.605, 0.654] & 0.664\newline [0.639, 0.687] & 0.520\newline [0.461, 0.580] & 0.515\newline [0.454, 0.577]\\
\multicolumn{3}{@{}>{\scriptsize\raggedright\arraybackslash}p{7.72cm}}{H2 permutation $p$ (null mean)} & $5.0\times10^{-4}$ (0.494) & $5.0\times10^{-4}$ (0.493) & 0.16 (0.485) & 0.19 (0.486)\\
\multicolumn{3}{@{}>{\scriptsize\raggedright\arraybackslash}p{7.72cm}}{H3 [95\% CI]} & 0.512\newline [0.481, 0.547] & 0.511\newline [0.476, 0.551] & 0.608\newline [0.537, 0.699] & 0.614\newline [0.546, 0.707]\\
\multicolumn{3}{@{}>{\scriptsize\raggedright\arraybackslash}p{7.72cm}}{H1a: largest $r$ (directions with max-$T$ $p<0.05$)} & 0.277 (10/10) & 0.282 (10/10) & 0.620 (10/10) & 0.616 (10/10)\\
\multicolumn{3}{@{}>{\scriptsize\raggedright\arraybackslash}p{7.72cm}}{Linear CKA between the 10- and 40-epoch models} & \multicolumn{2}{>{\scriptsize}c}{0.944} & \multicolumn{2}{>{\scriptsize}c@{}}{0.945}\\
\multicolumn{3}{@{}>{\scriptsize\raggedright\arraybackslash}p{7.72cm}}{H2, Word2Vec minus scGPT at its best layer, null-centred} & $^{a}$ & $+$0.000 ($L_{5}$)\newline [$-$0.024, 0.022] & $^{a}$ & $-$0.202 ($L_{2}$)\newline [$-$0.283, $-$0.116]\\
\multicolumn{3}{@{}>{\scriptsize\raggedright\arraybackslash}p{7.72cm}}{H2, Word2Vec minus scGPT at $L_{11}$, null-centred} & $^{a}$ & $+$0.039\newline [0.014, 0.063] & $^{a}$ & $-$0.049\newline [$-$0.119, 0.024]\\
\multicolumn{3}{@{}>{\scriptsize\raggedright\arraybackslash}p{7.72cm}}{H2, layers whose null-centred CI includes 0} & $^{a}$ & 11/12 & $^{a}$ & 3/12\\
\midrule
\multicolumn{7}{@{}p{16.4cm}@{}}{\footnotesize\textbf{C. Window calibration: Word2Vec trained on the extraction cells only (256 cells), seed 0}}\\
Quantity & lung: window 5, 10 epochs & lung: window 5, matched updates & lung: window 512, 10 epochs & immune: window 5, 10 epochs & immune: window 5, matched updates & immune: window 512, 10 epochs\\
\cmidrule(lr){1-7}
Epochs; pair updates & 10; $7.8\times10^{6}$ & 571; $4.5\times10^{8}$ & 10; $4.5\times10^{8}$ & 10; $3.6\times10^{6}$ & 365; $1.3\times10^{8}$ & 10; $1.3\times10^{8}$\\
H1b [95\% CI, delete-$d$] & 0.551\newline [0.535, 0.567] & 0.571\newline [0.555, 0.586] & 0.571\newline [0.555, 0.588] & 0.616\newline [0.569, 0.663] & 0.689\newline [0.640, 0.735] & 0.656\newline [0.605, 0.707]\\
H2 (permutation $p$) & 0.487\newline (0.64) & 0.611\newline ($5.0\times10^{-4}$) & 0.594\newline ($5.0\times10^{-4}$) & 0.443\newline (0.29) & 0.526\newline (0.12) & 0.492\newline (0.43)\\
H3 [95\% CI] & 0.502\newline [0.476, 0.528] & 0.516\newline [0.490, 0.545] & 0.520\newline [0.495, 0.549] & 0.512\newline [0.445, 0.598] & 0.592\newline [0.523, 0.673] & 0.569\newline [0.498, 0.661]\\
H1a: largest $r$ (directions with max-$T$ $p<0.05$) & 0.217 (7/10) & 0.193 (9/10) & 0.159 (10/10) & 0.439 (7/10) & 0.555 (9/10) & 0.464 (10/10)\\
Linear CKA with the window-512 model & 0.262 & 0.795 & -- & 0.401 & 0.830 & --\\
\end{longtable}
\endgroup

\setcounter{table}{11}
\begingroup
\footnotesize
\setlength{\tabcolsep}{2.4pt}
\setlength{\LTpre}{6pt}\setlength{\LTpost}{6pt}
\begin{longtable}{@{}>{\scriptsize\raggedleft\arraybackslash}p{0.42cm}>{\scriptsize\raggedright\arraybackslash}p{1.3cm}>{\scriptsize\raggedright\arraybackslash}p{4.25cm}>{\scriptsize\raggedright\arraybackslash}p{0.7cm}>{\scriptsize\raggedright\arraybackslash}p{1.55cm}>{\scriptsize\raggedright\arraybackslash}p{1.95cm}>{\scriptsize\raggedleft\arraybackslash}p{1.45cm}>{\scriptsize\raggedleft\arraybackslash}p{1.45cm}>{\scriptsize\raggedright\arraybackslash}p{0.8cm}>{\scriptsize\raggedright\arraybackslash}p{1.5cm}@{}}
\caption{\textbf{The 183 hypothesis records of the automated screen (iterations 1--63).} Each record is one hypothesis of the two-agent screen as archived in the iteration's JSON file, parsed together with the iteration's code. The first 63 iterations produced these 183 records, whose headline findings Kendiukhov (2026b) (arXiv:2602.22247) reported. \emph{Panel A}, per JSON family (the 8 families the records use; Kendiukhov (2026b) regroups them into 13): $n$, records; spec., records involving a spectral decomposition; extr., records whose result string yields a primary $p$-value; $=0$, of these, $p$ stored as exactly 0 (permutation tests with $N = 50$--1,000 draws); Bonf.\ rec., records passing Bonferroni at $0.05/183$ with the recorded $p$ ($p \leq 0.05/183$, zeros kept); Bonf.\ fl., passing with the floored $p$ ($p < 0.05/183$), where a zero is replaced by its resolution floor $1/(N+1)$, with $N$ from the record's method text where stated and 1,000 otherwise; BH, passing Benjamini--Hochberg at $q = 0.05$ over the 112 extractable floored $p$-values; kid.~A, records using the kidney~A matrix (a Tabula Sapiens kidney extraction, 4,803 genes, described as immune in Kendiukhov (2026b)), with the number whose gene labels were misaligned in parentheses (mis.; iterations 27 and 29--39, where blank lines of the gene list were dropped so that all 209 labels sat on rows 0--208); imm., records using the immune extraction (one record uses both); other, neither (one Gaussian null matrix, one Geneformer token-embedding matrix); directions and bands, the five most frequent spectral directions or bands named in the family's spectral records (SV24 and SV57 are the records' own labels for $\SV{2}$--$\SV{4}$ and $\SV{5}$--$\SV{7}$; PC, principal component). The last two rows repeat the audit over all 409 records of the loop (137 JSON files; 3 records belong to an aborted duplicate run), with the Bonferroni denominator equal to the number of extractable $p$-values; this is a sensitivity analysis. \emph{Panel B}, every record in iteration order: It., iteration; ID, hypothesis identifier within the iteration ($^{\star}$, record keyed to one of the 11 headline findings in the multiplicity audit); Fam., JSON family code (Panel A); Matrix, embedding matrix loaded (3 draws, the main extraction and two further 256-cell draws; lung attention, scGPT attention scores from a lung extraction; $^{\mathrm{m}}$, misaligned gene labels); Directions recorded, spectral directions or bands named in the record (---, none); $p$ recorded, primary $p$ as extracted from the result string (---, none extractable); $p$ floored, after replacing zeros by $1/(N+1)$; Passes: R, Bonferroni at $0.05/183$ with the recorded $p$; F, Bonferroni at $0.05/183$ with the floored $p$; H, Benjamini--Hochberg at $q = 0.05$ with the floored $p$; Decision, the screen's own decision for the record. Of the 11 starred records, 8 carry R and 3 carry F; the headline-level count after flooring (5 of the 11) scores 5 of the headline findings by a $p$-value recomputed for the claim itself instead of the record's primary $p$ shown here (for the transcription-factor classifier, a Stouffer combination of its 12 per-layer $p$-values, which is invalid because the layers share genes), and 2 of these pass; replacing the $t$-approximation $p = 1.4\times10^{-24}$ of the STRING confidence gradient by its exact value, 0.0167, and setting aside the Stouffer combination leaves 3 (Table~S13, C23 and C28). Titles are truncated to about 55 characters. No intervals are reported: $p$-values are as computed by the screen and are displayed to three significant digits. Because each iteration's hypotheses were chosen after seeing earlier results, no correction over the recorded family bounds the error rate of the screen as a whole. The file also gives, for every record, its gene set and gene-row mapping, decomposition convention, reference sets, script provenance and the claims of Table~S13 it supports, and the same fields for the 226 later records. Full values: \texttt{Table\_S12\_screen\_records.csv}.}\label{tab:S12}\\
\multicolumn{10}{@{}l}{\footnotesize\textbf{A.} Records per JSON family and multiple-testing passes}\\[2pt]
\multicolumn{10}{@{}l@{}}{\scriptsize\setlength{\tabcolsep}{2.4pt}\begin{tabular}{@{}>{\raggedright\arraybackslash}p{3.3cm}rrrrrrrrrr>{\raggedright\arraybackslash}p{4.3cm}@{}}
\toprule
 & & & \multicolumn{2}{c}{$p$-values} & \multicolumn{3}{c}{passes} & \multicolumn{3}{c}{matrix used} & \\
\cmidrule(lr){4-5}\cmidrule(lr){6-8}\cmidrule(lr){9-11}
JSON family (code) & $n$ & spec. & extr. & $=0$ & Bonf.\ rec. & Bonf.\ fl. & BH & kid.~A (mis.) & imm. & other & directions and bands involved (records) \\
\midrule
\texttt{module\_structure} (MOD) & 52 & 37 & 33 & 5 & 12 & 7 & 21 & 35 (7) & 17 & 0 & $\SV{2}$--$\SV{4}$ (16), $\SV{5}$--$\SV{7}$ (16), $\SV{2}$ (7), $\SV{1}$ (6), $\SV{2}$--$\SV{7}$ (4), \ldots{} \\
\texttt{manifold\_distance} (MAN) & 42 & 11 & 29 & 4 & 10 & 6 & 19 & 22 (10) & 20 & 0 & $\SV{5}$--$\SV{7}$ (9), $\SV{2}$--$\SV{4}$ (3), $\SV{5}$ (2), $\SV{7}$ (2), $\SV{2}$ (1), \ldots{} \\
\texttt{intrinsic\_dimensionality} (IDM) & 41 & 32 & 24 & 4 & 10 & 6 & 16 & 20 (6) & 21 & 0 & $\SV{1}$ (13), $\SV{2}$ (6), $\SV{1}$/$\SV{2}$ (5), PC1 (4), $\SV{2}$/$\SV{3}$ (2), \ldots{} \\
\texttt{graph\_topology} (GRT) & 16 & 11 & 9 & 1 & 3 & 2 & 4 & 14 (3) & 2 & 0 & $\SV{2}$ (3), $\SV{3}$ (2), $\SV{2}$+$\SV{3}$+$\SV{4}$ (2), $\SV{1}$ (1), $\SV{4}$ (1) \\
\texttt{null\_sensitivity} (NUL) & 15 & 6 & 7 & 2 & 3 & 1 & 5 & 13 (7) & 1 & 1 & $\SV{2}$ (3), $\SV{3}$/$\SV{4}$ (1), $\SV{3}$ (1), $\SV{4}$ (1), PC1 (1) \\
\texttt{topology\_stability} (TST) & 8 & 6 & 6 & 2 & 2 & 0 & 2 & 6 (3) & 2 & 0 & $\SV{1}$ (2), $\SV{5}$--$\SV{7}$ (2), SV57 (2), $\SV{2}$--$\SV{4}$ (2), $\SV{2}$/$\SV{3}$ (1), \ldots{} \\
\texttt{cross\_model\_alignment} (CMA) & 5 & 3 & 0 & 0 & 0 & 0 & 0 & 3 (0) & 2 & 1 & $\SV{5}$--$\SV{7}$ (1) \\
\texttt{persistent\_homology} (PHO) & 4 & 2 & 4 & 0 & 3 & 3 & 3 & 2 (0) & 2 & 0 & $\SV{2}$--$\SV{4}$ (2) \\
\midrule
All 183 records (iterations 1--63) & 183 & 108 & 112 & 18 & 43 & 25 & 70 & 115 (36) & 67 & 2 & Bonferroni threshold $0.05/183 = 2.73{\times}10^{-4}$ \\
All 409 records (iterations 1--138) & 409 & --- & 130 & 18 & --- & 28 & 79 & --- & --- & --- & Bonferroni at $0.05/130$ \\
\quad excl.\ aborted duplicate run & 406 & --- & 127 & 18 & --- & 27 & 76 & --- & --- & --- & Bonferroni at $0.05/127$ \\
\bottomrule
\end{tabular}}\\[4pt]
\multicolumn{10}{@{}l}{\footnotesize\textbf{B.} The 183 records}\\
\toprule
It. & ID & Hypothesis (short title) & Fam. & Matrix & Directions recorded & $p$ recorded & $p$ floored & Passes & Decision \\
\midrule
\endfirsthead
\multicolumn{10}{@{}l}{\footnotesize\textit{Table~\thetable{} (continued)}}\\
\toprule
It. & ID & Hypothesis (short title) & Fam. & Matrix & Directions recorded & $p$ recorded & $p$ floored & Passes & Decision \\
\midrule
\endhead
\midrule
\multicolumn{10}{@{}r@{}}{\scriptsize\textit{continued on next page}}\\
\endfoot
\bottomrule
\endlastfoot
2 & H01\_\allowbreak{}graph\_\allowbreak{}topology\_\allowbreak{}cc & kNN Clustering Coefficient vs Feature-Shuffle Null & GRT & kidney~A (3 draws) & --- & --- & --- &  & promising \\
2 & H02\_\allowbreak{}graph\_\allowbreak{}topology\_\allowbreak{}tr & kNN Transitivity vs Feature-Shuffle Null & GRT & kidney~A (3 draws) & --- & --- & --- &  & negative \\
4 & H01 & TwoNN intrinsic dimensionality per layer & IDM & kidney~A & --- & --- & --- &  & promising \\
4 & H02 & TRRUST co-target distance clustering (random-group\ldots{} & MOD & kidney~A & --- & --- & --- &  & inconclusive \\
4 & H03 & Cross-layer linear CKA alignment & CMA & kidney~A & --- & --- & --- &  & promising \\
5 & H01 & GO term embedding clustering across layers & MOD & kidney~A & --- & --- & --- &  & inconclusive \\
5 & H02 & Per-gene residual drift + GO enrichment & MOD & kidney~A & --- & 0.0044 & 0.0044 & H & promising \\
5 & H03 & Effective rank per layer (spectral cross-check vs\ldots{} & IDM & kidney~A & --- & 0.002 & 0.002 & H & promising \\
6 & H01 & Full-vocabulary residual drift + GO enrichment & MOD & kidney~A & --- & 0.0078 & 0.0078 & H & promising \\
6 & H02$^{\star}$ & SVD biology of near-rank-1 layer-11 subspace & IDM & kidney~A & $\SV{1}$, $\SV{1}$/$\SV{2}$ & $3{\times}10^{-4}$ & $3{\times}10^{-4}$ & H & promising \\
6 & H03 & Effective rank curvature and compression breakpoint & IDM & kidney~A & --- & 0.0021 & 0.0021 & H & promising \\
7 & H01 & $\SV{1}$ feature-shuffle and gene-label-shuffle null & IDM & kidney~A & $\SV{1}$ & 0.004 & 0.004 & H & promising \\
7 & H02 & Layer-wise SVD trajectory (12 layers) & IDM & kidney~A & $\SV{1}$/$\SV{2}$, $\SV{1}$ & --- & --- &  & promising \\
7 & H03 & Drift gene-label-shuffle null + $\SV{1}$ negative-pole GO\ldots{} & MOD & kidney~A & $\SV{1}$ & 0.124 & 0.124 &  & neutral \\
8 & H01 & $\SV{2}$/$\SV{3}$ axes GO enrichment + gene-label null at\ldots{} & IDM & kidney~A & $\SV{2}$/$\SV{3}$, $\SV{2}$, $\SV{3}$ & 0 & $9.99{\times}10^{-4}$ & R H & promising \\
8 & H02$^{\star}$ & Layer-specific compartment transients: mitochondrion\ldots{} & IDM & kidney~A & $\SV{1}$ & 0 & $9.99{\times}10^{-4}$ & R H & promising \\
8 & H03 & Signal peptide proxy vs $\SV{1}$ top-pole (layer-11) & MOD & kidney~A & $\SV{1}$ & 0.002 & 0.002 & H & promising \\
9 & H01 & $\SV{2}$ bottom-pole extracellular vesicle null (N=1000) & IDM & kidney~A & $\SV{2}$ & 0 & $9.99{\times}10^{-4}$ & R H & promising \\
9 & H02 & 12-layer x 8-compartment systematic $\SV{1}$ scan & MOD & kidney~A & $\SV{1}$ & --- & --- &  & promising \\
9 & H03 & Cross-layer $\SV{1}$ rank stability (Spearman) & TST & kidney~A & $\SV{1}$ & 0 & $9.99{\times}10^{-4}$ & R H & promising \\
10 & H01 & $\SV{2}$/$\SV{3}$ 12-layer $\times$ 9-compartment scan & IDM & kidney~A & $\SV{2}$/$\SV{3}$, $\SV{2}$, $\SV{3}$ & 0.05 & 0.05 &  & promising \\
10 & H02$^{\star}$ & TRRUST TF-target co-pole test in SVD geometry & MOD & kidney~A & $\SV{1}$, $\SV{2}$ & 0 & $9.99{\times}10^{-4}$ & R H & promising \\
10 & H03 & $\SV{1}$/$\SV{2}$/$\SV{3}$ spectral ratio profile + annotation\ldots{} & IDM & kidney~A & $\SV{1}$/$\SV{2}$/$\SV{3}$, $\SV{1}$/$\SV{2}$, $\SV{1}$, $\SV{2}$ & --- & --- &  & promising \\
11 & H01$^{\star}$ & TRRUST co-pole x 12 layers + Activation/Repression\ldots{} & MOD & kidney~A & $\SV{2}$ & 0 & $9.99{\times}10^{-4}$ & R H & promising \\
11 & H02 & $\SV{2}$ mito pole-flip gene tracking at $L_{3}$-to-$L_{4}$ transition & IDM & kidney~A & $\SV{2}$ & --- & --- &  & neutral \\
11 & H03 & GO Biological Process enrichment in $\SV{2}$ poles (broad\ldots{} & MOD & kidney~A & $\SV{2}$ & 1 & 1 &  & negative \\
12 & H01 & Bootstrap CI on activation vs repression co-pole\ldots{} & MOD & kidney~A & --- & 0.05 & 0.05 &  & negative \\
12 & H02 & Mean pairwise $\SV{2}$ distance --- spatial concentration\ldots{} & MAN & kidney~A & $\SV{2}$ & 0.05 & 0.05 &  & promising \\
12 & H03 & STRING PPI co-pole enrichment across 12 layers & MOD & kidney~A & $\SV{2}$ & 0.05 & 0.05 &  & promising \\
13 & H01 & SV axis specificity: STRING PPI co-pole on $\SV{1}$ vs $\SV{2}$\ldots{} & GRT & kidney~A & $\SV{1}$, $\SV{2}$, $\SV{3}$ & --- & --- &  & promising \\
13 & H02 & Repression anti-pole: TRRUST repression pairs occupy\ldots{} & MOD & kidney~A & $\SV{2}$ & --- & --- &  & negative \\
13 & H03 & Hub-degree control: does PPI co-pole signal survive in\ldots{} & NUL & kidney~A & $\SV{2}$ & --- & --- &  & inconclusive \\
14 & H01 & GO enrichment of $\SV{3}$ poles across 12 layers & MOD & kidney~A & $\SV{3}$ & --- & --- &  & promising \\
14 & H02 & Layer-depth Spearman trend on $\SV{2}$/$\SV{3}$ PPI z-scores & TST & kidney~A & $\SV{2}$/$\SV{3}$, $\SV{1}$/$\SV{2}$/$\SV{3}$, $\SV{2}$, $\SV{3}$, $\SV{1}$ & 0.059 & 0.059 &  & neutral \\
14 & H03 & Expanded hub-degree control: non-hub $\SV{2}$ co-pole\ldots{} & NUL & kidney~A & $\SV{2}$ & --- & --- &  & promising \\
15 & H01 & $\SV{4}$/$\SV{5}$ PPI co-pole test & MOD & kidney~A & $\SV{4}$/$\SV{5}$, $\SV{4}$, $\SV{5}$ & --- & --- &  & promising \\
15 & H02$^{\star}$ & STRING confidence gradient: co-pole z vs score quintile & GRT & kidney~A & $\SV{2}$ & $1.4{\times}10^{-24}$ & $1.4{\times}10^{-24}$ & R F H & promising \\
15 & H03 & GO co-annotation vs PPI as driver of $\SV{2}$ geometry & MOD & kidney~A & $\SV{2}$ & --- & --- &  & promising \\
16 & H01 & $\SV{3}$ STRING confidence gradient & GRT & kidney~A & $\SV{3}$ & 0.037 & 0.037 &  & promising \\
16 & H02 & $\SV{2}$/$\SV{3}$/$\SV{4}$ axis independence test & MOD & kidney~A & $\SV{2}$/$\SV{3}$/$\SV{4}$, $\SV{2}$, $\SV{3}$, $\SV{4}$, $\SV{2}$--$\SV{3}$, $\SV{2}$--$\SV{4}$, $\SV{3}$--$\SV{4}$ & --- & --- &  & promising \\
16 & H03 & $\SV{4}$ GO biology profile & MOD & kidney~A & $\SV{4}$ & $8.2{\times}10^{-4}$ & $8.2{\times}10^{-4}$ & H & promising \\
17 & H01 & $\SV{4}$ confidence quintile gradient & GRT & kidney~A & $\SV{4}$ & 0.0374 & 0.0374 &  & promising \\
17 & H02 & Axis-dominant pair GO composition (biological\ldots{} & MOD & kidney~A & $\SV{2}$/$\SV{3}$/$\SV{4}$, $\SV{2}$--$\SV{3}$, $\SV{2}$--$\SV{4}$, $\SV{3}$--$\SV{4}$ & --- & --- &  & promising \\
17 & H03 & TRRUST signed regulation in $\SV{3}$/$\SV{4}$ & NUL & kidney~A & $\SV{3}$/$\SV{4}$, $\SV{3}$, $\SV{4}$ & --- & --- &  & negative \\
18 & H01 & Random Gaussian embedding null control for $\SV{2}$ co-pole & NUL & Gaussian null & $\SV{2}$ & --- & --- &  & negative \\
18 & H02 & Out-of-sample PPI precision@k prediction benchmark & GRT & kidney~A & $\SV{2}$ & --- & --- &  & neutral \\
18 & H03 & scGPT attention co-occurrence geometry for\ldots{} & GRT & kidney~A + lung attention & --- & 0.87 & 0.87 &  & promising \\
19 & H01 & Geneformer static token embedding cross-model alignment & CMA & kidney~A + Geneformer & --- & --- & --- &  & promising \\
19 & H02 & Multi-axis composite co-pole P@k enrichment & GRT & kidney~A & $\SV{2}$+$\SV{3}$+$\SV{4}$ & --- & --- &  & promising \\
19 & H03 & TRRUST repression vs activation in scGPT attention & NUL & kidney~A + lung attention & --- & $9.4{\times}10^{-9}$ & $9.4{\times}10^{-9}$ & R F H & promising \\
20 & H01 & Shuffle null for multi-axis composite co-pole\ldots{} & NUL & kidney~A & --- & 0 & $9.99{\times}10^{-4}$ & R H & promising \\
20 & H02 & Attention-SVD joint ROC predictor for STRING/TRRUST & GRT & kidney~A + lung attention & $\SV{2}$+$\SV{3}$+$\SV{4}$ & --- & --- &  & neutral \\
20 & H03 & Persistent homology H0 and pairwise distance geometry & PHO & kidney~A & --- & $1{\times}10^{-72}$ & $1{\times}10^{-72}$ & R F H & promising \\
21 & H01 & STRING confidence quintile x pairwise distance gradient & MAN & kidney~A & --- & 0.05 & 0.05 &  & inconclusive \\
21 & H02 & TRRUST pairs distance --- specificity control & MAN & kidney~A & --- & --- & --- &  & promising \\
21 & H03 & H1 persistence trajectory + bootstrap CIs for\ldots{} & PHO & kidney~A & --- & $1{\times}10^{-4}$ & $1{\times}10^{-4}$ & R F H & promising \\
22 & H01 & STRING continuous AUROC + TRRUST-exclusive proximity\ldots{} & MAN & kidney~A & --- & $1{\times}10^{-5}$ & $1{\times}10^{-5}$ & R F H & promising \\
22 & H02$^{\star}$ & Cell-type marker gene cluster separation & MOD & kidney~A & --- & $1{\times}10^{-6}$ & $1{\times}10^{-6}$ & R F H & promising \\
22 & H03 & Per-layer bootstrap CIs for co-polarity enrichment +\ldots{} & TST & kidney~A & $\SV{2}$+$\SV{3}$+$\SV{4}$ & 0.16 & 0.16 &  & promising \\
23 & H01$^{\star}$ & Cell-type marker separation --- expanded panel +\ldots{} & MOD & kidney~A & --- & 0.004 & 0.004 & H & promising \\
23 & H02 & GO Biological Process proximity --- 3rd independent\ldots{} & MOD & kidney~A & --- & $1{\times}10^{-19}$ & $1{\times}10^{-19}$ & R F H & promising \\
23 & H03 & TRRUST Activation vs Repression proximity asymmetry & MAN & kidney~A & --- & 0.002 & 0.002 & H & promising \\
24 & H01 & Dorothea confidence stratification & MAN & kidney~A & --- & 0 & $9.99{\times}10^{-4}$ & R H & promising \\
24 & H02 & TF activation hub centrality & MOD & kidney~A & --- & 0.41 & 0.41 &  & negative \\
24 & H03 & GO ontology comparison BP vs MF vs CC & MOD & kidney~A & --- & $9{\times}10^{-56}$ & $9{\times}10^{-56}$ & R F H & promising \\
25 & H01 & Multi-predictor joint model (STRING + Dorothea + GO CC\ldots{} & MAN & kidney~A & --- & --- & --- &  & promising \\
25 & H02 & Layer-resolved encoding timeline across 5 biological\ldots{} & MAN & kidney~A & --- & --- & --- &  & promising \\
25 & H03 & Chromosomal proximity negative control & NUL & kidney~A & --- & 0.096 & 0.096 &  & negative \\
26 & H01 & Immune gene family clustering AUROC & MOD & kidney~A & --- & --- & --- &  & promising \\
26 & H02 & Dorothea confidence-tier proximity gradient & MAN & kidney~A & --- & --- & --- &  & promising \\
26 & H03 & Intrinsic dimensionality progression (participation\ldots{} & IDM & kidney~A & --- & $1{\times}10^{-8}$ & $1{\times}10^{-8}$ & R F H & promising \\
27 & H01 & PC1 Biological Identity at $L_{11}$ & IDM & kidney~A$^{\mathrm{m}}$ & PC1 & 0.42 & 0.42 &  & negative \\
27 & H02 & TRRUST Activator/Repressor Geometric Polarity & MOD & kidney~A$^{\mathrm{m}}$ & PC1 & 0.178 & 0.178 &  & negative \\
27 & H03 & Dorothea Threshold Sensitivity + PR Elbow Alignment & NUL & kidney~A$^{\mathrm{m}}$ & --- & --- & --- &  & negative \\
28 & H01 & Hub Gene Embedding Centrality & MAN & kidney~A & --- & 0.00135 & 0.00135 & H & promising \\
28 & H02 & PC1 T-cell vs APC Cell-type Axis & IDM & kidney~A & PC1 & 0.12 & 0.12 &  & negative \\
28 & H03 & kNN Graph Spectral Gap Across Layers & GRT & kidney~A & --- & $1.3{\times}10^{-10}$ & $1.3{\times}10^{-10}$ & R F H & promising \\
29 & H01 & kNN Connected Components Biological Enrichment at $L_{11}$ & GRT & kidney~A$^{\mathrm{m}}$ & --- & 0.023 & 0.023 & H & promising \\
29 & H02 & Gene Trajectory Clustering by Distance-to-Centroid\ldots{} & MAN & kidney~A$^{\mathrm{m}}$ & --- & 0.352 & 0.352 &  & promising \\
29 & H03 & Spectral Gap k-Robustness Sweep (k=5--30) & TST & kidney~A$^{\mathrm{m}}$ & --- & --- & --- &  & promising \\
30 & H01 & GO and TRRUST enrichment of 14 diverging genes & MOD & kidney~A$^{\mathrm{m}}$ & --- & $2.34{\times}10^{-11}$ & $2.34{\times}10^{-11}$ & R F H & positive \\
30 & H02 & Layer-of-bifurcation anatomy (inter/intra ratio) & TST & kidney~A$^{\mathrm{m}}$ & --- & 0 & $9.99{\times}10^{-4}$ & R H & negative \\
30 & H03 & Trajectory slope vs STRING degree (continuous) & GRT & kidney~A$^{\mathrm{m}}$ & --- & 0.683 & 0.683 &  & negative \\
30 & H04 & OOV gene discovery in scGPT named gene set & NUL & kidney~A$^{\mathrm{m}}$ & --- & --- & --- &  & promising \\
31 & H01 & kNN Spectral Gap Monotonic Decrease (195 in-vocab\ldots{} & GRT & kidney~A$^{\mathrm{m}}$ & --- & 0 & $9.99{\times}10^{-4}$ & R H & promising \\
31 & H02 & STRING Score $\to$ Embedding Distance (195 in-vocab genes\ldots{} & MAN & kidney~A$^{\mathrm{m}}$ & --- & --- & --- &  & negative \\
31 & H03 & Participation Ratio (Intrinsic Dim) Collapse Across\ldots{} & IDM & kidney~A$^{\mathrm{m}}$ & --- & 0 & $9.99{\times}10^{-4}$ & R H & promising \\
32 & H01 & Community Detection at $L_{11}$ $\to$ Biological Annotation & MOD & kidney~A$^{\mathrm{m}}$ & --- & 0.09 & 0.09 &  & promising \\
32 & H02 & Dorothea Regulatory Pairs $\to$ Embedding Distance Decay\ldots{} & MAN & kidney~A$^{\mathrm{m}}$ & --- & $4{\times}10^{-4}$ & $4{\times}10^{-4}$ & H & promising \\
32 & H03 & PC1 at $L_{11}$ Separates TFs from B-cell Markers\ldots{} & IDM & kidney~A$^{\mathrm{m}}$ & PC1 & 0.041 & 0.041 &  & inconclusive \\
33 & H01 & $L_{11}$ Community Differential Analysis (B-cell vs\ldots{} & MOD & kidney~A$^{\mathrm{m}}$ & PC1 & $1{\times}10^{-9}$ & $1{\times}10^{-9}$ & R F H & promising \\
33 & H02 & Community Partition Stability: k-sweep and Layer Sweep & TST & kidney~A$^{\mathrm{m}}$ & --- & 0.167 & 0.167 &  & promising \\
33 & H03 & PC1 Polarity Expanded Gene Sets --- B-cell Stable Across\ldots{} & IDM & kidney~A$^{\mathrm{m}}$ & PC1 & 0.001 & 0.001 & H & promising \\
34 & H01 & Layer-wise B-cell community purity emergence & MOD & kidney~A$^{\mathrm{m}}$ & PC1 & 0.804 & 0.804 &  & inconclusive \\
34 & H02 & PC2/PC3 cell-type axes at $L_{11}$ & IDM & kidney~A$^{\mathrm{m}}$ & PC2/PC3, PC2, PC3 & 0.168 & 0.168 &  & negative \\
34 & H03 & Expression-level confound test for B-cell PC1 signal & NUL & kidney~A$^{\mathrm{m}}$ & PC1 & 0 & $9.99{\times}10^{-4}$ & R H & promising \\
35 & H01 & TRRUST co-regulation + GO co-annotation vs embedding\ldots{} & MOD & kidney~A$^{\mathrm{m}}$ & --- & $1{\times}10^{-4}$ & $1{\times}10^{-4}$ & R F H & negative \\
35 & H02 & B-cell marker kNN precision@10 profile across layers & MAN & kidney~A$^{\mathrm{m}}$ & --- & 0 & $9.99{\times}10^{-4}$ & R H & promising \\
35 & H03 & B-cell vs non-B-cell centroid separation trajectory\ldots{} & IDM & kidney~A$^{\mathrm{m}}$ & --- & 0.088 & 0.088 &  & promising \\
36 & H01 & Multi-cell-type kNN precision@10 specificity panel & MAN & kidney~A$^{\mathrm{m}}$ & --- & --- & --- &  & promising \\
36 & H02 & Gene-name permutation null for B-cell kNN precision & NUL & kidney~A$^{\mathrm{m}}$ & --- & 0.005 & 0.005 & H & promising \\
36 & H03 & B-cell neighborhood functional characterization & MOD & kidney~A$^{\mathrm{m}}$ & --- & --- & --- &  & promising \\
37 & H01 & Extended cell-type panel kNN precision@10 + centroid\ldots{} & MAN & kidney~A$^{\mathrm{m}}$ & --- & 0 & $9.99{\times}10^{-4}$ & R H & promising \\
37 & H02 & Held-out B-cell marker generalization (germinal center\ldots{} & MAN & kidney~A$^{\mathrm{m}}$ & --- & --- & --- &  & promising \\
37 & H03 & T-cell permutation null + B-cell neighborhood TF\ldots{} & NUL & kidney~A$^{\mathrm{m}}$ & --- & 0.96 & 0.96 &  & promising \\
38 & H01 & GC-TF vs Plasma-TF centroid proximity divergence\ldots{} & MAN & kidney~A$^{\mathrm{m}}$ & --- & 0.024 & 0.024 & H & promising \\
38 & H02 & B-cell centroid directional drift alignment with\ldots{} & MAN & kidney~A$^{\mathrm{m}}$ & --- & $1{\times}10^{-4}$ & $1{\times}10^{-4}$ & R F H & inconclusive \\
38 & H03 & Master B-cell TF proximity + NK/myeloid specificity\ldots{} & NUL & kidney~A$^{\mathrm{m}}$ & --- & --- & --- &  & inconclusive \\
39 & H01 & Drift target identification --- B-cell centroid at $L_{11}$ & MAN & kidney~A$^{\mathrm{m}}$ & --- & --- & --- &  & promising \\
39 & H02 & Extended plasma panel + BCL6/PAX5 GC screen\ldots{} & MAN & immune & --- & --- & --- &  & mixed \\
39 & H03 & LOO ablation of B-cell anchor genes on GC-TF proximity\ldots{} & NUL & kidney~A$^{\mathrm{m}}$ & --- & --- & --- &  & promising \\
40 & H01 & Full 12-layer GC attractor onset scan + IRF4\ldots{} & MAN & immune & --- & $2{\times}10^{-6}$ & $2{\times}10^{-6}$ & R F H & promising \\
40 & H02 & Geneformer cross-model B-cell precision@10 (input\ldots{} & CMA & Geneformer & --- & --- & --- &  & negative \\
40 & H03 & GC-plasma subspace principal angle + CD19+BLK minimal\ldots{} & MAN & immune & --- & $1{\times}10^{-4}$ & $1{\times}10^{-4}$ & R F H & promising \\
41 & H01 & Multi-lineage attractor screen (T-cell + myeloid) & MAN & immune & --- & $1{\times}10^{-6}$ & $1{\times}10^{-6}$ & R F H & promising \\
41 & H02 & BCL6 and PAX5 neighborhood characterization & MAN & immune & --- & --- & --- &  & promising \\
41 & H03 & TwoNN intrinsic dimensionality by layer and lineage & IDM & immune & --- & $1{\times}10^{-4}$ & $1{\times}10^{-4}$ & R F H & promising \\
42 & H01 & TwoNN change-point analysis at $L_{3}$ (B-cell ID\ldots{} & IDM & immune & --- & 0.386 & 0.386 &  & promising \\
42 & H02 & Metabolic isolation generalization --- BCL6 specificity\ldots{} & MAN & immune & --- & --- & --- &  & promising \\
42 & H03$^{\star}$ & BATF/BACH2 convergence trajectory toward B-cell context & MAN & immune & --- & $1{\times}10^{-4}$ & $1{\times}10^{-4}$ & R F H & promising \\
43 & H01 & GC repression circuit PRDM1/IRF4 anti-convergence test & MAN & immune & --- & --- & --- &  & negative \\
43 & H02 & BCL6 divergence from B-cell centroid & MAN & immune & --- & 0 & $9.99{\times}10^{-4}$ & R H & negative \\
43 & H03 & T-cell subtype ID stratification at $L_{7}$ breakpoint & IDM & immune & --- & --- & --- &  & promising \\
44 & H01 & Cross-lineage ID compression law & IDM & immune & --- & --- & --- &  & negative \\
44 & H02 & TCR activation circuit attractor test & MAN & immune & --- & 0.275 & 0.275 &  & inconclusive \\
44 & H03 & BCL6 rank B-cell vs metabolic co-movement & MAN & immune & --- & 0.095 & 0.095 &  & negative \\
45 & H01 & kNN lineage community purity at $L_{2}$ vs $L_{11}$ & GRT & immune & --- & 0.958 & 0.958 &  & negative \\
45 & H02 & Full-manifold TwoNN intrinsic dimensionality across 12\ldots{} & IDM & immune & --- & --- & --- &  & promising \\
45 & H03 & TCR signaling circuit convergence to T-cell identity\ldots{} & MAN & immune & --- & 0.32 & 0.32 &  & inconclusive \\
46 & H01 & Cross-seed ID compression replication + bootstrap CI & IDM & immune (3 draws) & --- & --- & --- &  & promising \\
46 & H02 & Singular value spectrum concentration across layers & IDM & immune & --- & --- & --- &  & promising \\
46 & H03 & Lineage centroid orthogonality vs null & GRT & immune & --- & --- & --- &  & negative \\
47 & H01 & Zero-norm validity check: TwoNN and SVD on\ldots{} & IDM & immune & --- & --- & --- &  & promising \\
47 & H02 & $\SV{1}$ identity: what does the dominant direction encode? & IDM & immune & $\SV{1}$, $\SV{1}$/$\SV{2}$, $\SV{2}$ & --- & --- &  & negative \\
47 & H03 & ID compression breakpoint detection across layers & IDM & immune & --- & --- & --- &  & negative \\
48 & H01 & Gene annotation via edge dataset + $\SV{1}$/$\SV{2}$ identity & IDM & immune & $\SV{1}$/$\SV{2}$ & --- & --- &  & neutral \\
48 & H02 & Full spectral decay curve across 12 layers & IDM & immune & $\SV{1}$ & --- & --- &  & promising \\
48 & H03 & Wasserstein/transport distance between consecutive\ldots{} & MAN & immune & --- & 0.15 & 0.15 &  & inconclusive \\
49 & H01 & $\SV{1}$ gene loading direction stability across layers & IDM & immune & $\SV{1}$ & --- & --- &  & promising \\
49 & H02 & SWD vs feature-shuffle null comparison & NUL & immune & --- & 0.006 & 0.006 & H & promising \\
49 & H03 & $\SV{2}$--$\SV{4}$ subspace encodes TRRUST TF-target regulatory\ldots{} & MOD & immune & $\SV{2}$--$\SV{4}$, $\SV{1}$, $\SV{1}$--$\SV{4}$ & $1{\times}10^{-4}$ & $1{\times}10^{-4}$ & R F H & promising \\
50 & H01 & $\SV{2}$--$\SV{4}$ co-expression specificity control & MOD & immune & $\SV{2}$--$\SV{4}$ & --- & --- &  & negative \\
50 & H02 & $\SV{1}$ identity: expression correlation across layers & IDM & immune & $\SV{1}$ & --- & --- &  & promising \\
50 & H03 & $\SV{5}$--$\SV{10}$ spectral axis biological screen (TF/target\ldots{} & MOD & immune & $\SV{5}$--$\SV{10}$, $\SV{1}$, $\SV{2}$--$\SV{4}$, $\SV{5}$--$\SV{7}$, $\SV{8}$--$\SV{10}$ & 0.05 & 0.05 &  & promising \\
51 & H01 & Residual TF signal after co-expression regression in\ldots{} & MOD & immune & $\SV{2}$--$\SV{4}$ & 0.0016 & 0.0016 & H & negative \\
51 & H02 & $\SV{1}$-high genes are depleted of TFs at both $L_{0}$ and $L_{11}$ & IDM & immune & $\SV{1}$ & $1{\times}10^{-4}$ & $1{\times}10^{-4}$ & R F H & promising \\
51 & H03 & 0-dim PH cluster topology for circuit genes in $\SV{2}$--$\SV{4}$\ldots{} & PHO & immune & $\SV{2}$--$\SV{4}$ & $1{\times}10^{-6}$ & $1{\times}10^{-6}$ & R F H & inconclusive \\
52 & H01 & Housekeeping gene enrichment in $\SV{1}$-high & IDM & immune & $\SV{1}$ & --- & --- &  & inconclusive \\
52 & H02 & H1 Betti loops on circuit genes at $L_{8}$ & PHO & immune & $\SV{2}$--$\SV{4}$ & 0.953 & 0.953 &  & negative \\
52 & H03 & $\SV{5}$--$\SV{7}$ regulatory signal after co-expression\ldots{} & MOD & immune & $\SV{5}$--$\SV{7}$, $\SV{2}$--$\SV{4}$ & 0.002 & 0.002 & H & promising \\
53 & H01 & Bootstrap + spectral decay scan of SV regulatory\ldots{} & MOD & immune & $\SV{2}$--$\SV{4}$, $\SV{5}$--$\SV{7}$, $\SV{8}$--$\SV{10}$, $\SV{11}$--$\SV{14}$, $\SV{15}$--$\SV{20}$, $\SV{21}$--$\SV{30}$, $\SV{8}$--$\SV{14}$ & --- & --- &  & promising \\
53 & H02 & TF$\to$Target directional asymmetry in $\SV{5}$--$\SV{7}$ & MAN & immune & $\SV{5}$--$\SV{7}$, $\SV{5}$, $\SV{6}$, $\SV{7}$ & $8{\times}10^{-4}$ & $8{\times}10^{-4}$ & H & promising \\
53 & H03 & $\SV{1}$ loading vs TRRUST network degree & IDM & immune & $\SV{1}$ & 0.001 & 0.001 & H & neutral \\
54 & H01 & LR classifier: TF vs target-only from $\SV{5}$--$\SV{7}$ $L_{0}$ coords & MOD & immune & $\SV{5}$--$\SV{7}$, $\SV{2}$--$\SV{4}$ & 0 & 0.0099 & R H & promising \\
54 & H02 & BFS hierarchy depth vs $\SV{5}$--$\SV{7}$ dominant axis ($L_{0}$) & MAN & immune & $\SV{5}$--$\SV{7}$, $\SV{5}$, $\SV{7}$ & 0.0044 & 0.0044 & H & neutral \\
54 & H03 & Layer-resolved TF$\to$target directionality trajectory\ldots{} & MAN & immune & $\SV{5}$--$\SV{7}$ & 0.05 & 0.05 &  & promising \\
55 & H01 & Cross-seed replication: TF vs target-only AUROC\ldots{} & MOD & immune (3 draws) & $\SV{5}$--$\SV{7}$ & 0.0099 & 0.0099 & H & promising \\
55 & H02 & Layer-resolved AUROC: $\SV{5}$--$\SV{7}$ vs $\SV{2}$--$\SV{4}$ across all 12\ldots{} & MOD & immune & $\SV{5}$--$\SV{7}$, $\SV{2}$--$\SV{4}$ & --- & --- &  & promising \\
55 & H03 & Directionality trajectory $\SV{2}$--$\SV{4}$ vs $\SV{5}$--$\SV{7}$ (opposite\ldots{} & MAN & immune & $\SV{2}$--$\SV{4}$, $\SV{5}$--$\SV{7}$ & 0.001 & 0.001 & H & promising \\
56 & H01$^{\star}$ & Joint $\SV{2}$--$\SV{7}$ (6D) Layer-Stable TF/Target Classifier & MOD & immune & $\SV{2}$--$\SV{7}$, $\SV{2}$--$\SV{4}$, $\SV{5}$--$\SV{7}$, SV24, SV57 & 0 & 0.0099 & R H & promising \\
56 & H02 & $L_{9}$ Directionality Crossover Cross-Seed Replication +\ldots{} & MAN & immune (3 draws) & $\SV{2}$--$\SV{4}$, $\SV{5}$--$\SV{7}$, SV57, SV24 & --- & --- &  & inconclusive \\
56 & H03 & $\SV{5}$--$\SV{7}$ Subspace Principal Angle vs AUROC Trajectory & TST & immune & $\SV{5}$--$\SV{7}$, SV57, $\SV{2}$--$\SV{4}$ & 0.417 & 0.417 &  & negative \\
57 & H01$^{\star}$ & Joint 6D ($\SV{2}$--$\SV{7}$) TF/target classifier cross-seed\ldots{} & MOD & immune (3 draws) & $\SV{2}$--$\SV{7}$, $\SV{2}$--$\SV{4}$, $\SV{5}$--$\SV{7}$ & 0 & 0.0196 & R H & promising \\
57 & H02 & Seed43 SV basis alignment via principal angles & TST & immune (3 draws) & $\SV{2}$--$\SV{4}$, $\SV{5}$--$\SV{7}$, SV57 & --- & --- &  & negative \\
57 & H03 & SV energy fraction as AUROC proxy & IDM & immune & $\SV{1}$, $\SV{2}$--$\SV{4}$, $\SV{5}$--$\SV{7}$, $\SV{2}$--$\SV{7}$ & $1{\times}10^{-5}$ & $1{\times}10^{-5}$ & R F H & neutral \\
58 & H01 & TF boundary gene identity and family enrichment at\ldots{} & MOD & immune (3 draws) & $\SV{2}$--$\SV{7}$ & --- & --- &  & promising \\
58 & H02 & TRRUST graph Laplacian spectral alignment with $\SV{5}$--$\SV{7}$ & MOD & immune & $\SV{5}$--$\SV{7}$, $\SV{2}$--$\SV{4}$ & --- & --- &  & neutral \\
58 & H03 & Effective rank as cross-seed AUROC predictor & IDM & immune (3 draws) & $\SV{2}$--$\SV{7}$ & $1{\times}10^{-4}$ & $1{\times}10^{-4}$ & R F H & promising \\
59 & H01 & Cross-seed TF boundary anchor stability + out-degree\ldots{} & MOD & kidney~A (3 draws) & $\SV{5}$--$\SV{7}$ & 0.9 & 0.9 &  & neutral \\
59 & H02 & Pairwise TF-target 6D distance as edge probability\ldots{} & MAN & kidney~A & $\SV{5}$--$\SV{7}$ & 0.002 & 0.002 & H & negative \\
59 & H03 & Partial Spearman: effective rank vs AUROC controlling\ldots{} & IDM & immune (3 draws) & $\SV{5}$--$\SV{7}$ & 0.795 & 0.795 &  & negative \\
60 & H01 & Signed displacement projection AUROC (directional\ldots{} & MAN & kidney~A & $\SV{5}$--$\SV{7}$ & 0.001 & 0.001 & H & negative \\
60 & H02 & TF family-stratified margin trajectory across 12 layers & MOD & kidney~A & $\SV{5}$--$\SV{7}$ & --- & --- &  & neutral \\
60 & H03 & Feed-forward loop triangle geometry in $\SV{5}$--$\SV{7}$ & MOD & kidney~A & $\SV{5}$--$\SV{7}$ & 0.002 & 0.002 & H & inconclusive \\
61 & H01 & Layer-to-layer CKA trajectory ($L_{8}$ boundary test) & CMA & immune (3 draws) & --- & --- & --- &  & negative \\
61 & H02 & Dual-role gene (TF intersect target) margin trajectory & MOD & immune (3 draws) & $\SV{5}$--$\SV{7}$ & 0.084 & 0.084 &  & neutral \\
61 & H03 & FFL geometry expanded (N=264) with multi-subspace and\ldots{} & MOD & immune (3 draws) & $\SV{2}$--$\SV{4}$, $\SV{5}$--$\SV{7}$, $\SV{2}$--$\SV{7}$ & 0.297 & 0.297 &  & negative \\
62 & H01$^{\star}$ & cycle4\_immune edge-level AUROC (735 pairs\ldots{} & MAN & immune (3 draws) & $\SV{5}$--$\SV{7}$/$\SV{2}$--$\SV{4}$, $\SV{5}$--$\SV{7}$, $\SV{2}$--$\SV{4}$ & 0.045 & 0.045 &  & promising \\
62 & H02 & AUROC mechanism diagnostic: sparsity vs depth as\ldots{} & IDM & immune & --- & $9.5{\times}10^{-7}$ & $9.5{\times}10^{-7}$ & R F H & promising \\
62 & H03 & Cross-cycle Procrustes alignment transfer (cycle1 to\ldots{} & CMA & kidney~A + immune & $\SV{5}$--$\SV{7}$ & --- & --- &  & negative \\
63 & H63-A & $\SV{2}$--$\SV{4}$ repulsion mechanistic probe: TF class separation & MOD & immune & $\SV{2}$--$\SV{4}$, $\SV{5}$--$\SV{7}$ & --- & --- &  & promising \\
63 & H63-G & Signed regulation AUROC split: Activation vs\ldots{} & MOD & immune & $\SV{5}$--$\SV{7}$, $\SV{2}$--$\SV{4}$ & --- & --- &  & promising \\
63 & H63-B & Systematic subspace scan SV(k,k+2) for k=0..9 at $L_{0}$\ldots{} & MAN & immune & $\SV{k}$--$\SV{k+2}$, $\SV{8}$--$\SV{10}$, $\SV{5}$--$\SV{7}$ & 0.28 & 0.28 &  & inconclusive \\
\end{longtable}
\endgroup

\setcounter{table}{12}
\begingroup
\footnotesize
\setlength{\tabcolsep}{2.4pt}
\setlength{\LTpre}{6pt}\setlength{\LTpost}{6pt}
\begin{longtable}{@{}>{\scriptsize\raggedright\arraybackslash}p{0.55cm}>{\scriptsize\raggedright\arraybackslash}p{1.85cm}>{\scriptsize\raggedright\arraybackslash}p{1.9cm}>{\scriptsize\raggedright\arraybackslash}p{2.75cm}>{\scriptsize\raggedright\arraybackslash}p{3.5cm}>{\scriptsize\raggedright\arraybackslash}p{2.2cm}>{\scriptsize\raggedright\arraybackslash}p{2.95cm}@{}}
\caption{\textbf{Claims ledger: the 61 numerical claims reported from the screen, traced and recomputed.} One row per numerical claim of Kendiukhov (2026b) (abstract, results, methods and tables), numbered C01--C61. ID, claim identifier (the file of Table~S12 lists, for each record, the claims it supports); Claim, what was claimed; Source records, iteration/hypothesis of the record(s) behind the claim, with the matrix the record used (kidney~A, the Tabula Sapiens kidney extraction with 4,803 genes, described as immune in Kendiukhov (2026b); immune, the immune extraction; 3 draws, the main extraction and two further 256-cell draws; earlier loop, iterations of an earlier run of the loop whose folders were overwritten) and its decomposition: 209-gene subset (22 zero rows), decomposition fitted to the 209 named genes at their true rows, 22 of which are never-observed all-zero rows; misaligned rows 0--208, the same with every label on the wrong row (blank lines of the gene list dropped before numbering); observed rows, only genes observed in at least one sampled cell, mean-centred over them (the convention of this paper); full matrix (zero rows), all rows including the never-observed ones, whose first direction is then the observation indicator. As reported, the value in Kendiukhov (2026b). Screen's procedure, recomputed, the value recomputed here with the screen's own code conventions on the matrix and gene rows the record used (brackets in C02: minimum--maximum over 50 random subsamples). Status: \textit{reproduces}, exactly or within resampling error (25 claims); \textit{partly} (2); \textit{does not reproduce} (5); \textit{provenance only}, traced to its record and code but not recomputed (27); \textit{not recomputable}, records not preserved (2). Observed-gene value; this paper: the same quantity recomputed on the observed genes or on the labels' true rows, where computed, and the section and table of the main text that examine the corresponding structure under the convention of this paper (---, not examined here). Abbreviations: OR, odds ratio; ER, effective rank; TwoNN, two-nearest-neighbour intrinsic dimension; emp $p$, empirical permutation $p$; $r_{\mathrm{rb}}$ and rbc, rank-biserial correlation; OOV, out of vocabulary; GC, germinal centre; FFL, feed-forward loop; CKA, centred kernel alignment; U, left singular vectors; S, singular values; TF, transcription factor; PPI, protein--protein interaction; GO BP, Gene Ontology biological process; EV, extracellular vesicle; mito, mitochondrion; DC, dendritic cell; CV, cross-validation; OOF, out-of-fold. No new intervals are computed for this table. Full values: \texttt{Table\_S13\_claims\_ledger.csv}.}\label{tab:S13}\\
\toprule
ID & Claim & Source records; data & As reported & Screen's procedure, recomputed & Status & Observed-gene value; this paper \\
\midrule
\endfirsthead
\multicolumn{7}{@{}l}{\footnotesize\textit{Table~\thetable{} (continued)}}\\
\toprule
ID & Claim & Source records; data & As reported & Screen's procedure, recomputed & Status & Observed-gene value; this paper \\
\midrule
\endhead
\midrule
\multicolumn{7}{@{}r@{}}{\scriptsize\textit{continued on next page}}\\
\endfoot
\bottomrule
\endlastfoot
C01 & Size of the screen & all records, iterations 1--63 & 63 iterations; 183 hypotheses; 13 families & 183 records in 61 JSON files (63 iteration dirs; iter\_\allowbreak{}0001 and iter\_\allowbreak{}0003 have no JSON); 8 JSON families & \textit{reproduces}: 183/\allowbreak{}63; no (families: 8 in the JSON, 13 is Kendiukhov (2026b)'s own regrouping) & Sec.~2.9; Sec.~4.16 (Table~S12) \\
C02 & Effective-rank collapse & 46/H02\newline immune; 2,000-row subsample (zero rows) & 23.6 ($L_{0}$) $\to$ 1.6 ($L_{11}$), 14.4-fold, `computed on the full 4,803-gene vocabulary' & 50 random subsamples: $L_{0}$ 23.92 [22.44--25.24], $L_{11}$ 1.67; fold 14.3; full 4,941 rows: 25.46$\to$1.67 (15.2-fold); kidney~A full 4,803 rows (what the text says): 19.54$\to$1.63 (12.0-fold) & \textit{reproduces}: on immune subsample (approx.; seed unknown); no on the matrix the text names & \emph{Observed genes / true rows:} immune observed 2,039 rows: 236.9$\to$48.7 (4.86-fold); Sec.~2.2; Table~2; Table~3 (Table~S1) \\
C03 & $\SV{1}$ variance fraction & 46/H02, 7/H02\newline immune; zero rows included & 93.4\% at $L_{11}$ (53.7\% at $L_{0}$), `full vocabulary' & immune subsample mean $L_{0}$ 53.5\%, $L_{11}$ 93.1\%; kidney~A full 4,803: $L_{0}$ 57.2\%, $L_{11}$ 93.2\% & \textit{reproduces}: approx. & \emph{Observed genes / true rows:} immune observed: 5.6\% $\to$ 18.6\%; kidney~A observed: 8.1\% $\to$ 23.1\%; Sec.~2.2; Table~2; Table~3 \\
C04 & $\SV{1}$ variance on the named-gene submatrix & 9/H03, 10/H03\newline kidney~A; 209-gene subset (22 zero rows) & 19\% ($L_{0}$) $\to$ 77\% ($L_{11}$), `195-gene submatrix' & 19.05\% $\to$ 76.74\% & \textit{reproduces}: on 209 rows incl. 22 zero rows & \emph{Observed genes / true rows:} 187 observed named genes: 7.1\% $\to$ 27.5\%; misaligned `195': 8.0\% $\to$ 25.9\%; Sec.~2.2 (all observed genes) \\
C05 & TwoNN intrinsic dimension & 45/H02, 46/H01\newline immune; not spectral & 32.6 $\to$ 18.1 (44.6\% decrease), 2,000 random genes, three seeds & iter\_\allowbreak{}0042 estimator, seeds 42/\allowbreak{}0/\allowbreak{}1: $L_{0}$ 28.4/\allowbreak{}27.8/\allowbreak{}34.4, $L_{11}$ 18.7/\allowbreak{}16.9/\allowbreak{}20.3 & \textit{does not reproduce}: not exactly; seed and estimator of iter\_\allowbreak{}0045 unknown, value depends on the sample: $L_{0}$ range 27.8--34.4 across 3 seeds & --- \\
C06 & Participation ratio & 31/H03\newline kidney~A; misaligned rows 0--208 & 58 $\to$ 9.5 (6.1-fold), `195-gene submatrix' & 58.09 $\to$ 9.53 (misaligned rows, as executed) & \textit{reproduces}: exactly, on the misaligned rows & \emph{Observed genes / true rows:} true 187 observed named genes: 59.47 $\to$ 10.28; Sec.~2.2 (participation ratio, all observed genes) \\
C07 & Feature-shuffle control for effective rank & 46/H02\newline immune; 2,000-row subsample (zero rows) & ER 28.9 at $L_{11}$ (17.6$\times$ above observed) & mean of 10 shuffles: 29.2 & \textit{reproduces}: approx. & Sec.~2.2 (feature-permutation control) \\
C08 & $\SV{1}$ extracellular pole & 6/H02\newline kidney~A; full matrix (zero rows) & GO:0005615 OR = 6.37, $p$ = $2.6{\times}10^{-4}$ ($\SV{1}$, $L_{11}$) & OR = 6.3677, $p$ = $2.57{\times}10^{-4}$ (21 vs 5 genes) & \textit{reproduces}: exactly & \emph{Observed genes / true rows:} observed-gene SVD, 187 genes, quartiles: OR = 1.00, $p$ = 0.60; Sec.~2.3; Table~3 (Table~S2) \\
C09 & $\SV{1}$ cytosolic pole & 7/H03\newline kidney~A; full matrix (zero rows) & GO:0005829 OR = 2.96, $p$ = 0.010 & OR = 2.956, $p$ = 0.0104 & \textit{reproduces}: exactly & Sec.~2.3 (Table~S2) \\
C10 & $\SV{1}$ gene-label-shuffle null & 7/H01\newline kidney~A; full matrix (zero rows) & $N$ = 500, empirical $p$ = 0.004 & --- & \textit{provenance only}: null = permuting named genes over rows of a matrix whose $\SV{1}$ is the zero-row axis; the statistic is the min Fisher p over 296 terms & Sec.~2.3 (max-$T$ over 288 tests) \\
C11 & Compartment transients on $\SV{1}$ & 8/H02, 9/H02\newline kidney~A; 209-gene subset (22 zero rows) & mito OR 23.3 at $L_{3}$; endoplasmic-reticulum lumen OR 4.7 $\to$ 18.5 ($L_{1}$--$L_{11}$) & --- & \textit{provenance only} & Sec.~2.3 (compartment gradient) \\
C12 & $\SV{2}$ /\allowbreak{} compartment enrichments & 8/H01, 9/H01, 10/H01\newline kidney~A; 209-gene subset (22 zero rows) & $\SV{2}$ IL-4 emp $p$$<$0.001; EV OR 7.99; cytoskeleton OR 20.4 & --- & \textit{provenance only} & Sec.~2.3 (first twelve directions tested) \\
C13 & TRRUST co-pole on $\SV{2}$ & 10/H02, 11/H01\newline kidney~A; 209-gene subset (22 zero rows) & observed 0.206 vs null 0.122; $p$ $<$ 0.001 at all 12 layers & $L_{11}$ $\SV{2}$ co-pole = 0.2055 over 326 rows ($\SV{1}$: 0.1626) & \textit{reproduces}: 0.206 & Sec.~2.6 (within-factor edge benchmark) \\
C14 & Sign specificity of TRRUST co-pole & 11/H01\newline kidney~A; 209-gene subset (22 zero rows) & activation 12/\allowbreak{}12 layers; repression 1/\allowbreak{}12 & --- & \textit{provenance only} & --- \\
C15 & $\SV{3}$ signalling axis & 14/H01\newline kidney~A; 209-gene subset (22 zero rows) & all layers p $\leq$ 0.017 & --- & \textit{provenance only} & --- \\
C16 & $\SV{1}$ stability and annotation-density check & 9/H03, 10/H03\newline kidney~A; 209-gene subset (22 zero rows) & mean adjacent-layer Spearman $r$ = 0.929; zero annotation-density confounders & mean adjacent-layer Spearman (raw signs) = 0.9283 & \textit{reproduces} & --- \\
C17 & $\SV{1}$ depleted of TFs & 51/H02\newline immune; observed rows & top-quartile $\SV{1}$: $\sim$9$\times$ fewer TFs, OR = 0.108, $p$ $<$ $1{\times}10^{-4}$ & --- & \textit{provenance only} & Sec.~2.5 (single directions, Table~S6) \\
C18 & Cell-type marker clustering & 23/H01\newline kidney~A; not spectral & AUROC = 0.851 (12/\allowbreak{}12 layers $p$ $<$ $1{\times}10^{-8}$); random-partition control 0.488 & mean AUROC = 0.8508 (float32, as executed; float64 gives 0.8504) & \textit{reproduces}: exactly in float32 & \emph{Observed genes / true rows:} dropping the 3 never-observed markers (DCN, COL1A1, MUC1): mean AUROC 0.843 (T and B only) \\
C19 & Marker AUROC trend with depth & 23/H01\newline kidney~A; not spectral & slope = +0.036 per layer, $r$ = 0.765, $p$ = 0.004 (`increases with depth') & slope of logit(AUROC) on layer = $-$0.0357, $r$ = $-$0.765, $p$ = 0.0038 & \textit{does not reproduce}: sign reversed in Kendiukhov (2026b): AUROC decreases with depth & Table~3 \\
C20 & Gene-family clustering & 26/H01\newline kidney~A; not spectral & HLA-I 1.000; AP1 0.969; RUNX 0.925; IL-2 0.869 & --- & \textit{provenance only} & --- \\
C21 & STRING PPI co-pole by axis (Tables 1 and S4 of Kendiukhov 2026b) & 12/H03, 13/H01\newline kidney~A; 209-gene subset (22 zero rows) & 1,022 pairs ($\geq$ 0.7); $\SV{2}$ mean 0.226 vs null 0.124 (12/\allowbreak{}12); $\SV{3}$ 0.198 (12/\allowbreak{}12); $\SV{1}$ 0.119 (2/\allowbreak{}12) & pairs = 1022; mean co-pole $\SV{2}$ 0.226, $\SV{3}$ 0.198, $\SV{1}$ 0.128; $\SV{2}$ per-layer = Table S4 of Kendiukhov (2026b) exactly & \textit{reproduces}: for $\SV{2}$/\allowbreak{}$\SV{3}$ and N pairs; no for $\SV{1}$ (0.128 vs 0.119) & Sec.~2.4; Table~3 (Tables~S3, S4) \\
C22 & $\SV{4}$ /\allowbreak{} $\SV{5}$ PPI co-pole & 15/H01\newline kidney~A; 209-gene subset (22 zero rows) & $\SV{4}$ 0.158 (7/\allowbreak{}12); $\SV{5}$ 0.149 (6/\allowbreak{}12) on 3,092 pairs ($\geq$ 0.4) & pair count = 3092; rates not recomputed & \textit{partly}: pair count & Sec.~2.9 (Table~S14) \\
C23 & STRING confidence gradient & 15/H02\newline kidney~A; 209-gene subset (22 zero rows) & quintile z 1.00/\allowbreak{}1.48/\allowbreak{}2.09/\allowbreak{}3.18/\allowbreak{}5.02; Spearman $\rho$ = 1.000, $n$ = 5, $p$ = 0.017 & exact two-sided p for $\rho$ = 1 with $n$ = 5: 0.0167; the record's $p$ = $1.4{\times}10^{-24}$ is scipy's t-approximation ($1.40{\times}10^{-24}$) at $\rho$ = 1 & \textit{reproduces}: $p$ = 0.017 is right; the screen's $1.4{\times}10^{-24}$ is an artefact & Sec.~2.4; Table~3 (STRING confidence bands) \\
C24 & PPI-only vs GO-only pairs & 15/H03\newline kidney~A; 209-gene subset (22 zero rows) & PPI-only $N$ = 670, $z$ = 2.66 (12/\allowbreak{}12); GO-only $z$ = 1.89 (7/\allowbreak{}12) & --- & \textit{provenance only} & --- \\
C25 & Hub control & 14/H03\newline kidney~A; 209-gene subset (22 zero rows) & low-degree $z$ = 3.1 vs hub $z$ = 2.6 & --- & \textit{provenance only} & Sec.~2.4 (hub removal, degree-stratified null; Table~S3) \\
C26 & Orthogonality of $\SV{2}$--$\SV{4}$ PPI axes & 16/H02\newline kidney~A; 209-gene subset (22 zero rows) & pairwise Pearson $r$ $<$ 0.25 & --- & \textit{provenance only} & --- \\
C27 & TRRUST pairs absent from STRING & 22/H01\newline kidney~A; not spectral & $N$ = 141, AUROC = 0.573, 12/\allowbreak{}12 layers $p$ $<$ 0.007 & --- & \textit{provenance only} & Sec.~2.6 (within-factor edge benchmark) \\
C28 & TF vs target classifier, joint 6-D $\SV{2}$--$\SV{7}$ & 56/H01\newline immune; observed rows & AUROC = 0.744 (mean over layers), max 0.789 at $L_{3}$, all 12 layers $p$ $<$ 0.01 (100 perms); $n$ = 51 TFs vs 144 targets, cosine-to-TF score, 195-gene set & mean 0.747, max 0.792 ($L_{2}$); pooled-OOF mean 0.745; $L_{3}$ 100-perm $p$ = 0.0099 (= floor 1/\allowbreak{}101) & \textit{reproduces}: approx.; $\pm$0.01, CV split seed unknown & Sec.~2.5; Table~3 (Table~S6) \\
C29 & Joint beats single subspaces; cross-seed & 56/H01, 57/H01\newline immune (3 draws); observed rows & joint $>$ both in 11/\allowbreak{}12 layers; cross-seed means 0.744/\allowbreak{}0.753/\allowbreak{}0.757 (SD $\sim$0.016) & joint $>$ both in 10/\allowbreak{}12; seed43 mean 0.749, seed44 0.756 & \textit{reproduces}: approx. & Sec.~2.5 (all 78 contiguous bands) \\
C30 & Depth division of labour & 55/H02, 56/H01\newline immune; observed rows & $\SV{5}$--$\SV{7}$ dominates $L_{0}$--$L_{3}$ ($\sim$0.71); $\SV{2}$--$\SV{4}$ $L_{4}$--$L_{8}$ ($\sim$0.72) & $\SV{5}$--$\SV{7}$ $L_{0}$--$L_{3}$ mean 0.716; $\SV{2}$--$\SV{4}$ $L_{4}$--$L_{8}$ mean 0.737 & \textit{reproduces}: approx. & Sec.~2.5 (all 78 contiguous bands) \\
C31 & $\SV{2}$--$\SV{4}$ signal explained by co-expression & 51/H01\newline immune; observed rows & residual non-significant at 0/\allowbreak{}12 layers & layers with $p$ $<$ 0.05: 1/\allowbreak{}12 ($L_{8}$ rbc = 0.083, $p$ = 0.0016) & \textit{does not reproduce}: 1/\allowbreak{}12, not 0/\allowbreak{}12 (the JSON record says 1/\allowbreak{}12) & Sec.~2.6 (co-expression comparator) \\
C32 & $\SV{5}$--$\SV{7}$ regulatory proximity after co-expression residualisation & 52/H03, 53/H01\newline immune; observed rows & $r_{\mathrm{rb}}$ = 0.148, $p$ $<$ 0.001 at $L_{0}$; 100/\allowbreak{}100 bootstrap positive & $L_{0}$ rbc = 0.148, $L_{1}$ 0.119, $L_{2}$ 0.083; significant layers 3/\allowbreak{}12; $L_{5}$--$L_{11}$ rbc $<$ 0 & \textit{reproduces}: exactly & Sec.~2.6; Sec.~2.7 \\
C33 & Edge-level TF-target AUROC & 62/H01\newline immune; observed rows & 0.602 at $L_{0}$ ($\SV{5}$--$\SV{7}$); 589 positives vs 2,351 negatives & $L_{0}$ 0.6017, $L_{1}$ 0.6006 (U coords, 1,523 candidates); with U*S: 0.6053 & \textit{reproduces}: exactly, with 1,523 candidate pairs & Sec.~2.6; Table~3 (Table~S7) \\
C34 & Edge AUROC decays with depth & 62/H01, 62/H02\newline immune; observed rows & Spearman $\rho$ = $-$0.958, $p$ = $9.5{\times}10^{-7}$; $L_{0}$--$L_{8}$ significant, $L_{9}$--$L_{11}$ not & $\rho$ = $-$0.958, $p$ = $9.54{\times}10^{-7}$ & \textit{reproduces}: exactly & Sec.~2.6 (every layer, Table~S7) \\
C35 & $\SV{2}$--$\SV{4}$ edge AUROC & 62/H01\newline immune; observed rows & mean 0.492 (near/\allowbreak{}below chance) & mean 0.492 & \textit{reproduces} & Sec.~2.6; Sec.~2.7 \\
C36 & Non-zero embeddings & 62/H02\newline immune & constant at 2,039 of the 4,803 vocabulary genes & immune: 2039 of 4941; kidney~A: 1972 of 4803 & \textit{does not reproduce}: 2,039 is the immune count; the 4,803-gene matrix has 1,972 & Sec.~2.1; Table~1 \\
C37 & Repression more separable than activation & 63/H63-G\newline immune; observed rows & $\Delta$ = +0.021 ($\SV{5}$--$\SV{7}$), +0.084 ($\SV{2}$--$\SV{4}$) & --- & \textit{provenance only} & --- \\
C38 & Mechanistic probe & 63/H63-A\newline immune; observed rows & TF-TF 0.539 /\allowbreak{} TF-target 0.485 ($\SV{2}$--$\SV{4}$); 0.528 /\allowbreak{} 0.599 ($\SV{5}$--$\SV{7}$) & --- & \textit{provenance only} & --- \\
C39 & B-cell kNN clustering & 37/H01\newline kidney~A; not spectral; misaligned rows & precision@10 $z$ = 7.55 ($p$ $<$ 0.001); T-cell $z$ = $-$1.37; DC $z$ = $-$0.87 & misaligned (as executed): $L_{2}$ $z$ = 7.55 & \textit{reproduces}: exactly, on misaligned rows & \emph{Observed genes / true rows:} true rows (187 observed): $z$ = $L_{0}$ 7.16, $L_{2}$ 8.22, $L_{5}$ 8.30, $L_{8}$ 7.65, $L_{11}$ 6.31; Table~3 \\
C40 & B-cell signal not driven by norm & 34/H03\newline kidney~A; misaligned rows 0--208 & survives $L_{2}$-norm regression (residual $p$ = 0.002) & --- & \textit{provenance only} & --- \\
C41 & BATF/\allowbreak{}BACH2 convergence to the B-cell centroid & 42/H03\newline immune; not spectral & BATF rank 1,510 $\to$ 189, dist to PAX5 18.3 $\to$ 5.1, $\rho$ = $-$0.972; BACH2 611 $\to$ 146, 17.4 $\to$ 5.2, $\rho$ = $-$0.844 & BATF ranks 1510 $\to$ 189, dist 18.32 $\to$ 5.13, $\rho$ $-$0.972; BACH2 611 $\to$ 146, 17.41 $\to$ 5.15, $\rho$ $-$0.844 & \textit{reproduces}: exactly & --- \\
C42 & PAX5 as stable anchor & 40/H01, 41/H02\newline immune; not spectral & rank $\sim$66 from $L_{0}$; convergence onset $L_{3}$ & PAX5 rank (genes closer) $L_{0}$ = 65, range 28--72 & \textit{reproduces}: 65 = $66 - 1$ & --- \\
C43 & Whole GC circuit converges & 43/H01\newline immune; not spectral & PRDM1 $\rho$ = $-$1.000; BACH2-PRDM1 tightest pair (3.94) at $L_{11}$ & --- & \textit{provenance only} & --- \\
C44 & GC vs plasma programmes diverge & 38/H02\newline kidney~A; not spectral; misaligned rows & angle 77$^\circ$ ($L_{0}$) $\to$ 94$^\circ$ ($L_{11}$) & misaligned (as executed): 77.1 $\to$ 94.4$^\circ$ & \textit{reproduces}: exactly, on misaligned rows & \emph{Observed genes / true rows:} true rows: 62.7 $\to$ 55.0$^\circ$ (the trend reverses); Table~3 \\
C45 & BCL6 metabolic neighbourhood & 41/H02, 42/H02\newline immune; not spectral & 10 of 20 nearest neighbours metabolic at every layer & --- & \textit{provenance only} & --- \\
C46 & B-cell manifold compression & 41/H03\newline immune; not spectral & TwoNN 8.2 $\to$ 5.1, $\rho$ = $-$0.951; T-cell $\rho$ = +0.287 ($p$ = 0.37); myeloid +0.699 ($p$ = 0.011) & --- & \textit{provenance only} & --- \\
C47 & Lineage-specific delay & 41/H01\newline immune; not spectral & T-cell /\allowbreak{} myeloid TFs start at ranks 170 /\allowbreak{} 86; GC $\rho$ = $-$0.993 ($p$ $<$ $1{\times}10^{-6}$) & --- & \textit{provenance only} & --- \\
C48 & Persistent homology collapses under rewiring & historical loop iter\_\allowbreak{}0003--0008 (folders overwritten)\newline earlier loop & 11/\allowbreak{}12 layers $p$ $<$ 0.05 (feature shuffle); 0/\allowbreak{}24 under degree-preserving rewiring & --- & \textit{not recomputable}: iteration folders overwritten by an earlier loop & --- \\
C49 & Cross-model alignment partial & historical loop iter\_\allowbreak{}0003\newline earlier loop & cosine 0.825; permutation $p$ $>$ 0.3 & --- & \textit{not recomputable}: iteration folders overwritten by an earlier loop & --- \\
C50 & Geneformer static embeddings encode STRING & 19/H01\newline Geneformer + kidney~A & $p$ = $7.8{\times}10^{-127}$ & --- & \textit{provenance only} & Sec.~2.7 (static embeddings as comparators) \\
C51 & Geneformer B-cell structure absent & 40/H02\newline Geneformer & precision@10 = 0.000 & --- & \textit{provenance only} & --- \\
C52 & ER-AUROC correlation confounded by layer & 58/H03, 59/H03\newline immune (3 draws); observed rows & $\rho$ = 0.855; partial $\rho$ = $-$0.045 & --- & \textit{provenance only} & --- \\
C53 & GO BP not encoded in $\SV{2}$ poles & 11/H03\newline kidney~A; 209-gene subset (22 zero rows) & 591 terms, zero significant & --- & \textit{provenance only} & --- \\
C54 & Feed-forward-loop geometry rejected & 61/H03\newline immune; observed rows & 264 FFLs, no significant result after permutation correction & --- & \textit{provenance only} & --- \\
C55 & Cross-seed CKA & 61/H01\newline immune (3 draws); not spectral & 0.979 ($L_{0}$) $\to$ 0.779 ($L_{11}$) & --- & \textit{provenance only} & Sec.~2.8 (variance across cell draws) \\
C56 & Attention vs residual geometry & 18/H03\newline lung attention, kidney~A labels & attention TRRUST 2$\times$ ($p$ = $9.9{\times}10^{-9}$); STRING $p$ = 0.091 & --- & \textit{provenance only} & --- \\
C57 & Graph topology & 2/H01\_\allowbreak{}graph\_\allowbreak{}topology\_\allowbreak{}cc\newline kidney~A (3 draws); not spectral & clustering coefficient $z$ $>$ 9 at all 36 tests & --- & \textit{provenance only} & --- \\
C58 & TRRUST signed subsets on the edge benchmark & 63/H63-G\newline immune & 589 pairs: 270 activation-only, 141 repression-only, 178 mixed/\allowbreak{}ambiguous & any-Activation 270, any-Repression 141 (overlap 31); activation-only 239, repression-only 110, unknown-only 209 & \textit{partly}: 589 yes; 270/\allowbreak{}141 are `any' counts that overlap by 31, so they are not `only' sets and 178 = $589-270-141$ is not a count of anything & --- \\
C59 & Named genes & 30/H04\newline kidney~A & 209 named; 14 OOV (zero norm); 195 in-vocabulary genes used for all biological analyses & true rows: 22 zero, 187 observed; misaligned mapping: 14 zero (FOS, HLA-A, HLA-DPB1, JUNB, KLF6, LDHA, LGALS1, NCAM1, NCOA3, NR4A3, PAX5, PTGS2, TBXAS1, TNF) $\to$ 195 & \textit{does not reproduce}: 195 /\allowbreak{} 14 OOV are artefacts of the misaligned mapping; truth is 187 observed /\allowbreak{} 22 never-observed & Sec.~2.1; Sec.~2.9 (observed genes) \\
C60 & TRRUST pairs in the named set & 10/H02\newline kidney~A & 326 TF-target pairs (screen record; Kendiukhov (2026b) uses the resulting co-pole rate) & 326 TRRUST rows = 288 distinct pairs; 245 with both genes observed & \textit{reproduces}: as a row count; 288 distinct pairs & Sec.~2.6; Table~1 (edge benchmark counts) \\
C61 & Front-loaded compression & 6/H03\newline kidney~A; 209-gene subset (22 zero rows) & effective rank below half its initial value by layer 4 & --- & \textit{provenance only} & Sec.~2.2 \\
\end{longtable}
\endgroup

\setcounter{table}{13}
\begingroup
\footnotesize
\setlength{\tabcolsep}{2.4pt}
\setlength{\LTpre}{6pt}\setlength{\LTpost}{6pt}
\begin{longtable}{@{}>{\scriptsize\raggedright\arraybackslash}p{0.4cm}>{\scriptsize\raggedright\arraybackslash}p{1.35cm}>{\scriptsize\raggedright\arraybackslash}p{1.3cm}>{\scriptsize\raggedright\arraybackslash}p{2.35cm}>{\scriptsize\raggedright\arraybackslash}p{5.1cm}>{\scriptsize\raggedright\arraybackslash}p{5.4cm}@{}}
\caption{\textbf{Correspondence of other spectral bases with the observed-gene basis, and provenance of the screen's three bands.} \emph{Panel A}, correspondence of single directions, from the right singular vectors (512-dimensional feature space; $|\cos|$ is sign-free) at each layer. kidney~A, full matrix: the decomposition of all 4,803 rows, mean-centred over all of them, including the 2,831 never-observed zero rows, whose first direction is therefore the observation indicator; it is compared with the decomposition of the 1,972 observed genes (this paper's convention). For its $\SV{k}$, $k = 2, \ldots, 12$: best $= \SV{k-1}$, number of directions whose best-matching observed direction (largest $|\cos|$ among observed $\SV{1}$--$\SV{40}$) is $\SV{k-1}$; mean and range of that best $|\cos|$ over the 11 directions; $\SV{1}$: best $|\cos|$ of its first direction with any observed direction (and that direction), showing that it has no counterpart. Pooled over the 132 layer--direction cases, the mean $|\cos|$ between full-matrix $\SV{k}$ and observed $\SV{k-1}$ exactly is 0.956 (independent re-derivation); at layers where all 11 best matches are $\SV{k-1}$ it equals the mean shown. kidney~A, 209-gene: mean best $|\cos|$ over $\SV{1}$--$\SV{12}$ for the decomposition fitted to the 209 named genes at their true rows (22 of them zero rows), on which the screen assembled $\SV{2}$--$\SV{4}$. immune, screen's basis: the screen's decomposition of the 2,039 observed rows (``centered nonzero embeddings'') against this paper's; best $= \SV{k}$, number of $\SV{1}$--$\SV{12}$ whose best match is the same index, with the mean $|\cos|$ at the same index. \emph{Panel B}, mapping of bands: band, the band's indices in that decomposition; matrix and decomposition (rows), the matrix and the rows the decomposition was fitted to (mean, mean embeddings; summed, the summed-embedding matrix the screen loaded in iterations 59--60); used in, the iterations of the screen whose records use the band in that decomposition (---, a reference basis: the full kidney~A decomposition, never-observed rows included); mapped, the contiguous observed-gene band of the same width (ending at or before $\SV{12}$) with the largest mean $\cos^2$ of the principal angles between the two subspaces, averaged over the 12 layers ($\overline{\cos^2}$; range, its range over the 12 layers); identity, the same quantity for the band with the same indices; runner-up, the second-best band; layers won, number of layers at which the mapped band is the best (other winning bands in parentheses). The mapping uses the geometry of the matrices only; no outcome data (interaction, regulatory or label sets) were read to fix it. \emph{Panel C}, provenance of the bands from the iteration files (JSON records, iteration reports, the language-model brainstormer's briefs and next-step notes): \#, step; iteration, iteration and hypothesis or file; band, the band the step bears on (all, every direction); text in quotation marks is verbatim from those files. ``sig'' is the screen's count of layers with empirical permutation $p < 0.05$; emp\_p and perm\_p, empirical permutation $p$; mean\_z, mean $z$ over layers; rbc, rank-biserial correlation; TRRUST, STRING, reference networks of documented regulation and protein interaction; PPI, protein--protein interaction; TF, transcription factor. None of the three bands was specified before the results that motivated it were seen. No sampling intervals apply: the matrices are fixed and all values are exact geometric quantities. Full values: \texttt{Table\_S14\_band\_crosswalk.csv}.}\label{tab:S14}\\
\multicolumn{6}{@{}l}{\footnotesize\textbf{A.} $\SV{k}$ of the full-matrix, 209-gene and screen's immune decompositions against the observed-gene $\SV{k}$ and $\SV{k-1}$, per layer}\\[2pt]
\multicolumn{6}{@{}l@{}}{\scriptsize\setlength{\tabcolsep}{3.2pt}\begin{tabular}{@{}lrrrrrrr@{}}
\toprule
 & \multicolumn{4}{c}{kidney~A: full matrix, zero rows included} & \multicolumn{1}{c}{kidney~A: 209-gene} & \multicolumn{2}{c}{immune: screen's basis} \\
\cmidrule(lr){2-5}\cmidrule(lr){6-6}\cmidrule(lr){7-8}
Layer & best $= \SV{k-1}$ & mean $|\cos|$ & range $|\cos|$ & $\SV{1}$: best $|\cos|$ & mean best $|\cos|$ & best $= \SV{k}$ & mean $|\cos|$ \\
\midrule
$L_{0}$ & 10/11 & 0.971 & 0.824--0.999 & 0.219 ($\SV{3}$) & 0.488 & 12/12 & 1.000 \\
$L_{1}$ & 11/11 & 0.994 & 0.983--1.000 & 0.233 ($\SV{12}$) & 0.483 & 12/12 & 1.000 \\
$L_{2}$ & 11/11 & 0.992 & 0.952--1.000 & 0.270 ($\SV{12}$) & 0.485 & 12/12 & 1.000 \\
$L_{3}$ & 9/11 & 0.941 & 0.686--1.000 & 0.224 ($\SV{10}$) & 0.501 & 12/12 & 1.000 \\
$L_{4}$ & 11/11 & 0.982 & 0.916--1.000 & 0.261 ($\SV{8}$) & 0.492 & 12/12 & 1.000 \\
$L_{5}$ & 9/11 & 0.946 & 0.702--1.000 & 0.309 ($\SV{7}$) & 0.473 & 12/12 & 1.000 \\
$L_{6}$ & 9/11 & 0.969 & 0.825--1.000 & 0.356 ($\SV{6}$) & 0.474 & 12/12 & 1.000 \\
$L_{7}$ & 11/11 & 0.958 & 0.785--1.000 & 0.309 ($\SV{5}$) & 0.473 & 12/12 & 1.000 \\
$L_{8}$ & 11/11 & 0.973 & 0.901--1.000 & 0.265 ($\SV{5}$) & 0.524 & 12/12 & 1.000 \\
$L_{9}$ & 11/11 & 0.985 & 0.957--0.998 & 0.232 ($\SV{5}$) & 0.543 & 12/12 & 1.000 \\
$L_{10}$ & 11/11 & 0.976 & 0.904--0.999 & 0.211 ($\SV{5}$) & 0.584 & 12/12 & 1.000 \\
$L_{11}$ & 11/11 & 0.994 & 0.987--1.000 & 0.155 ($\SV{3}$) & 0.555 & 12/12 & 1.000 \\
\midrule
all layers & 125/132 (94.7\%) & 0.973 & 0.686--1.000 & 0.254 (mean) & 0.506 & 144/144 & 1.000 \\
\bottomrule
\end{tabular}}\\[6pt]
\multicolumn{6}{@{}l}{\footnotesize\textbf{B.} The three bands in each decomposition, mapped to the observed-gene basis}\\[2pt]
\multicolumn{6}{@{}l@{}}{\scriptsize\setlength{\tabcolsep}{2.4pt}\begin{tabular}{@{}l l >{\raggedright\arraybackslash}p{2.2cm} >{\raggedright\arraybackslash}p{2.25cm} l r r r l >{\raggedright\arraybackslash}p{1.85cm}@{}}
\toprule
Band & Matrix & Decomposition (rows) & Used in & Mapped & $\overline{\cos^2}$ & range & identity & runner-up & layers won \\
\midrule
$\SV{2}$--$\SV{4}$ & kidney~A & full matrix (2,831 zero rows) & --- (reference basis) & $\SV{1}$--$\SV{3}$ & 0.990 & 0.966--0.999 & 0.659 & $\SV{2}$--$\SV{4}$ (0.659) & 12/12 \\
$\SV{5}$--$\SV{7}$ & kidney~A & full matrix (2,831 zero rows) & --- (reference basis) & $\SV{4}$--$\SV{6}$ & 0.938 & 0.697--0.993 & 0.649 & $\SV{3}$--$\SV{5}$ (0.650) & 12/12 \\
$\SV{2}$--$\SV{7}$ & kidney~A & full matrix (2,831 zero rows) & --- (reference basis) & $\SV{1}$--$\SV{6}$ & 0.966 & 0.848--0.992 & 0.820 & $\SV{2}$--$\SV{7}$ (0.820) & 12/12 \\
$\SV{2}$--$\SV{4}$ & kidney~A & 209 named genes (22 zero rows) & iterations 13--19 (co-pole axes) & $\SV{1}$--$\SV{3}$ & 0.489 & 0.350--0.640 & 0.452 & $\SV{2}$--$\SV{4}$ (0.452) & 8/12 ($\SV{2}$--$\SV{4}$ 3, $\SV{5}$--$\SV{7}$ 1) \\
$\SV{5}$--$\SV{7}$ & kidney~A & observed (1,972), mean & later iterations on kidney~A & $\SV{5}$--$\SV{7}$ & 1.000 & 1.000--1.000 & 1.000 & $\SV{4}$--$\SV{6}$ (0.667) & 12/12 \\
$\SV{5}$--$\SV{7}$ & kidney~A & observed (1,972), summed & iterations 59--60 & $\SV{3}$--$\SV{5}$ & 0.212 & 0.146--0.319 & 0.198 & $\SV{4}$--$\SV{6}$ (0.210) & 2/12 (7 other bands) \\
$\SV{2}$--$\SV{4}$ & immune & observed (2,039) & iterations 49+ (TRRUST, TF/target) & $\SV{2}$--$\SV{4}$ & 1.000 & 1.000--1.000 & 1.000 & $\SV{1}$--$\SV{3}$ (0.667) & 12/12 \\
$\SV{5}$--$\SV{7}$ & immune & observed (2,039) & iterations 50+ & $\SV{5}$--$\SV{7}$ & 1.000 & 1.000--1.000 & 1.000 & $\SV{4}$--$\SV{6}$ (0.667) & 12/12 \\
$\SV{2}$--$\SV{7}$ & immune & observed (2,039) & iterations 56+ (joint classifier) & $\SV{2}$--$\SV{7}$ & 1.000 & 1.000--1.000 & 1.000 & $\SV{1}$--$\SV{6}$ (0.833) & 12/12 \\
\bottomrule
\end{tabular}}\\[6pt]
\multicolumn{6}{@{}l}{\footnotesize\textbf{C.} How the bands were chosen}\\
\toprule
\# & Iteration & Band & Matrix, gene rows, convention & What was done or decided & Result or criterion (quotations verbatim) \\
\midrule
\endfirsthead
\multicolumn{6}{@{}l}{\footnotesize\textit{Table~\thetable{} (continued)}}\\
\toprule
\# & Iteration & Band & Matrix, gene rows, convention & What was done or decided & Result or criterion (quotations verbatim) \\
\midrule
\endhead
\midrule
\multicolumn{6}{@{}r@{}}{\scriptsize\textit{continued on next page}}\\
\endfoot
\bottomrule
\endlastfoot
1 & 8 H01 & $\SV{2}$--$\SV{4}$ & kidney~A, 209 true rows, subset-fitted & first use of $\SV{2}$ and $\SV{3}$: ``SVD of layer-11 mean-centered 209-gene embedding. Project onto $\SV{2}$ and $\SV{3}$.'' & $\SV{2}$ IL-4 term emp\_\allowbreak{}p=0.000; $\SV{3}$ emp\_\allowbreak{}p=0.088 \\
2 & 10 H02 & $\SV{2}$--$\SV{4}$ & same as above & TRRUST co-pole on $\SV{1}$ and $\SV{2}$ at $L_{11}$ (326 TRRUST rows, i.e. 288 distinct pairs) & ``$\SV{1}$: obs=0.163 vs null=0.123, emp\_\allowbreak{}p=0.016; $\SV{2}$: obs=0.206 vs null=0.122, emp\_\allowbreak{}p=0.000'' \\
3 & 13 H01 & $\SV{2}$--$\SV{4}$ & same as above & STRING (score $\geq$ 0.7, 1,022 API pairs) co-pole on $\SV{1}$, $\SV{2}$ and $\SV{3}$ & ``$\SV{1}$ $z$=0.39 (2/\allowbreak{}12 sig), $\SV{2}$ $z$=10.04 (12/\allowbreak{}12 sig), $\SV{3}$ $z$=7.18 (12/\allowbreak{}12 sig)'' \\
4 & 14 brief & $\SV{2}$--$\SV{4}$ & --- & brainstormer proposes extending to the next axes: ``H01: $\SV{4}$/\allowbreak{}$\SV{5}$ PPI co-pole test'', with the success criterion ``$\SV{4}$ or $\SV{5}$ $z$ $>$ 3 in $\geq$6/\allowbreak{}12 layers'' & --- \\
5 & 15 H01 & $\SV{2}$--$\SV{4}$ & same, but 3,092 pairs at STRING $\geq$ 0.4, not the brief's 1,022 at $\geq$ 0.7 & $\SV{4}$ and $\SV{5}$ tested & ``$\SV{4}$ mean\_\allowbreak{}z=2.21, sig=7/\allowbreak{}12; $\SV{5}$ mean\_\allowbreak{}z=1.65, sig=6/\allowbreak{}12'' (``sig'' = emp\_\allowbreak{}p $<$ 0.05). Against the stated criterion ``$z$ $>$ 3 in $\geq$6/\allowbreak{}12'', $\SV{4}$ passes at 6/\allowbreak{}12 with $L_{1}$ $z$ = 3.00003, and $\SV{5}$ fails at 2/\allowbreak{}12 \\
6 & 15 brief & $\SV{2}$--$\SV{4}$ & --- & $\SV{5}$ dropped: ``$\SV{5}$ (deprioritized; rescue only if $\SV{6}$+ shows new pattern)'' & $\SV{4}$ kept by a z-margin of $3\times10^{-5}$ at one layer; the group becomes $\SV{2}$, $\SV{3}$, $\SV{4}$ \\
7 & 16 H02, 19 H02, 20 & $\SV{2}$--$\SV{4}$ & same as above & ``$\SV{2}$/\allowbreak{}$\SV{3}$/\allowbreak{}$\SV{4}$ axis independence''; ``multi-axis co-polarity count ($\SV{2}$+$\SV{3}$+$\SV{4}$)'' & the $\SV{2}$--$\SV{4}$ triple becomes a unit, defined by PPI co-pole on the kidney 209-gene subset \\
8 & 47 H01/\allowbreak{}H02 & all & immune, 2,039 observed & the screen finds the zero rows: ``Found 2902 genes have zero norm at ALL layers'', and that the earlier ``93.4\%'' $\SV{1}$ figure ``was dominated by zero vectors'' & switch to observed rows only \\
9 & 48 brief & $\SV{2}$--$\SV{4}$ & immune, observed & the kidney band is carried over unchanged: ``H02 (HIGH VALUE): $\SV{2}$--$\SV{4}$ PPI Geometry in 2039-Gene Active Set'', comparing with ``prior $\rho$ values ($\SV{2}$: 1.0, $\SV{3}$: 0.9, $\SV{4}$: 0.7)'' from the kidney matrix & --- \\
10 & 49 H03 & $\SV{2}$--$\SV{4}$ & immune, observed; coordinates = left singular vectors & first band-level test of $\SV{2}$--$\SV{4}$ on TRRUST pairs & ``Best: $L_{8}$ $p$$<$0.0001 effect=0.268; significant ($p$$<$0.05) at 8/\allowbreak{}12 layers'' \\
11 & 49 next steps & $\SV{5}$--$\SV{7}$ & --- & \textbf{origin of $\SV{5}$--$\SV{7}$}: ``Q: Is the effect specific to $\SV{2}$--$\SV{4}$, or does it extend further into the spectrum? Test $\SV{5}$--$\SV{7}$, $\SV{8}$--$\SV{10}$ subspaces'' & fixed-width windows of three, chosen after the $\SV{2}$--$\SV{4}$ result \\
12 & 50 H03 & $\SV{5}$--$\SV{7}$ & immune, observed & first recorded $\SV{5}$--$\SV{7}$ (with $\SV{1}$, $\SV{2}$--$\SV{4}$, $\SV{8}$--$\SV{10}$) & ``$\SV{2}$--$\SV{4}$: rbc=+0.229 (10/\allowbreak{}12 layers $p$$<$0.05); $\SV{5}$--$\SV{7}$: rbc=+0.152 (7/\allowbreak{}12 layers); $\SV{8}$--$\SV{10}$: rbc=+0.123 (5/\allowbreak{}12 layers)'' \\
13 & 50 H01, 51 H01 & $\SV{2}$--$\SV{4}$ & immune, observed & $\SV{2}$--$\SV{4}$ TRRUST proximity tested against ``co-expression'' (in practice the cosine of the full 512-d embeddings, not expression correlation) & ``0/\allowbreak{}12 layers specific''; residualised: ``1/\allowbreak{}12 layers significant ($L_{8}$: rbc=0.083, $p$=0.0016)'' \\
14 & 51 brainstormer & $\SV{5}$--$\SV{7}$ & --- & $\SV{5}$--$\SV{7}$ promoted as the alternative: ``$\SV{5}$--$\SV{7}$ regulatory signal (independent of co-expression) \textbar{} UNTESTED \textbar{} High potential'' & --- \\
15 & 52 H03 & $\SV{5}$--$\SV{7}$ & immune, observed & $\SV{5}$--$\SV{7}$ residualised & ``rbc\_\allowbreak{}resid=0.148 ($L_{0}$), 0.119 ($L_{1}$), 0.083 ($L_{2}$); all $p$$<$0.002. $L_{3}$+ lose significance''. Recomputed exactly (0.148 /\allowbreak{} 0.119 /\allowbreak{} 0.083); layers 5--11 have \emph{negative} rbc \\
16 & 53 H01 & $\SV{5}$--$\SV{7}$ & immune, observed & other bands scanned only after $\SV{5}$--$\SV{7}$ had been chosen: ``$\SV{2}$--$\SV{4}$, $\SV{5}$--$\SV{7}$, $\SV{8}$--$\SV{10}$, $\SV{11}$--$\SV{14}$, $\SV{15}$--$\SV{20}$, $\SV{21}$--$\SV{30}$'' at $L_{0}$/\allowbreak{}$L_{1}$/\allowbreak{}$L_{2}$/\allowbreak{}$L_{8}$ & ``$\SV{5}$--$\SV{7}$ dominant at $L_{0}$--$L_{2}$; $\SV{2}$--$\SV{4}$ dominant at $L_{8}$'' \\
17 & 55 H02 & $\SV{2}$--$\SV{7}$ & immune, observed & layer profiles of $\SV{5}$--$\SV{7}$ and $\SV{2}$--$\SV{4}$ classifiers & ``Complementary encoding: $\SV{5}$--$\SV{7}$ high early+late, $\SV{2}$--$\SV{4}$ high mid-depth'' \\
18 & 55 brainstormer & $\SV{2}$--$\SV{7}$ & --- & \textbf{origin of $\SV{2}$--$\SV{7}$}: ``can we dissect the joint $\SV{2}$--$\SV{7}$ space to produce a layer-stable classifier?''; decision threshold ``joint AUROC $\geq$ 0.72 at 10+ of 12 layers'' & --- \\
19 & 56 H01 & $\SV{2}$--$\SV{7}$ & immune, observed & joint 6-D classifier & ``joint $\geq$ 0.72 at 9/\allowbreak{}12 layers''. The threshold was then relaxed: ``Lower threshold for criterion since 9/\allowbreak{}12 vs 10/\allowbreak{}12 is minor; note claim should state `most layers''' \\
20 & 63 H63-B & all & immune, observed & the screen's own systematic scan of $\SV{k}$--$\SV{k+2}$, k = 0\ldots{}9, at $L_{0}$ with Bonferroni over 10 windows & ``Best raw AUROC: $\SV{8}$--$\SV{10}$=0.619 (perm\_\allowbreak{}p=0.67); $\SV{5}$--$\SV{7}$=0.611 (perm\_\allowbreak{}p=0.28); no window survives Bonferroni'' \\
\end{longtable}
\endgroup

\end{document}
